%\documentclass{article} % For LaTeX2e
%\usepackage{iclr2027_conference,times}

\documentclass[11pt]{article}

\usepackage[review]{acl}
\usepackage{times}
% Optional math commands from https://github.com/goodfeli/dlbook_notation.
\input{math_commands.tex}

\usepackage{booktabs}
\usepackage{graphicx}
\usepackage{hyperref}
\usepackage{url}
\usepackage{multirow}
\usepackage{adjustbox}
\usepackage{placeins}
\usepackage{comment}
\usepackage{capt-of}
\usepackage{graphicx}
\usepackage{wrapfig}


\usepackage{booktabs,multirow,makecell,graphicx,adjustbox}

\usepackage{xcolor}
\usepackage{colortbl}

\definecolor{table_green}{RGB}{34,139,34}
\definecolor{table_red}{RGB}{139,34,34}
\newcommand{\xiyan}[1]{\textcolor{table_green}{#1}}
%\newcommand{\mingzhe}[1]{\textcolor{table_red}{#1}}
% --- table styling ---

\definecolor{oursbg}{HTML}{E8F0FB}
\definecolor{tabgreen}{HTML}{1A7F37}
\definecolor{tabred}{HTML}{C0392B}
\definecolor{tabgray}{gray}{0.45}
\newcommand{\tabsd}[1]{{\color{tabgray}\scriptsize$\pm$#1}}
\newcommand{\tabup}[1]{{\color{tabgreen}\scriptsize$\uparrow$#1}}
\newcommand{\tabdn}[1]{{\color{tabred}\scriptsize$\downarrow$#1}}
\newcommand{\tabdng}[1]{{\color{tabgreen}\scriptsize$\downarrow$#1}}
\newcommand{\tabspd}[1]{{\color{tabgreen}\scriptsize$\times$#1}}
\newcommand{\tabpub}{$^\dagger$}


%\title{Kernelized Alignment for Latent Reasoning and Communication}
\title{Prototype-Based Latent Alignment for Multi-Agent Systems}
% Authors must not appear in the submitted version. They should be hidden
% as long as the \iclrfinalcopy macro remains commented out below.
% Non-anonymous submissions will be rejected without review.

\begin{comment}
    
\end{comment}
\author{
  First Author \\
  Department of Computer Science \\
  University A \\
  \texttt{first@example.com}
  \And
  Second Author \\
  Department of Computational \\
  Neuroscience, University B \\
  \texttt{second@example.com}
  \AND
  Third Author \\
  University C \\
  \texttt{third@example.com}
}



% The \author macro works with any number of authors. There are two commands
% used to separate the names and addresses of multiple authors: \And and \AND.
%
% Using \And between authors leaves it to \LaTeX{} to determine where to break
% the lines. Using \AND forces a linebreak at that point. So, if \LaTeX{}
% puts 3 of 4 authors names on the first line, and the last on the second
% line, try using \AND instead of \And before the third author name.

\newcommand{\fix}{\marginpar{FIX}}
\newcommand{\new}{\marginpar{NEW}}

\definecolor{table_green}{RGB}{34,139,34}
\definecolor{table_red}{RGB}{139,34,34}
\newcommand{\mingzhe}[1]{\textcolor{table_red}{#1}}

%\iclrfinalcopy % Uncomment for camera-ready version, but NOT for submission.
\begin{document}


\maketitle

\begin{comment}
\begin{abstract}
Multi agent systems rely on an intermediate medium to carry information across reasoning steps and between agents. Continuous latent representations offer an alternative to discrete text, but require aligning a sender’s hidden states with a receiver’s input embedding space. Existing training-free soft alignment methods solve this issue via vocabulary level probability distributions, yet incur substantial online costs from full vocabulary projection and embedding aggregation. We introduce \textit{Prototype Align}, an efficient, training-free alignment method for latent reasoning and multi-agent communication. Our key observation is that soft alignment can be expressed as attention over the sender’s unembedding matrix and the receiver’s embedding matrix. Using orthogonal random features, Prototype Align approximates this operation with a fixed-dimensional feature map, allowing vocabulary dependent statistics to be precomputed offline. Across six benchmarks, Prototype Align achieves 2.93× and 2.52× inference speedups over exact soft alignment while improving average accuracy by 2.97 and 2.28 percentage points on Qwen3-8B and Qwen3-14B, respectively. Cross-model experiments further demonstrate its applicability to agents with different hidden dimensions and compatible vocabularies. Further analyses provide insights into the behavior and effectiveness of Prototype Align. Our implementations are available at \url{https://anonymous.4open.science/r/kernel_align-258F}.
\end{abstract}
\end{comment}


\begin{abstract}
Continuous latent representations offer an alternative to text for reasoning and communication between agents, but require aligning a sender's hidden states with a receiver's input embeddings. Existing training-free methods rely on fixed linear mappings or probability-weighted embedding averages. 
%{\color{red}The former limits the expressiveness of alignment, while the latter uses distributions trained for discrete token prediction rather than continuous feedback.} 
%{\color{blue}
The former restricts nonlinear expressiveness, while the latter relies on token-prediction probabilities whose suitability for continuous latent feedback remains unclear.
%} 
Therefore, we propose \textit{Prototype Align}, which constructs a compact set of receiver-space prototypes from pretrained vocabulary parameters using orthogonal random features. Specifically, each hidden state determines a nonlinear mixture of these prototypes, which can be precomputed offline without additional training. Across six benchmarks, Prototype Align improves average accuracy over the strongest baseline by 3.19 and 3.15 percentage points under fixed-budget and entropy-based stopping on Qwen3-8B, respectively. 
Cross-model experiments further examine its applicability to agents with different hidden dimensions and compatible vocabularies. 
Further analyses provide insights into the behavior and effectiveness of Prototype Align. Our implementation is available at \url{https://anonymous.4open.science/r/kernel_align-258F}.
\end{abstract}


\section{Introduction}
%Agentic AI is becoming an increasingly important paradigm in solving complicated tasks, such as mathematics, coding, and scientific problems \citep{agentic_ai}. Prior works include a single agent which tackles tasks by combining its multi-step Chain-of-Thought reasoning ability \citep{COT_Jason2022}, i.e., intra-model reasoning . As tasks grow more difficult, more recent work extends this single agent paradigm into a collaborative one built from multiple agents system (MAS) \citep{hong2024metagpt}. Its performance depends heavily on effective coordination between agents, i.e., inter-model communication. Although intra-model reasoning and inter-model communication are typically studied as separate capabilities, we observe that they in fact stem from the same underlying mechanism. Each new step depends on the relevant information from the past. We therefore define this shared process as transmitting information across steps.

Agentic AI is becoming an increasingly important paradigm in solving complicated tasks, such as mathematics, coding, and scientific problems \citep{agentic_ai}. Early agentic systems typically relied on a single LLM to solve a task through multi-step reasoning \citep{wei2022cot}, self-reflection \citep{shinn2023reflexion}, or interactions with external tools \citep{schick2023toolformer}. As tasks grow more difficult, more recent work extends this single agent paradigm into a collaborative one built from multi-agent system (MAS) \citep{hong2024metagpt}. The effectiveness and efficiency of such systems consequently depend on how information is represented and transmitted between agents, which we refer to as the ``medium'' \citep{multiagent_survey,tran2025multi}.


As task complexity continues to grow and model capabilities improve, the context required to pass information across reasoning steps becomes increasingly long, further emphasizing the importance of the choice of medium. Early works use discrete \textit{text tokens} as the medium \citep{wei2022cot,hong2024metagpt,zhang2024chainofagents}, which are widely adopted and highly interpretable. However, sampling only a single token at each step discards useful distributional information, often leading to incomplete exploration of reasoning paths. Instead, \textit{latent representation} offers an alternative by transmitting continuous representations without explicitly decoding \citep{coconut}. Transmission via hidden states requires transferring representations from the unembedding space of the sender to the embedding space of the receiver, which induces a distribution compatibility problem. The mismatch exists even when agents use the same model and becomes more challenging when they use models of different scales. 

%\citet{shen2024collaborative,soft_thinking} further explore using logits as the medium to preserve this information. However, such methods can be computationally expensive due to the large vocabulary size, as decoding requires vocabulary-sized matrix multiplications \citep{decoding_time}. To address this issue, recent works \citep{hao2025training} introduce hidden states as the medium, which are both information-rich and decoding-free.

%Transmission via hidden states requires transferring representations from the unembedding space to the embedding space, which introduces a distribution shift problem. 

%{\color{blue}
A common approach to addressing this mismatch is training-based alignment, which adapts the model or learns an auxiliary adapter \citep{coconut,c2c,interlat,latent_cache_flow,cache_alignment}. Such adaptation requires additional training for new model configurations. Training-free alignment methods instead use existing model parameters to connect hidden states with input embeddings. Representative approaches include linear alignment \citep{latentmas} and soft alignment \citep{soft_thinking}. Linear alignment constructs a direct map from pretrained vocabulary parameters, but restricts the interface to a fixed linear transformation. This limits its ability to capture nonlinear, hidden-state-dependent weighting of receiver embeddings. Soft alignment forms a continuous input by averaging receiver embeddings according to the sender's next-token probabilities. Although it preserves the nonlinear distributional information, these probabilities are trained for discrete token prediction, and their exact embedding average is not necessarily the most effective input for subsequent latent reasoning. These limitations motivate a training-free alignment that retains nonlinear expressiveness while supporting effective continuous feedback for subsequent reasoning.

%Based on that, we propose ... (I am not sure whether the name "kernel" is still appropriate or not)}



%A common approach to addressing this shift is training-based alignment, which further trains the model to adapt to the shifted representation distribution, either by fine-tuning the entire model \citep{coconut} or by training an adapter \citep{c2c,interlat,latent_cache_flow,cache_alignment}. However, such approaches are computationally expensive, particularly for large-scale or frequently changing systems. More recently, training-free alignment methods have been proposed which solely rely on existing model parameters \citep{soft_thinking}. They project a sender's hidden state onto its full vocabulary distribution and then use this distribution to form a weighted combination of the receiver's input embeddings. Using the vocabulary as a shared interface, soft alignment connects models with different hidden dimensions without additional training. However, the computational cost of soft alignment scales with vocabulary size, making it expensive for modern LLMs with large vocabularies.
%Moreover, this alignment is only functionally effective for hidden states following common distributions, ignoring some vital outliers. 

%However, its full-vocabulary projection and embedding aggregation incur substantial online costs over long latent rollouts and repeated agent interactions.

%\citep{latentmas,later}. 
%Given the high efficiency, these approaches suffer from limited alignment accuracy in certain cases \citep{on_latentmas}, which can lead to instability in downstream reasoning and inter-agent communication. These limitations motivate us to develop a training-free yet accurate alignment method for hidden-state based agentic systems.

%In this work, we propose a kernelized alignment framework for hidden state based information transmission problem. In specific, we aim to find an efficient closed-form alignment function that directly approximates the ideal distribution mapping from sender to reciever. 
%To address this bottleneck, we introduce \emph{Prototype Align}, an efficient training-free function that directly approximates the ideal distribution mapping from sender to receiver. 
%To address this bottleneck, 

%We introduce \emph{Prototype Align}, an efficient training-free function that directly approximates the full vocabulary projection from sender to receiver. We observe that the vocabulary projection takes the form of a softmax attention operation constructed from the sender's unembedding matrix and the receiver's embedding matrix. Building on this observation, we apply kernel-based linear attention \citep{linear_attentions} with orthogonal random features \citep{performer} to replace the softmax over the full vocabulary $\mathcal{V}$ with a fixed-size feature map of dimension $m \ll |\mathcal{V}|$. In this way, the alignment depends only on the model's fixed embedding matrices and can therefore be precomputed and cached offline. This yields a training-free alignment procedure that avoids both the information loss of discrete sampling and the computational overhead of full-vocabulary decoding. It reduces the per-step alignment cost from $O(|\mathcal{V}|(d_A+d_B))$ to $O(m(d_A+d_B))$, where $d_A$, $d_B$ denote the embedding dimension of the sender and the receiver agent respectively.

We introduce \emph{Prototype Align}, a training-free nonlinear alignment. 
%{\color{red}Our key idea is to construct a compact set of receiver-space prototypes from the vocabulary correspondences between the sender and receiver, and combine them with hidden-state-dependent weights.} 
%{\color{blue}
To organize continuous feedback around reusable representations, we construct vocabulary-derived prototypes and mix them according to the sender's hidden state, providing an alternative to direct token-probability weighting.
%} 
Specifically, orthogonal random directions define weighted averages of receiver embeddings, while the sender's current hidden state determines how these prototypes are mixed to form the receiver's input. This construction is equivalent to an orthogonal random-feature approximation of vocabulary-mediated soft alignment \citep{linear_attentions,performer}, providing a mathematical basis for the prototypes and their mixing weights. Since the prototypes and normalization statistics depend only on fixed model parameters and sampled directions, they can be precomputed and cached offline. Online alignment therefore operates over $m$ prototypes rather than the full vocabulary $\mathcal{V}$, reducing the per-step alignment cost from $O(|\mathcal{V}|(d_A+d_B))$ to $O(m(d_A+d_B))$, where $m\ll|\mathcal{V}|$ and $d_A$, $d_B$ denote the sender and receiver hidden dimensions, respectively.


To empirically evaluate prototype alignment, we conduct experiments on six benchmarks spanning mathematical reasoning \citep{aime2024,aime2025}, code generation \citep{evalplus}, and scientific and medical question answering \citep{gpqa,medqa}. 
%Compared with soft alignment, Prototype Align achieves inference speedups of 2.93× and 2.52× while improving average accuracy by 2.97 and 2.28 percentage points on Qwen3-8B and Qwen3-14B, respectively. Beyond same model collaboration, cross model experiments show that Prototype Align also supports effective communication between agents with different model scales. Further analyses of aligned representations and latent reasoning trajectories provide insight into how the approximation affects information transmission. The sensitivity experiments suggest the stability of Prototype Alignment.
On Qwen3-8B, Prototype Align improves average accuracy over the strongest baseline by 3.19 and 3.15 percentage points under fixed-budget and entropy-based stopping, respectively. 
%Under entropy-based stopping, it achieves a 3.10$\times$ end-to-end inference speedup over Exact-Soft Align on Qwen3-8B, while using slightly more time under fixed budgets to achieve higher accuracy. 
Beyond same-model collaboration, cross-model experiments show that Prototype Align also supports effective communication between agents with different model scales. Further analyses of aligned representations and latent reasoning trajectories provide insight into how the approximation affects information transmission, and sensitivity experiments show that Prototype Align is stable across random-feature seeds and feature dimensions.

%Compared with Exact-Soft Align, Prototype Align improves average accuracy by 2.97 and 2.28 percentage points on Qwen3-8B and Qwen3-14B, respectively, while using 3.99$\times$ and 4.90$\times$ fewer tokens and achieving 2.93$\times$ and 2.52$\times$ faster inference. It also outperforms competitive baseline Linear Align in average accuracy by 3.19 and 2.33 percentage points at the two model scales.



\section{Related Work}
\textbf{Multi-Agent Collaboration} Multi-agent systems organize collaboration through role-based task decomposition \citep{hong2024metagpt}, multi-agent debate \citep{du2024debate}, and iterative feedback and refinement \citep{liang2024divergent} to jointly solve complex tasks. Early works primarily rely on text tokens for these systems \citep{wu2024autogen,wang2024survey}, which is time consuming and causes information loss \citep{liu2026beyond}. Recent work therefore explores collaboration beyond discrete text. Existing approaches enable latent collaboration by recovering structured latent factors and injecting them as prefixes \citep{zheng2025thought}, fusing activations within intermediate layers \citep{ramesh2025communicating}, or augmenting text with token wise hidden state differences \citep{tang-etal-2025-augmenting}. These methods differ in both the information they transmit and where it enters the receiver's computation. Complementary to these multi-agent collaboration mechanisms, our work focuses on final layer hidden states as the transmission medium and investigates their efficient alignment with the receiver's input embedding space. 



\noindent\textbf{Latent Alignment} Using latent representations for subsequent reasoning and agent communication relies on adapting them to the receiver's computation. Existing approaches achieve this adaptation through different levels. Training based approaches can adapt the LLMs itself to use latent representations, by curriculum learning \citep{coconut} or self-distillation \citep{shen-etal-2025-codi}. Other approaches learn auxiliary adapters to transform latent representations, either at the hidden-state interface \citep{softcot,interlat} or within the receiver's KV caches \citep{c2c,cache_alignment,latent_cache_flow}. Despite their strong performance, both approaches require additional training, making adaptation to new model configurations time-consuming. This limitation motivates training-free approaches that enable latent computation using existing model parameters. For example, \citet{pham2024let,soft_thinking} propagate probability weighted token embeddings, while \citet{latentmas} use a vocabulary derived linear map for hidden state feedback in latent multi-agent systems. Our method belongs to the training-free category. 
%{\color{red}It approximates the nonlinear vocabulary mediated transition with kernel features, enabling offline precomputation and reducing online alignment costs by avoiding full-vocabulary projection.} 
%{\color{blue}
It constructs vocabulary-derived prototypes for nonlinear latent feedback, with activations grounded in a readout--embedding approximation and fixed statistics that can be precomputed offline.
%}

%\textbf{Representations for reasoning and communication.}Intermediate representations play a central role in both multi-step reasoning and inter-agent communication. Chain-of-thought prompting expresses intermediate reasoning as natural-language tokens \citep{zelikman2022star,wei2022cot,hsieh2023distilling,deepseek_r1_distill_llama_8b}, while text-based multi-agent systems exchange messages to coordinate task decomposition, feedback, and refinement \citep{hong2024metagpt,zhang2024chainofagents}. Beyond discrete text, distribution-based approaches retain probabilities over multiple vocabulary items and use their weighted embeddings to continue reasoning \citep{shen2024collaborative}. Soft Thinking exemplifies this approach by propagating soft concepts instead of selecting a single token at each intermediate step \citep{soft_thinking}. Another line of work uses hidden states directly as intermediate representations, enabling continuous latent reasoning within a model and latent information exchange between agents \citep{coconut,latentmas}. By avoiding discrete token selection and explicit vocabulary decoding at intermediate steps, hidden states based methods can retain richer representations while reducing the computational overhead of reasoning and communication.



%\textbf{Intra-model reasoning.} Early works \cite{wei2022cot,kojima2022zeroshot} on intra-model reasoning use prompts to elicit CoT reasoning, which extends the thinking process and also improves the quality of the final answer. Following works \cite{zelikman2022star,hsieh2023distilling,deepseek_r1_distill_llama_8b} improve CoT reasoning through supervised fine-tuning with CoT generated by stronger models or humans, and end-to-end reinforcement learning. To further extend the exploration capability of reasoning, \cite{yao2023tot,besta2024got} propose to use tree- or graph-structured reasoning to search for multiple paths in parallel; to improve the reliability and stability of reasoning, \cite{shinn2023reflexion,yao2023react} use self-reflection mechanism or tools to re-evaluate the reasoning process and correct errors of previous steps. Recently, it has been found that recursive reasoning in one step without explicit CoT can also achieve comparable performance to CoT reasoning \cite{coconut}, which can also be regarded as a special case of latent reasoning.

%\textbf{Inter-model communication.} One single agent is insufficient to handle complex tasks, but a multi-agent system (MAS) that consists of different experts assist to bridge the gap. \cite{zhang2024chainofagents,hong2024metagpt} use sequential message flows to enable inter-model communication, and \cite{wang2024moa,liu2024dylan,hu2025owl} create more complex communication topologies, such as hierarchical, star, and more general graph structures. In order to overcome the hallucination and error propagation inherent to individual models, \cite{du2024debate,estornell2024multillm,liang2024divergent} introduce iterative communication rounds, where agents can discuss and refine their responses via adversarial or cooperative dynamics. Recent works such as \cite{anthropic2024mcp,google2025a2a} aim to establish a protocol for communication in MAS, which standardizes the interaction and tool usage of LLM agents.

%\xiyan{an illustration graph is necessary in the introduction}


%We formalize intra-model reasoning and inter-model communication as a same information transmission problem, which align the hidden state  before the unembedding matrix to the embedding space in Section~\ref{sec:problem_formulation}. Building on this unified view, we propose a kernelized alignment method to solve this transmission problem in Section~\ref{sec:kernelized_alignment}. Figure~\ref{fig:method_overview} illustrates the method.
\section{Prototype Align}
%We formulate text based and direct information transmission from a sender's hidden state space to a receiver's input embedding space in Section~\ref{sec:problem_formulation}. We then introduce Prototype Align, an efficient approximation to the direct transmission mapping, in Section~\ref{sec:kernelized_alignment}. Figure~\ref{fig:method_overview} shows an overview. Finally, Section~\ref{sec:agent_integration} shows how the Prototype Alignment is applied in multi-agent collaboration. 
We study how a hidden state produced by a sender $A$ can be transformed into an input embedding for a receiver $B$. Here, we consider sender and receiver with a shared vocabulary $\mathcal{V}$, while allowing their hidden dimensions $d_A$ and $d_B$ to differ. Figure~\ref{fig:method_overview} provides an overview. 
%We first formulate text-based and latent information transfer per transmission step in Section~\ref{sec:problem_formulation}. We then introduce the Prototype Align, an efficient approximation to the direct transmission mapping, in Section~\ref{sec:kernelized_alignment}. Figure~\ref{fig:method_overview} provides an overview. Finally, we apply the single step latent transmission to iterative reasoning within an agent and latent sequence transmission between agents in Section~\ref{sec:agent_integration}.

\begin{figure*}[t]
\centering
\includegraphics[width=\textwidth]{figs/model.pdf}
\caption{\textbf{Top:} Text-mediated transmission projects the sender's hidden state onto the vocabulary, samples a token, and retokenizes and embeds it for the receiver. \textbf{Bottom:} 
%Under compatible vocabularies, our method approximates the corresponding soft transition, defined as the probability-weighted average of receiver embeddings, using positive orthogonal random features. Vocabulary-dependent numerator and normalization statistics are precomputed offline; online alignment only evaluates the hidden-state feature map and combines it with these cached statistics, avoiding explicit vocabulary decoding.
Prototype Align precomputes vocabulary-derived prototypes offline and mixes them according to the sender's hidden state online. The equivalent kernel-feature implementation shown here avoids full-vocabulary computation during alignment.
}
\label{fig:method_overview}
\end{figure*}


%\xiyan{\textbf{Unified View: Information Transmission}. Intra-model reasoning and inter-model communication are often treated as distinct capabilities. However, we observe that they share the same underlying operation, where a hidden state produced by one step is decoded into a textual representation and subsequently encoded by its next step. The next step may be in the same agent, in which case the transition constitutes an autoregressive reasoning step, or in a different agent, in which case it constitutes a communication. We therefore formalize this shared processing as transmitting information problem across steps.}


\begin{comment}
\paragraph{Text Based Information Transmission.}
The conventional text based transmission projects the
hidden state of sender \(A\) onto its vocabulary, samples a discrete token, and subsequently retokenizes and embeds the sampled output using receiver \(B\). Specifically, at each generation step \(t\), an input text token is mapped into the embedding space by the embedding matrix \(W^{\mathrm{in}}\) and processed by the Transformer blocks \(\Phi\) to obtain the corresponding hidden state $h_t$:
\begin{equation}
        e_t = W^{\mathrm{in}}x_t,
        \qquad
        h_t = \Phi(e_t;h_{<t})
    \label{eq:text_encoding}
\end{equation}
where $e_t$ is the token embedding and \(h_{<t}\) denotes the preceding context. The hidden state is subsequently decoded into a probability distribution $p_t$ over the vocabulary:
\begin{equation}
        p_{t}
        =
        \operatorname{Softmax}
        \left(
            W^{\mathrm{out}}h_t
        \right)
    \label{eq:text_decoding}
\end{equation}
The output token $y_{t}$ is then sampled from this distribution and transmitted to the next step $t+1$. Formally, the input embedding received by receiver \(B\) is given by
\(e_{t+1}^{B}
=
W_B^{\mathrm{in}}T_{A\rightarrow B}(y_{t}^{A})\), where $W_B^{\mathrm{in}}$ is the embedding matrix and $T_{A\rightarrow B}$ is the detokenization process.
\end{comment}

\subsection{Text and Latent Information Transmission}
\label{sec:problem_formulation}
\paragraph{Text Based Information Transmission.}
%Conventional text-based transmission projects a sender's hidden state onto its vocabulary, samples a discrete token, and converts the sampled token into text for the receiver to retokenize and embed. Specifically, 
At each generation step \(t\), an input text token is mapped into the embedding space by the embedding matrix \(W^{\mathrm{in}}_A\) and processed by the Transformer blocks \(\Phi\) to obtain the corresponding hidden state $h_t^A$:
\begin{equation}
    e_t^A = W_A^{\mathrm{in}}x_t^A,
    \qquad
    h_t^A = \Phi_A(e_t^A; C_t^A),
    \label{eq:text_encoding}
\end{equation}
where $e_t^A$ is the input embedding and $C_t^A$ is the preceding context. The resulting hidden state is projected through sender's unembedding matrix to obtain a probability distribution over its vocabulary:
\begin{equation}
    p_t^A
    =
    \operatorname{Softmax}
    \left(
        W_A^{\mathrm{out}}h_t^A
    \right).
    \label{eq:text_decoding}
\end{equation}
A token $y_t^A \sim p_t^A$ is then sampled and transmitted to receiver $B$. Assuming that the sampled token corresponds to a single receiver token, the input embedding at the next step is
\begin{equation}
    e_{t+1}^B
    =
    W_B^{\mathrm{in}}y_t^A,
    \label{eq:text_transmission}
\end{equation}
where $W_B^{\mathrm{in}}$ is the receiver's embedding matrix.
%, and $T_{A\rightarrow B}$ denotes detokenization with the sender's tokenizer followed by retokenization with the receiver's tokenizer.
Thus, information is transmitted through explicitly decoded text tokens, which are subsequently embedded by the receiver.

\begin{comment}
The conventional text based transmission process projects the
hidden state of sender \(A\) onto its vocabulary, samples a discrete token, and subsequently retokenizes and embeds the sampled output using receiver \(B\). Formally, the input embedding received by agent \(B\) is given by
\(e_{t+1}^{B}
=
W_B^{\mathrm{in}}T_{A\rightarrow B}(y_{t+1}^{A})\), where $W_B^{\mathrm{in}}$ is the embedding matrix and $T_{A\rightarrow B}$ is the 
\end{comment}

\paragraph{Latent Information Transmission.}
Text-based transmission selects a single token embedding according to the sender's next-token distribution. A continuous alternative is to replace this sampled embedding with its probability weighted average \citep{soft_thinking}. 
%We consider sender and receiver models with a shared vocabulary $\mathcal{V}$ and matching token indices, while allowing their hidden dimensions $d_A$ and $d_B$ to differ.
Let $w_{B,j}^{\mathrm{in}}\in\mathbb{R}^{d_B}$ denote the receiver's input embedding for token $j$, we define the soft alignment mapping as:
\begin{equation}
\begin{aligned}
    F_{A\rightarrow B}(h_t^A)
    &:= \sum_{j\in\mathcal{V}}
        p_t^A(j)\,w_{B,j}^{\mathrm{in}}, \\
    p_t^A
    &= \operatorname{Softmax}\!
        \left(\frac{W_A^{\mathrm{out}}h_t^A}{\tau}\right).
\end{aligned}
\label{eq:distribution_transition}
\end{equation}
where $h_t^A$ is the sender's final layer hidden state at step $t$, and $\tau>0$ is a temperature parameter controlling the sharpness of the transmitted distribution. The receiver uses $e_{t+1}^B=F_{A\rightarrow B}(h_t^A)$ as its continuous input. 
%This mapping connects the sender's hidden state space to the receiver's input embedding space through their shared vocabulary, without selecting a discrete token. 
This mapping avoids discrete token sampling and requires no additional training, but each evaluation computes a full-vocabulary distribution and aggregates receiver embeddings, costing $O(|\mathcal{V}|(d_A+d_B))$.
%Although \(F_{A\rightarrow B}\) avoids the information loss and additional training, each evaluation computes a full-vocabulary distribution and aggregates the corresponding receiver embeddings, costing $O(|\mathcal{V}|(d_A+d_B))$. Since modern LLMs typically have large vocabularies, repeating these vocabulary-wide operations at every transmission step incurs substantial computational overhead. The difficulty is that the softmax couples all $|\mathcal{V}|$ tokens nonlinearly: filtering (e.g., top-$k$) still requires the full logits, and linear surrogates \citep{latentmas} lose the softmax nonlinearity. We therefore seek to \textbf{preserve the nonlinearity of soft alignment while eliminating vocabulary-wide computation from each transmission step}, which, as we show next, becomes possible by viewing soft alignment as attention over fixed vocabulary parameters.
%We therefore aim to \textbf{approximate this soft alignment mapping without performing vocabulary level computation at every step}.
%Since modern LLMs typically have large vocabularies, repeating these vocabulary wide operations at every transmission step incurs substantial computational overhead. We therefore aim to \textbf{approximate this soft alignment mapping while eliminating vocabulary wide computation from each transmission step}.

%Since modern LLMs typically have large vocabularies, repeating these vocabulary-wide operations at every transmission step incurs substantial computational overhead. Avoiding this cost is non-trivial. The softmax couples all $|\mathcal{V}|$ tokens through its normalization, making $F_{A\rightarrow B}$ a nonlinear function of $h_t^A$ whose vocabulary-dependent part cannot simply be precomputed. Truncating the distribution (e.g., top-$k$) still requires the full vocabulary logits to select the candidates, while replacing the mapping with a linear map \citep{latentmas} removes the softmax nonlinearity and over-smooths the transmitted distribution. The key challenge is therefore to \textbf{preserve the distribution-level nonlinearity of soft alignment while eliminating vocabulary-wide computation from each transmission step}. As we show next, this becomes possible once soft alignment is viewed as attention over fixed vocabulary parameters.

\begin{comment}
\paragraph{Latent Information Transmission.}
By compressing the complete vocabulary distribution \(p_{t}^{A}\) into a single discrete token, sampling inevitably discards part of the information encoded in the distribution. To avoid this information loss, an ideal alternative is to omit the sampling operation and directly propagate the complete vocabulary distribution. After aligning \(p_{t}^{A}\) with the vocabulary of agent \(B\), the ideal information transmission $F_{A\rightarrow B}$ is denoted as:
\begin{equation}
    e_{t+1}^{B}
    =
    W_B^{\mathrm{in}}
    T_{A\rightarrow B}
    \left(
        \operatorname{Softmax}
        \left(
            \frac{W_A^{\mathrm{out}}h_t^A}{\tau}
        \right)
    \right).
\label{eq:distribution_transition}
\end{equation}
Here $\tau$ is a temperature parameter controlling the sharpness of the transmitted distribution. Although \(F_{A\rightarrow B}\) avoids the information loss caused by discrete sampling, its full vocabulary projection and re-embedding at every generation step, resulting in substantial computational and communication overhead. We therefore \textbf{seek an efficient closed-form alignment function \(\widehat{F}_{A\rightarrow B}\) that directly approximates the ideal distribution level mapping}.
\end{comment}

%\paragraph{Unified View of Intra-model Reasoning and Inter-model Communication.} We now associate each generation step with the agent that executes it. Generation step \(t\) is executed by agent \(A\), while the resulting output is consumed by agent \(B\) at step \(t+1\). Specifically, the token generated by \(A\) is decoded into text and subsequently transmitted to \(B\), a process denoted by \(T_{A\rightarrow B}\). The identities of agents \(A\) and \(B\) give rise to two cases: i) when $A=B$, the two consecutive steps are executed by the same agent, and the transmission corresponds to an intra-model autoregressive reasoning step; ii) when $A \neq B$, the output is transmitted between different agents, and the transmission corresponds to inter-model communication. Therefore, intra-model reasoning and inter-model communication can be unified as the same information-transmission problem between consecutive generation steps.

\begin{comment}
\paragraph{Direct Information Transmission.}
The conventional text-mediated transmission process projects the
hidden state of agent \(A\) onto its vocabulary, samples a discrete token,
and subsequently retokenizes and embeds the sampled output using agent
\(B\). Formally, the input embedding received by agent \(B\) is given by
\(e_{t+1}^{B}
=
W_B^{\mathrm{in}}T_{A\rightarrow B}(y_{t+1}^{A})\).
By compressing the complete vocabulary distribution \(p_{t+1}^{A}\) into a single discrete token, sampling inevitably discards part of the information encoded in the distribution. To avoid this information loss, an ideal alternative is to omit the sampling operation and directly propagate the complete vocabulary distribution. After aligning \(p_{t+1}^{A}\) with the vocabulary of agent \(B\), the ideal information transmission $F_{A\rightarrow B}$ is denoted as:
\begin{equation}
    e_{t+1}^{B}
    =
    W_B^{\mathrm{in}}
    T_{A\rightarrow B}
    \left(
        \operatorname{Softmax}
        \left(
            \frac{W_A^{\mathrm{out}}h_t^A}{\tau}
        \right)
    \right).
\label{eq:distribution_transition}
\end{equation}
Here $\tau$ is a temperature parameter controlling the sharpness of the transmitted distribution. Although \(F_{A\rightarrow B}\) avoids the information loss caused by discrete sampling, its full vocabulary projection and re-embedding at every generation step, resulting in substantial computational and communication overhead. We therefore \textbf{seek an efficient closed-form alignment function \(\widehat{F}_{A\rightarrow B}\) that directly approximates the ideal distribution level mapping}.
\end{comment}


\subsection{Prototype-Based Align}
\label{sec:prototype_alignment}
%The shared vocabulary provides paired correspondences between the sender's unembedding vectors and the receiver's input embeddings. We use these correspondences to precompute a compact set of receiver space prototypes, then combine them with weights determined by the sender's hidden state.

\paragraph{Vocabulary-Derived Prototypes.}
We construct a shared set of continuous prototypes from the sender's unembedding vectors and the receiver's input embeddings. Let $q=h_t^A/\sqrt{\tau}$, $k_j=w_{A,j}^{\mathrm{out}}/\sqrt{\tau}$, and $v_j=w_{B,j}^{\mathrm{in}}$, where $\tau>0$ is the alignment temperature. We sample $m$ Gaussian probe vectors $\omega_i\in\mathbb{R}^{d_A}$ without additional training. To spread the probes across distinct directions, we use orthogonal random features (ORFs) \citep{performer}, making the vectors orthogonal within each block while preserving their Gaussian marginals. Each probe defines a prototype by aggregating receiver embeddings according to their corresponding sender keys:
\begin{equation}
\begin{aligned}
    b_{ij}
    &=\exp\!\left(\omega_i^\top k_j-\frac{\|k_j\|^2}{2}\right), \\
    Z_i&=\sum_{j\in\mathcal{V}}b_{ij},
    \qquad
    u_i=\frac{\sum_{j\in\mathcal{V}}b_{ij}v_j}{Z_i}.
\end{aligned}
\label{eq:prototype_construction}
\end{equation}
The normalized weights satisfy $b_{ij} \propto \operatorname{Softmax}_j(-\|\omega_i-k_j\|^2/2)$. Thus, each prototype averages receiver embeddings using normalized Gaussian radius weights computed in the sender's key space.

\paragraph{Query-Dependent Prototype Mixing.}
Our transmission protocol forms a query-dependent mixture of the shared prototypes. Given positive activations $\{a_i(q)\}_{i=1}^{m}$, the receiver's continuous input is defined as:
\begin{equation}
    \widehat{F}_{A\rightarrow B}(h_t^A)
    =\frac{\sum_{i=1}^{m}a_i(q)u_i}
            {\sum_{i=1}^{m}a_i(q)}
    =\sum_{i=1}^{m}\alpha_i(q)u_i.
\label{eq:prototype_mixing}
\end{equation}

The normalized weights $\alpha_i(q)=a_i(q)/\sum_{\ell=1}^{m}a_\ell(q)$ sum to one, restricting the transmitted representation to the convex hull of the shared prototypes. To instantiate this protocol without additional training, we design the activations as $a_i(q)=\exp(\omega_i^\top q)Z_i$. This design reduces the mismatch between the sender's output and the receiver's input, as formalized below.

\paragraph{Theorem (Training-Free Approximation).}
With $a_i(q)=\exp(\omega_i^\top q)Z_i$, the prototype mixture in Eq.~\ref{eq:prototype_mixing} is an estimator of the readout--embedding mapping:
\begin{equation}
    F_{A\rightarrow B}(h_t^A)
    =\frac{\sum_{j\in\mathcal{V}}\exp(q^\top k_j)v_j}
           {\sum_{j\in\mathcal{V}}\exp(q^\top k_j)}.
\label{eq:prototype_attention_target}
\end{equation}
This mapping is softmax attention over the sender's unembedding keys and the receiver's embedding values. For a fixed query and finite vocabulary, the estimator converges in probability to this mapping as the number of features increases. The derivation and proof are provided in Appendix~\ref{app:prototype_approximation}.

\paragraph{Offline Computing.}
The prototype mapping admits an accelerated implementation because $\{u_i\}$ and $\{Z_i\}$ depend only on pretrained vocabulary parameters, the temperature, and the sampled directions. These quantities can therefore be computed once and cached for each sender--receiver pair. Online alignment only evaluates the query activations and mixes the cached prototypes, reducing the per-step cost from $O(|\mathcal{V}|(d_A+d_B))$ to $O(m(d_A+d_B))$. This offline computation implements the same alignment without changing its transmitted representations.

\begin{comment}
    

\subsection{Kernelized Alignment}
\label{sec:kernelized_alignment}
%We propose a \textit{kernalized alignment} strategy, which is a closed form alignment framework designed to approximate the soft alignment \(F_{A\rightarrow B}\) denoted in Eq \ref{eq:distribution_transition}. 

\paragraph{Soft Alignment as Vocabulary Attention.} The soft alignment uses the sender's next token probabilities as weights to combine the receiver's token embeddings. To expose this structure, we extend the soft alignment in Eq.~\ref{eq:distribution_transition} as:
\begin{equation}
F_{A \rightarrow B}(h_t^A)
=
\frac{
\sum_{j=1}^{|\mathcal{V}|}
\exp\!\left(\frac{h_t^A(w_{A,j}^{out})^\top}{\tau}\right)w_{B,j}^{in}
}{
\sum_{j=1}^{|\mathcal{V}|}
\exp\!\left(\frac{h_t^A(w_{A,j}^{out})^\top}{\tau}\right)
}.
\label{eq:vocab_attention}
\end{equation}
%where $w_{A,j}^{out}$ is the row vector form of the $j$-th output column of unmbedding matrix $W_{A}^{out}$ and $w_{B,j}^{in}$ is the corresponding target embedding row.
where $w_{A,j}^{\mathrm{out}}$ denotes the transpose of the $j$-th row of the sender's unembedding matrix $W_A^{\mathrm{out}}$, and $w_{B,j}^{\mathrm{in}}$ denotes the $j$-th column of the receiver's input embedding matrix $W_B^{\mathrm{in}}$.
We observe that the extension form can be expressed as attention over fixed vocabulary parameters. Here, the sender's hidden state $h_t^A$ acts as the query, its unembedding vectors $\{w_{A,j}^{out}\}_{j=1}^{|\mathcal{V}|}$ act as keys, and the receiver's input embeddings $\{w_{B,j}^{in}\}_{j=1}^{|\mathcal{V}|}$ act as values. Here, only the query changes across transmission steps, while the keys and values are fixed pretrained parameters. 


 

\paragraph{Kernel Approximation.}
\label{sec:kernel}
The above characteristics motivate us to accelerate the computation using kernel based linear attention techniques \citep{linear_attentions}. Specifically, following \citet{performer}, we approximate the exponential kernel with a finite dimensional positive random feature map:
$\phi:\mathbb{R}^{d_A}\rightarrow\mathbb{R}^{m}$:
\begin{equation}
\begin{aligned}
    \exp(qk^{\top})
    &\approx \phi(q)\phi(k)^{\top}, \\
    \phi(x)
    &= \frac{1}{\sqrt{m}}
       \exp\!\left(-\frac{\|x\|^2}{2}\right) \\
    &\quad {}\times
       \left[
       e^{x\omega_1^\top},\ldots,
       e^{x\omega_m^\top}
       \right].
\end{aligned}
\label{eq:kernel_approximation}
\end{equation}
where $\{\omega_i\}_{i=1}^m$ are sampled orthogonally from $\mathcal{N}(0,I)$ rather than independently. It reduces the variance of the Monte-Carlo estimate $\mathbb{E}[\phi(q)\phi(k)^\top]=\exp(qk^\top)$ (Appendix \ref{app:orf_accuracy}) and thus is expected to improve approximation accuracy. Applying the kernel approximation in Eq.~\ref{eq:kernel_approximation} to Eq.~\ref{eq:vocab_attention}, the soft alignment function can be approximated as:
\begin{equation}
\widehat{F}_{A \rightarrow B}(h_t^A)
=
\frac{
\sum_{j=1}^{|\mathcal{V}|}
\phi\!\left(\frac{h_t^A}{\sqrt{\tau}}\right)
\phi\!\left(\frac{w_{A,j}^{out}}{\sqrt{\tau}}\right)^\top w_{B,j}^{in}
}{
\sum_{j=1}^{|\mathcal{V}|}
\phi\!\left(\frac{h_t^A}{\sqrt{\tau}}\right)
\phi\!\left(\frac{w_{A,j}^{out}}{\sqrt{\tau}}\right)^\top
}.
\label{eq:kernel_substituted}
\end{equation}

\paragraph{Offline Precomputation and Online Alignment.}
Since $\phi(h_t^A/\sqrt{\tau})$ does not depend on the summation index $j$, it can be factored out of both sums:
\begin{equation}
\begin{aligned}
\widehat{F}_{A \rightarrow B}(h_t^A)
&=
\frac{
  \phi\!\left(\frac{h_t^A}{\sqrt{\tau}}\right)
  S_{A \rightarrow B}
}{
  \phi\!\left(\frac{h_t^A}{\sqrt{\tau}}\right) z_A
}, \\[4pt]
S_{A \rightarrow B}
&=
\sum_{j=1}^{|\mathcal{V}|}
\phi\!\left(
  \frac{w_{A,j}^{\mathrm{out}}}{\sqrt{\tau}}
\right)^{\!\top}
w_{B,j}^{\mathrm{in}}, \\[4pt]
z_A
&=
\sum_{j=1}^{|\mathcal{V}|}
\phi\!\left(
  \frac{w_{A,j}^{\mathrm{out}}}{\sqrt{\tau}}
\right)^{\!\top}.
\end{aligned}
\label{eq:linear_alignment}
\end{equation}

The two sums in Eq.~\ref{eq:linear_alignment} depend only on the fixed matrices $W_A^{out}$ and $W_B^{in}$ which can be computed once and cached offline. Hence, each alignment query $h_t^A$ only needs to be projected through $\phi$ and matched against the cached sums, without repeating the full-vocabulary computation at every generation step. Moreover, since $W_A^{out}$ and $W_B^{in}$ can be directly retrieved from the model, this precomputation itself requires no additional training. As a result, the per-step complexity of the alignment is reduced from $O(|\mathcal{V}|(d_A+d_B))$ to $O(m(d_A+d_B))$ with $m\ll|\mathcal{V}|$, yielding significant speedup in inference time.

\end{comment}

%Prototype Align provides a common interface for latent reasoning within an agent and latent communication between agents. Let $A$ denote the agent producing a final-layer hidden state and $B$ the agent consuming its aligned representation. The transmission is defined as $e^B=\widehat{F}_{A\rightarrow B}(h^A)$. The identities of $A$ and $B$ determine how this interface is used: i) \textit{Intra-agent reasoning ($A=B$).} The aligned hidden state serves as the input embedding for the same agent's next reasoning step. Repeating this operation produces a sequence of continuous reasoning states without sampling intermediate text tokens. ii) \textit{Inter-agent communication ($A\neq B$).} Prototype Align maps the sender's hidden states into the receiver's input-embedding space. The resulting embeddings are incorporated into the receiver's input context alongside its task instructions.

%Prototype Align constructs each sender and receiver operator directly from the sender's unembedding matrix and the receiver's input embedding matrix, without additional training. This interface supports homogeneous and heterogeneous multi-agent systems. In homogeneous systems, agents share the same model checkpoint, allowing a single alignment operator to be reused across agent roles. In heterogeneous systems, agents use different model checkpoints with different hidden dimensions. 

\subsection{Integration into Multi-Agent Collaboration}
\label{sec:agent_integration}
\paragraph{Intra-Agent Reasoning.} Within an agent’s latent reasoning process, the sender and receiver for each step coincide, i.e, $A=B$. Prototype Align transforms the current hidden state into the input embedding for the next reasoning step. Repeating this operation produces a sequence of latent reasoning states without sampling text tokens.



\paragraph{Inter-Agent Communication.}
Communication from agent $M_i$ to agent $M_j$ transfers a sequence of latent reasoning steps. Let $H^i=(h_1^i,\ldots,h_L^i)$ denote the sequence of $L$ hidden steps to be transmitted. Using the alignment operator introduced in Section \ref{sec:prototype_alignment}, these steps can be mapped to the receiver's input embedding space to form a continuous message:
\begin{equation}
    E_{i\rightarrow j}
    =
    \left(
        \widehat{F}_{i\rightarrow j}(h_1^i),
        \ldots,
        \widehat{F}_{i\rightarrow j}(h_L^i)
    \right),
    \label{eq:inter_agent_message}
\end{equation}
where $\widehat{F}_{i\rightarrow j}$ denotes the alignment operator with $M_i$ as the sender and $M_j$ as the receiver. This formulation accommodates both homogeneous agents that share the same model checkpoint and heterogeneous agents that use models of different scales with a shared vocabulary.

\begin{comment}
    

The communication mechanism depends on the agents' model configurations. In the homogeneous case, all agents share the same model checkpoint.
To preserve the accumulated reasoning context while
avoiding redundant computation, we directly transfer
$M_i$'s KV cache to $M_j$, without explicitly constructing
and processing the aligned message above.
In the heterogeneous case, the agents use different
model checkpoints with a shared vocabulary, and their
cached representations are not directly interchangeable.
We therefore apply the sequence-level alignment in
Eq.~\ref{eq:inter_agent_message}.
The receiver incorporates the aligned embeddings into
its input context in their original order and continues
reasoning conditioned on this context.
\end{comment}

\begin{comment}
Consider information transmission between agents, e.g., $M_i \rightarrow M_j$, the communication mechanism depends on whether they use the same model checkpoint. In the homogeneous case, agents use the same model checkpoint. Hence, no additional alignment between agents is required. To preserve the preceding reasoning context while avoiding redundant computation, we transfer $M_i$'s cache directly to the $M_j$. In our heterogeneous setting, the sender and receiver usemodels of different scales with a shared vocabulary.Because their internal representations are not directly compatible, we transfer information through aligned input embeddings.
Specifically, let $H^i=(h_1^i,\ldots,h_L^i)$ denote the
sequence of $L$ hidden states transmitted from agent $i$
to agent $j$.
We align each state to agent $j$'s input-embedding space
to form the message:
\begin{equation}
    M_{i\rightarrow j}
    =
    \left(
        \widehat{F}_{i\rightarrow j}(h_1^i),
        \ldots,
        \widehat{F}_{i\rightarrow j}(h_L^i)
    \right),
    \label{eq:inter_agent_message}
\end{equation}
where $\widehat{F}_{i\rightarrow j}$ instantiates the
alignment operator with agent $i$ as the sender and
agent $j$ as the receiver.

\mingzhe{Intra-agent reasoning: only align and pass the latest hidden state, history of the past hidden states are cached in the form of kv.
\\
Inter-agent communication: align and pass all history hidden states. specially, inter-agent comm. is mathematically equivalent with intra-agent reasoning but involve a piece of new rule prompt, thus we reuse past kv in this case for the same reason as intra-model reasoning.}
\end{comment}





%For each sender and receiver pair, Prototype Align constructs the alignment operator from the sender's unembedding matrix and the receiver's input-embedding matrix, without additional training. The heterogeneous setting requires a shared vocabulary with matching token indices, so that each sender vocabulary entry corresponds to the appropriate receiver embedding.

 


\begin{comment}
\subsection{Kernelized Alignment}
\label{sec:kernelized_alignment}
We propose a \textit{kernalized alignment} strategy, which is a closed-form alignment framework designed to approximate the ideal distribution level mapping \(F_{A\rightarrow B}\) denoted in Eq \ref{eq:distribution_transition}. Here, we limit that all agent share a same vocabulary $\mathcal{V}$, where the retokenization function $T_{A \rightarrow B}$ is the identity map. The ideal alignment function is then extended as:
\begin{equation}
F_{A \rightarrow B}(h_t^A)
=
\frac{
\sum_{j=1}^{|\mathcal{V}|}
\exp\!\left(\frac{h_t^A(w_{A,j}^{out})^\top}{\tau}\right)w_{B,j}^{in}
}{
\sum_{j=1}^{|\mathcal{V}|}
\exp\!\left(\frac{h_t^A(w_{A,j}^{out})^\top}{\tau}\right)
}.
\label{eq:vocab_attention}
\end{equation}
where $w_{A,j}^{out}$ is the row-vector form of the $j$-th output column and $w_{B,j}^{in}$ is the corresponding target embedding row. Evaluating this expression requires computation over the entire vocabulary at each generation step, which is time consuming. Notably, we observe that Eq.~\ref{eq:vocab_attention} takes the form of a softmax attention operation, with $h_t^A$ acting as the query, $\{w_{A,j}^{out}\}_{j=1}^{|\mathcal{V}|}$ as keys, and $\{w_{B,j}^{in}\}_{j=1}^{|\mathcal{V}|}$ as values. This connection motivates us to accelerate the computation using kernel-based linear attention techniques \citep{linear_attentions}. Specifically, following \citet{performer}, we approximate the exponential kernel with a finite-dimensional positive random feature map:
\begin{equation}
    \exp(qk^{\top})
    \approx
    \phi(q)\phi(k)^{\top}, \qquad
    \phi(x)=\frac{1}{\sqrt{m}}\exp\!\left(-\frac{\|x\|^2}{2}\right)
    \left[\exp(x\omega_1^\top),\ldots,\exp(x\omega_m^\top)\right],
\label{eq:kernel_approximation}
\end{equation}
where $\{\omega_i\}_{i=1}^m$ are sampled orthogonally from $\mathcal{N}(0,I)$ rather than independently. It reduces the variance of the Monte-Carlo estimate $\mathbb{E}[\phi(q)\phi(k)^\top]=\exp(qk^\top)$ (Appendix \ref{app:orf_accuracy}) and thus is expected to improve approximation accuracy. Applying the kernel approximation in Eq.~\ref{eq:kernel_approximation} to Eq.~\ref{eq:vocab_attention}, with $q=h_t^A/\sqrt{\tau}$ and $k=w_{A,j}^{out}/\sqrt{\tau}$, the ideal alignment function can be approximated as:
\begin{equation}
\widehat{F}_{A \rightarrow B}(h_t^A)
=
\frac{
\sum_{j=1}^{|\mathcal{V}|}
\phi\!\left(\frac{h_t^A}{\sqrt{\tau}}\right)
\phi\!\left(\frac{w_{A,j}^{out}}{\sqrt{\tau}}\right)^\top w_{B,j}^{in}
}{
\sum_{j=1}^{|\mathcal{V}|}
\phi\!\left(\frac{h_t^A}{\sqrt{\tau}}\right)
\phi\!\left(\frac{w_{A,j}^{out}}{\sqrt{\tau}}\right)^\top
}.
\label{eq:kernel_substituted}
\end{equation}
Since $\phi(h_t^A/\sqrt{\tau})$ does not depend on the summation index $j$, it can be factored out of both sums:
\begin{equation}
\widehat{F}_{A \rightarrow B}(h_t^A)
=
\frac{\phi\!\left(\frac{h_t^A}{\sqrt{\tau}}\right) S_{A \rightarrow B}}
{\phi\!\left(\frac{h_t^A}{\sqrt{\tau}}\right) z_A},
\quad
S_{A \rightarrow B}
=
\sum_{j=1}^{|\mathcal{V}|}
\phi\!\left(\frac{w_{A,j}^{out}}{\sqrt{\tau}}\right)^\top w_{B,j}^{in},
\quad
z_A
=
\sum_{j=1}^{|\mathcal{V}|}
\phi\!\left(\frac{w_{A,j}^{out}}{\sqrt{\tau}}\right)^\top.
\label{eq:linear_alignment}
\end{equation}

The two sums in Eq.~\ref{eq:linear_alignment} depend only on the fixed matrices $W_A^{out}$ and $W_B^{in}$ which can be computed once and cached offline. Hence, each alignment query $h_t^A$ only needs to be projected through $\phi$ and matched against the cached sums, without repeating the full-vocabulary computation at every generation step. Moreover, since $W_A^{out}$ and $W_B^{in}$ can be directly retrieved from the model, this precomputation itself requires no additional training. As a result, the per-step complexity of the alignment is reduced from $O(|\mathcal{V}|(d_A+d_B))$ to $O(m(d_A+d_B))$ with $m\ll|\mathcal{V}|$, yielding significant speedup in inference time.
\end{comment}

%\xiyan{(TODO) The same alignment operator can be used for latent feedback and communication among agents sharing a model checkpoint. With different source and target checkpoints, it instead defines a cross-model interface, provided that their vocabularies are compatible.}

\section{Experiments}
\subsection{Setup}
\textbf{Models and Task configurations.} Our main experiments use a sequential multi-agent system consisting of a Planner, a Critic, a Refiner, and a Judger. All agents share either Qwen3-8B or Qwen3-14B \citep{qwen3}. Information is passed between successive agents, with the Judger producing the final answer. We additionally evaluate a hierarchical configuration. More task configurations are provided in Appendix \ref{app:config_task}.
%with its setup and results provided in Appendix~\ref{app:hierarchical}. 
To assess alignment across heterogeneous agents while controlling for the effects of multi-agent interaction, we use a two agent Planner and Judger configuration in which the two agents use Qwen3-8B and Qwen3-14B in both directions. 

\noindent\textbf{Stopping configurations.}
We evaluate two mechanisms for controlling latent reasoning: fixed-budget reasoning and entropy-based stopping. For reasoning with fixed budget, we select the latent-step budget for each configuration on a development set from $\{0,10,20,40,80\}$ following \citet{latentmas}. For entropy-based stopping, we terminate latent reasoning when the threshold is tuned using Prototype Align's development-set entropy trajectories following \citep{soft_thinking}. Details are provided in the Appendix~\ref{app:latent_step}. 

\noindent\textbf{Baselines.}
We compare prototype align against other training-free alignment baselines:
i) \emph{Linear Alignment}, which applies a vocabulary-derived linear map \citep{latentmas};
ii) \emph{Exact-Soft Alignment}, which computes the expected receiver embedding under the sender's full next-token distribution (Eq.~\ref{eq:distribution_transition}) \citep{soft_thinking}; iii) \emph{Filtered-Soft Align}, which applies top-$k$/top-$p$ filtering before embedding aggregation of Exact-Soft Alignment.
%\footnote{\citet{soft_thinking} additionally apply top-$k$/top-$p$ filtering to the distribution. We use the unfiltered form to isolate the effect of the kernel approximation. Filtering also requires the full vocabulary logits, which Prototype Align is designed to avoid.}
We include the text baselines for reference, in which agents communicate through generated natural-language messages.\footnote{\textit{Text} and \textit{Linear Align} are published results from LatentMAS~\citep{latentmas}. 
%Our reproduction results are provided in Appendix~\ref{app:reproduction}.
}

%\paragraph{Models.} We evaluate kernelized alignment across multiple model families and scales. For homogeneous multi-agent systems, all agents share the same model checkpoint, using either Qwen3-8B or Qwen3-14B\citep{qwen3} independently. For heterogeneous multi-agent systems, we evaluate bidirectional communication between Qwen3-8B and Qwen3-14B, which have different hidden dimensions but compatible vocabularies.

\noindent\textbf{Benchmarks.} Our primary evaluation covers six benchmarks across three
domains: i) mathematical reasoning with AIME 2024 and AIME 2025 \citep{aime2024,aime2025}; ii) code generation with HumanEval+ and MBPP+ \citep{evalplus}; and iii) scientific and medical question answering with  GPQA-Diamond, and MedQA \citep{gpqa,medqa}. 
We additionally evaluate on GSM8K \citep{gsm8k}, ARC-Easy and ARC-Challenge \citep{arc} as supplementary benchmarks.
%in Appendix \ref{app:supplement}. 

\noindent\textbf{Evaluation.} We measure task performance using the following two metrics: i) \textit{accuracy}, measured by exact match accuracy for mathematical reasoning, answer choice accuracy for scientific and medical question answering, and pass@1 for code generation. 
%ii) \textit{token}, which reports the average number of tokens generated by the final Judger, excluding input tokens, intermediate-agent outputs, and latent reasoning steps.
%ii) \textit{token}, which reports the total number of vocabulary level decoding operations per problem.
%which report the average number of explicit vocabulary decoding operations per problem. 
%For baseline results, we obtain tokens from the original paper. 
ii) \textit{time}, which we measure end-to-end time and report the average time per problem. 
%For baseline results, we inference time on our hardware for fair comparison. 
To reduce the effect of hardware differences, we estimate the inference time of baselines using published results of baselines and reference throughput measure on our hardware.
Specifically, we use $\widehat{T}_{b,s,m,d} = N_{b,s,m,d}/r_{b,s,m,d}$, where $N_{b,s,m,d}$ is the average token count per problem recorded in the published logs for baseline $b$, multi-agent setting $s$, model $m$, and benchmark $d$, and $r_{b,s,m,d}$ is the reference throughput, measured in tokens per second on our hardware under the corresponding model and benchmark config. Details are in Appendix \ref{app:inference_efficiency}.

%To assess computational efficiency, we report the number of explicit vocabulary decoding operations and end-to-end wall-clock inference time.



\subsection{Main Results}

\begin{table*}[t]
\centering
\setlength{\tabcolsep}{4pt}
\renewcommand{\arraystretch}{1.12}
\small
\resizebox{\textwidth}{!}{%
\begin{tabular}{@{}ll c ccc>{\columncolor{oursbg}}c cc>{\columncolor{oursbg}}c@{}}
\toprule
& & & \multicolumn{4}{c}{Fixed Steps} & \multicolumn{3}{c}{Entropy Stopping} \\
\cmidrule(lr){4-7}\cmidrule(lr){8-10}
 & Data & Text$^\dagger$ & \makecell{Linear\\Align$^\dagger$} & \makecell{Exact-Soft\\Align} & \makecell{Filtered-Soft\\Align} & \makecell{\textbf{Prototype}\\\textbf{Align}} & \makecell{Exact-Soft\\Align} & \makecell{Filtered-Soft\\Align} & \makecell{\textbf{Prototype}\\\textbf{Align}} \\
\midrule
\multirow{7}{*}{\rotatebox{90}{Qwen3-8B}} & AIME 2024 & 53.30 & 56.70 & 53.33\tabsd{0.00} & 53.33\tabsd{8.16} & 68.89\tabsd{1.57} & 62.22\tabsd{6.85} & \textbf{72.22}\tabsd{1.57} & 65.56\tabsd{1.57} \\
& AIME 2025 & 53.30 & 53.30 & 40.00\tabsd{5.44} & 48.89\tabsd{1.57} & 54.44\tabsd{4.16} & 52.22\tabsd{8.31} & 43.33\tabsd{0.00} & \textbf{57.78}\tabsd{4.16} \\
& HumanEval+ & 80.50 & 80.50 & 77.44\tabsd{1.00} & 78.66\tabsd{1.72} & 83.94\tabsd{0.57} & 83.13\tabsd{2.24} & \textbf{84.15}\tabsd{0.86} & 83.94\tabsd{1.60} \\
& MBPP+ & 69.50 & 74.60 & 73.02\tabsd{0.57} & 71.78\tabsd{0.45} & 75.40\tabsd{0.22} & 72.49\tabsd{0.75} & 74.69\tabsd{1.19} & \textbf{75.49}\tabsd{0.50} \\
& \makecell[l]{GPQA-Diamond} & 43.40 & 45.50 & 45.12\tabsd{2.68} & 46.63\tabsd{2.08} & 49.16\tabsd{1.04} & 48.48\tabsd{1.24} & 41.58\tabsd{3.15} & \textbf{50.17}\tabsd{1.04} \\
& MedQA & 75.00 & \textbf{75.30} & 71.33\tabsd{2.16} & 70.33\tabsd{0.82} & 73.22\tabsd{0.87} & 68.67\tabsd{1.44} & 68.22\tabsd{0.57} & 73.22\tabsd{1.29} \\
\cmidrule(lr){2-10}
& \textbf{AVG} & 62.50 & 64.32 & 60.04 & 61.60 & 67.51 & 64.54 & 64.03 & \textbf{67.69} \\
\midrule

\multirow{7}{*}{\rotatebox{90}{Qwen3-14B}} & AIME 2024 & 63.30 & 66.70 & 70.00\tabsd{2.72} & 68.89\tabsd{3.14} & \textbf{75.56}\tabsd{3.14} & 68.89\tabsd{5.67} & 65.56\tabsd{3.14} & 73.33\tabsd{2.72} \\
& AIME 2025 & 60.00 & \textbf{63.30} & 57.78\tabsd{3.14} & 55.56\tabsd{3.14} & 60.00\tabsd{2.72} & 54.44\tabsd{3.14} & 53.33\tabsd{2.72} & 58.89\tabsd{5.67} \\
& HumanEval+ & 81.10 & 86.50 & 88.82\tabsd{1.04} & 87.80\tabsd{0.50} & 88.21\tabsd{0.57} & \textbf{91.26}\tabsd{1.44} & 88.41\tabsd{0.50} & 84.55\tabsd{0.29} \\
& MBPP+ & 72.80 & 75.70 & 77.87\tabsd{1.47} & 77.69\tabsd{0.62} & 78.13\tabsd{0.12} & \textbf{79.98}\tabsd{0.12} & 79.54\tabsd{0.25} & 78.22\tabsd{0.33} \\
& \makecell[l]{GPQA-Diamond} & 51.50 & 52.00 & 54.55\tabsd{2.51} & 55.89\tabsd{1.04} & \textbf{56.23}\tabsd{0.48} & 52.53\tabsd{1.24} & 53.70\tabsd{0.86} & 54.71\tabsd{3.10} \\
& MedQA & 80.30 & \textbf{80.70} & 78.44\tabsd{1.10} & 78.78\tabsd{1.50} & 80.67\tabsd{0.72} & 78.11\tabsd{0.63} & 78.33\tabsd{1.66} & 79.22\tabsd{0.79} \\
\cmidrule(lr){2-10}
& \textbf{AVG} & 68.17 & 70.82 & 71.24 & 70.77 & \textbf{73.13} & 70.87 & 69.81 & 71.49 \\
\bottomrule
\end{tabular}}
\caption{Accuracy (\%, $\uparrow$) of Qwen3-8B and Qwen3-14B under the sequential configuration (mean\,$\pm$\,population std over available runs). Code-generation benchmarks use pass@1. AVG denotes the unweighted mean across all six benchmarks and is reported only when all six results are available. \textbf{Bold}: best among available latent alignment configurations within each model and benchmark. $^\dagger$Published results from LatentMAS~\citep{latentmas}.}
\label{tab:main_results_acc}
\end{table*}



\begin{comment}
\begin{table*}[t]
\centering
\caption{Performance and efficiency of Qwen3-8B and Qwen3-14B under the sequential configuration (mean\,$\pm$\,population std over available runs). \textit{Acc} (\%, $\uparrow$) denotes task accuracy, \textit{Token} ($\downarrow$) the total number of vocabulary level decoding operations per problem, and \textit{Time} ($\downarrow$) inference time in seconds per problem. \textit{Time} of Text and Linear Align is converted to the throughput of our machine (Appendix~\ref{app:inference_efficiency}). \textbf{Bold}: best among available latent alignment methods. $^\dagger$Published results from LatentMAS~\citep{latentmas}.}
\label{tab:main_results}
\setlength{\tabcolsep}{3pt}
\renewcommand{\arraystretch}{1.08}
\small
\resizebox{\textwidth}{!}{%
\begin{tabular}{@{}ll cccc>{\columncolor{oursbg}}c cc>{\columncolor{oursbg}}c cccc>{\columncolor{oursbg}}c cc>{\columncolor{oursbg}}c@{}}
\toprule
& & \multicolumn{8}{c}{Qwen3-8B} & \multicolumn{8}{c}{Qwen3-14B} \\
\cmidrule(lr){3-10}\cmidrule(l){11-18}
& && \multicolumn{4}{c}{Fixed Steps} & \multicolumn{3}{c}{Entropy Stopping} & &\multicolumn{4}{c}{Fixed Steps} & \multicolumn{3}{c}{Entropy Stopping} \\
\cmidrule(lr){4-7}\cmidrule(lr){8-10}\cmidrule(lr){12-15}\cmidrule(lr){16-18}
Data & Metric
& Text$^\dagger$
& \makecell{Linear\\Align$^\dagger$}
& \makecell{Exact-Soft\\fixed steps}
& \makecell{Filtered-Soft\\fixed steps}
& \raisebox{-0.3\height}{\shortstack{\textbf{Kernel}\\\textbf{Align}}}
& \makecell{Exact-Soft\\Align}
& \makecell{Filtered-Soft\\Align}
& \makecell{Kernel-Align\\entropy}
& Text$^\dagger$
& \makecell{Linear\\Align$^\dagger$}
& \makecell{Exact-Soft\\fixed steps}
& \makecell{Filtered-Soft\\fixed steps}
& \raisebox{-0.3\height}{\shortstack{\textbf{Kernel}\\\textbf{Align}}}
& \makecell{Exact-Soft\\Align}
& \makecell{Filtered-Soft\\Align}
& \makecell{Kernel-Align\\entropy} \\
\midrule
\multirow{3}{*}{\makecell[l]{AIME\\2024}}
 & Acc
& 53.30 & 56.70 & 53.33\tabsd{0.00} & 53.33\tabsd{8.16} & 68.89\tabsd{1.57} & 62.22\tabsd{6.85} & \textbf{72.22}\tabsd{1.57} & 65.56\tabsd{1.57}
& 63.30 & 66.70 & 70.00\tabsd{2.72} & 70.00\tabsd{3.33}\textsuperscript{($n=2$)} & \textbf{75.56}\tabsd{3.14} & 68.89\tabsd{5.67} & 65.56\tabsd{3.14} & 73.33\tabsd{2.72} \\
 & Token
& 38{,}596 & 8{,}953 & \textbf{5{,}910}\tabsd{331} & 6{,}514\tabsd{387} & 8{,}548\tabsd{344} & 22{,}496\tabsd{1178} & 17{,}850\tabsd{1100} & 8{,}552\tabsd{586}
& 32{,}092 & 10{,}593 & 8{,}233\tabsd{238} & \textbf{8{,}031}\tabsd{69} & 8{,}045\tabsd{310} & 19{,}493\tabsd{709} & 20{,}564\tabsd{743} & 8{,}776.7\tabsd{337.0} \\
 & Time
& 1177.2 & 269.3 & \textbf{187.7}\tabsd{13.6} & 194.1\tabsd{15.9} & 229.3\tabsd{3.3} & 575.3\tabsd{33.7} & 361.3\tabsd{50.6} & 216.0\tabsd{13.2}
& 993.2 & 327.8 & 247.6\tabsd{22.9} & 253.7\tabsd{1.6} & \textbf{243.3}\tabsd{5.9} & 515.1\tabsd{9.8} & 492.9\tabsd{9.4} & 272.9\tabsd{18.8} \\
\addlinespace[2pt]\cmidrule(lr){1-18}\addlinespace[1pt]
\multirow{3}{*}{\makecell[l]{AIME\\2025}}
 & Acc
& 53.30 & 53.30 & 40.00\tabsd{5.44} & 48.89\tabsd{1.57} & 54.44\tabsd{4.16} & 52.22\tabsd{8.31} & 43.33\tabsd{0.00} & \textbf{57.78}\tabsd{4.16}
& 60.00 & \textbf{63.30} & 57.78\tabsd{3.14} & -- & 60.00\tabsd{2.72} & 54.44\tabsd{3.14} & 53.33\tabsd{2.72} & 58.89\tabsd{5.67} \\
 & Token
& 45{,}088 & 8{,}699 & \textbf{8{,}396}\tabsd{305} & 8{,}527\tabsd{226} & 10{,}230\tabsd{286} & 26{,}193\tabsd{1716} & 19{,}930\tabsd{485} & 9{,}419\tabsd{427}
& 44{,}618 & 11{,}402 & 10{,}192\tabsd{343} & -- & \textbf{10{,}076}\tabsd{314} & 20{,}402\tabsd{151} & 21{,}101\tabsd{374} & 10{,}746.7\tabsd{183.8} \\
 & Time
& 1423.9 & \textbf{232.0} & 256.0\tabsd{4.8} & 250.2\tabsd{6.8} & 271.2\tabsd{5.8} & 634.6\tabsd{18.5} & 438.7\tabsd{30.6} & 244.2\tabsd{18.6}
& 1458.1 & 372.6 & \textbf{323.4}\tabsd{12.0} & -- & 328.7\tabsd{9.6} & 519.1\tabsd{13.7} & 495.7\tabsd{29.9} & 379.0\tabsd{17.8} \\
\addlinespace[2pt]\cmidrule(lr){1-18}\addlinespace[1pt]
\multirow{3}{*}{\makecell[l]{Human\\Eval+}}
 & Acc
& 80.50 & 80.50 & 77.44\tabsd{1.00} & 78.66\tabsd{1.72} & 83.94\tabsd{0.57} & 83.13\tabsd{2.24} & \textbf{84.15}\tabsd{0.86} & 83.94\tabsd{1.60}
& 81.10 & 86.50 & 88.82\tabsd{1.04} & -- & 88.21\tabsd{0.57} & \textbf{91.26}\tabsd{1.44} & -- & 84.55\tabsd{0.29} \\
 & Token
& 4{,}593 & 1{,}866 & 958\tabsd{57} & \textbf{951}\tabsd{33} & 1{,}639\tabsd{42} & 14{,}598\tabsd{308} & 15{,}540\tabsd{459} & 1{,}692\tabsd{34}
& 5{,}934 & 2{,}042 & \textbf{1{,}563}\tabsd{15} & -- & 1{,}598\tabsd{42} & 19{,}784\tabsd{655} & -- & 2{,}203.6\tabsd{68.8} \\
 & Time
& 65.2 & 25.0 & \textbf{17.1}\tabsd{0.7} & 17.2\tabsd{1.0} & 18.9\tabsd{0.2} & 107.1\tabsd{7.4} & 105.9\tabsd{3.3} & 21.5\tabsd{0.4}
& 79.7 & 27.4 & \textbf{23.5}\tabsd{0.6} & -- & 24.1\tabsd{0.8} & 158.2\tabsd{7.1} & -- & 29.0\tabsd{0.9} \\
\addlinespace[2pt]\cmidrule(lr){1-18}\addlinespace[1pt]
\multirow{3}{*}{\makecell[l]{MBPP+}}
 & Acc
& 69.50 & 74.60 & 73.02\tabsd{0.57} & 71.78\tabsd{0.45} & 75.40\tabsd{0.22} & 72.49\tabsd{0.75} & 74.69\tabsd{1.19} & \textbf{75.49}\tabsd{0.50}
& 72.80 & 75.70 & 77.87\tabsd{1.47} & -- & 78.13\tabsd{0.12} & \textbf{79.98}\tabsd{0.12} & 79.54\tabsd{0.25} & 78.22\tabsd{0.33} \\
 & Token
& 3{,}695 & 1{,}164 & 795\tabsd{10} & \textbf{710}\tabsd{10} & 1{,}463\tabsd{19} & 14{,}297\tabsd{370} & 13{,}755\tabsd{454} & 1{,}532\tabsd{11}
& 4{,}971 & 1{,}621 & \textbf{1{,}314}\tabsd{3} & -- & 1{,}339\tabsd{18} & 21{,}716\tabsd{12} & 21{,}440\tabsd{336} & 1{,}920.3\tabsd{23.5} \\
 & Time
& 56.7 & 17.6 & 14.4\tabsd{0.7} & \textbf{13.0}\tabsd{0.3} & 18.1\tabsd{0.2} & 101.6\tabsd{6.9} & 87.3\tabsd{2.5} & 19.6\tabsd{0.4}
& 76.2 & 24.9 & 22.0\tabsd{0.3} & -- & \textbf{21.4}\tabsd{0.2} & 166.3\tabsd{3.4} & 147.2\tabsd{2.2} & 26.8\tabsd{0.5} \\
\addlinespace[2pt]\cmidrule(lr){1-18}\addlinespace[1pt]
\multirow{3}{*}{\makecell[l]{GPQA-\\Diamond}}
 & Acc
& 43.40 & 45.50 & 45.12\tabsd{2.68} & 46.63\tabsd{2.08} & 49.16\tabsd{1.04} & 48.48\tabsd{1.24} & 41.58\tabsd{3.15} & \textbf{50.17}\tabsd{1.04}
& 51.50 & 52.00 & 54.55\tabsd{2.51} & -- & \textbf{56.23}\tabsd{0.48} & 52.53\tabsd{1.24} & 53.70\tabsd{0.86} & 54.71\tabsd{3.10} \\
 & Token
& 17{,}986 & 4{,}571 & 3{,}301\tabsd{59} & \textbf{3{,}091}\tabsd{63} & 3{,}574\tabsd{87} & 13{,}058\tabsd{315} & 19{,}239\tabsd{502} & 3{,}865\tabsd{77}
& 12{,}676 & 5{,}454 & \textbf{2{,}833}\tabsd{72} & -- & 2{,}999\tabsd{55} & 18{,}443\tabsd{848} & 18{,}329\tabsd{729} & 3{,}499.3\tabsd{32.4} \\
 & Time
& 293.9 & 83.6 & 57.3\tabsd{1.6} & \textbf{56.3}\tabsd{2.6} & 57.6\tabsd{0.6} & 257.4\tabsd{10.7} & 242.5\tabsd{8.4} & 60.0\tabsd{0.6}
& 250.2 & 107.7 & 63.0\tabsd{0.8} & -- & \textbf{59.2}\tabsd{1.8} & 229.7\tabsd{21.1} & 196.7\tabsd{14.9} & 67.0\tabsd{1.4} \\
\addlinespace[2pt]\cmidrule(lr){1-18}\addlinespace[1pt]
\multirow{3}{*}{\makecell[l]{MedQA}}
 & Acc
& 75.00 & \textbf{75.30} & 71.33\tabsd{2.16} & 70.33\tabsd{0.82} & 73.22\tabsd{0.87} & 68.67\tabsd{1.44} & 68.22\tabsd{0.57} & 73.22\tabsd{1.29}
& 80.30 & \textbf{80.70} & 78.44\tabsd{1.10} & -- & 80.67\tabsd{0.72} & 78.11\tabsd{0.63} & 78.33\tabsd{1.66} & 79.22\tabsd{0.79} \\
 & Token
& 4{,}260 & 1{,}555 & 1{,}011\tabsd{28} & \textbf{991}\tabsd{32} & 1{,}163\tabsd{20} & 15{,}480\tabsd{383} & 14{,}942\tabsd{451} & 1{,}148\tabsd{8}
& 3{,}444 & 1{,}841 & \textbf{821}\tabsd{15} & -- & 896\tabsd{22} & 22{,}480\tabsd{342} & 21{,}896\tabsd{191} & 1{,}509.5\tabsd{21.4} \\
 & Time
& 39.5 & 17.7 & 12.8\tabsd{0.7} & \textbf{12.4}\tabsd{0.8} & 13.3\tabsd{0.5} & 107.6\tabsd{2.7} & 99.0\tabsd{4.4} & 13.4\tabsd{0.1}
& 41.6 & 22.2 & \textbf{10.3}\tabsd{0.4} & -- & 12.7\tabsd{0.5} & 147.7\tabsd{3.9} & 129.1\tabsd{0.9} & 16.1\tabsd{0.2} \\
\midrule
\multirow{3}{*}{\textbf{AVG}}
 & Acc
& 62.50 & 64.32 & 60.04 & 61.60 & 67.51 & 64.54 & 64.03 & \textbf{67.69}
& 68.17 & 70.82 & 71.24 & -- & \textbf{73.13} & 70.87 & -- & 71.49 \\
 & Token
& 19{,}036 & 4{,}468 & \textbf{3{,}395} & 3{,}464 & 4{,}436 & 17{,}687 & 16{,}876 & 4{,}368
& 17{,}289 & 5{,}492 & \textbf{4{,}159} & -- & \textbf{4{,}159} & 20{,}386 & -- & 4{,}776.0 \\
 & Time
& 509.4 & 107.5 & 90.9 & \textbf{90.5} & 101.4 & 297.3 & 222.4 & 95.8
& 483.2 & 147.1 & 115.0 & -- & \textbf{114.9} & 289.4 & -- & 131.8 \\
\bottomrule
\end{tabular}}
\par\smallskip
%\parbox{\textwidth}{\footnotesize\raggedright Filtered-Soft uses top-$k=30$, top-$p=0.95$, and top-$n=10$. Kernel-Align entropy uses entropy-based stopping; the default Prototype Align uses fixed steps. Acc superscripts $(n=1)$ or $(n=2)$ indicate the repetition count for all three metrics; otherwise $n=3$. Single-run results have no standard deviation. Token is \texttt{output\_tokens}/\texttt{processed}; Time is \texttt{seconds\_per\_sample}. Fixed-step budgets vary by task and model (recorded in the source comments). AVG is the unweighted mean over all six tasks and is shown only when all six are available. Missing results or incomplete task coverage are marked --.}
\end{table*}
\end{comment}

\begin{comment}
\begin{table*}[t]
\centering
\caption{Performance and efficiency of Qwen3-8B and Qwen3-14B under the sequential configuration (mean\,$\pm$\,population std over available runs). \textit{Acc} (\%, $\uparrow$) denotes task accuracy, \textit{Token} ($\downarrow$) the total number of vocabulary level decoding operations per problem, and \textit{Time} ($\downarrow$) inference time in seconds per problem. \textit{Time} of Text and Linear Align is converted to the throughput of our machine (Appendix~\ref{app:inference_efficiency}). \textbf{Bold}: best among available latent alignment methods. $^\dagger$Published results from LatentMAS~\citep{latentmas}.}
\label{tab:main_results}
\setlength{\tabcolsep}{3pt}
\renewcommand{\arraystretch}{1.08}
\small
\resizebox{\textwidth}{!}{%
\begin{tabular}{@{}ll ccc>{\columncolor{oursbg}}c ccc ccc>{\columncolor{oursbg}}c c ccc@{}}
\toprule
& & \multicolumn{7}{c}{Qwen3-8B} & \multicolumn{8}{c}{Qwen3-14B} \\
\cmidrule(lr){3-9}\cmidrule(l){10-17}
Data & Metric
& Text$^\dagger$
& \makecell{Linear\\Align$^\dagger$}
& \makecell{Exact-Soft\\Align}
& \raisebox{-0.3\height}{\shortstack{\textbf{Kernel}\\\textbf{Align}}}
& \makecell{Filtered-Soft\\Align}
& \makecell{Filtered-Soft\\fixed steps}
& \makecell{Exact-Soft\\fixed steps}
& Text$^\dagger$
& \makecell{Linear\\Align$^\dagger$}
& \makecell{Exact-Soft\\Align}
& \raisebox{-0.3\height}{\shortstack{\textbf{Kernel}\\\textbf{Align}}}
& \makecell{Kernel-Align\\entropy}
& \makecell{Filtered-Soft\\Align}
& \makecell{Filtered-Soft\\fixed steps}
& \makecell{Exact-Soft\\fixed steps} \\
\midrule
\multirow{3}{*}{\makecell[l]{AIME\\2024}}
 & Acc
& 53.30 & 56.70 & 62.22\tabsd{6.85} & 68.89\tabsd{1.57} & \textbf{72.22}\tabsd{1.57} & 53.33\tabsd{8.16} & 53.33\tabsd{0.00}
& 63.30 & 66.70 & 68.89\tabsd{5.67} & \textbf{75.56}\tabsd{3.14} & 73.33\tabsd{2.72} & 65.56\tabsd{3.14} & 70.00\tabsd{3.33}\textsuperscript{($n=2$)} & 70.00\tabsd{2.72} \\
 & Token
& 38{,}596 & 8{,}953 & 22{,}496\tabsd{1178} & 8{,}548\tabsd{344} & 17{,}850\tabsd{1100} & 6{,}514\tabsd{387} & \textbf{5{,}910}\tabsd{331}
& 32{,}092 & 10{,}593 & 19{,}493\tabsd{709} & 8{,}045\tabsd{310} & 8{,}776.7\tabsd{337.0} & 20{,}564\tabsd{743} & \textbf{8{,}031}\tabsd{69} & 8{,}233\tabsd{238} \\
 & Time
& 1177.2 & 269.3 & 575.3\tabsd{33.7} & 229.3\tabsd{3.3} & 361.3\tabsd{50.6} & 194.1\tabsd{15.9} & \textbf{187.7}\tabsd{13.6}
& 993.2 & 327.8 & 515.1\tabsd{9.8} & \textbf{243.3}\tabsd{5.9} & 272.9\tabsd{18.8} & 492.9\tabsd{9.4} & 253.7\tabsd{1.6} & 247.6\tabsd{22.9} \\
\addlinespace[2pt]\cmidrule(lr){1-17}\addlinespace[1pt]
\multirow{3}{*}{\makecell[l]{AIME\\2025}}
 & Acc
& 53.30 & 53.30 & 52.22\tabsd{8.31} & \textbf{54.44}\tabsd{4.16} & 43.33\tabsd{0.00} & 48.89\tabsd{1.57} & 40.00\tabsd{5.44}
& 60.00 & \textbf{63.30} & 54.44\tabsd{3.14} & 60.00\tabsd{2.72} & 58.89\tabsd{5.67} & 53.33\tabsd{2.72} & -- & 57.78\tabsd{3.14} \\
 & Token
& 45{,}088 & 8{,}699 & 26{,}193\tabsd{1716} & 10{,}230\tabsd{286} & 19{,}930\tabsd{485} & 8{,}527\tabsd{226} & \textbf{8{,}396}\tabsd{305}
& 44{,}618 & 11{,}402 & 20{,}402\tabsd{151} & \textbf{10{,}076}\tabsd{314} & 10{,}746.7\tabsd{183.8} & 21{,}101\tabsd{374} & -- & 10{,}192\tabsd{343} \\
 & Time
& 1423.9 & \textbf{232.0} & 634.6\tabsd{18.5} & 271.2\tabsd{5.8} & 438.7\tabsd{30.6} & 250.2\tabsd{6.8} & 256.0\tabsd{4.8}
& 1458.1 & 372.6 & 519.1\tabsd{13.7} & 328.7\tabsd{9.6} & 379.0\tabsd{17.8} & 495.7\tabsd{29.9} & -- & \textbf{323.4}\tabsd{12.0} \\
\addlinespace[2pt]\cmidrule(lr){1-17}\addlinespace[1pt]
\multirow{3}{*}{\makecell[l]{Human\\Eval+}}
 & Acc
& 80.50 & 80.50 & 83.13\tabsd{2.24} & 83.94\tabsd{0.57} & \textbf{84.15}\tabsd{0.86} & 78.66\tabsd{1.72} & 77.44\tabsd{1.00}
& 81.10 & 86.50 & \textbf{91.26}\tabsd{1.44} & 88.21\tabsd{0.57} & 84.55\tabsd{0.29} & -- & -- & 88.82\tabsd{1.04} \\
 & Token
& 4{,}593 & 1{,}866 & 14{,}598\tabsd{308} & 1{,}639\tabsd{42} & 15{,}540\tabsd{459} & \textbf{951}\tabsd{33} & 958\tabsd{57}
& 5{,}934 & 2{,}042 & 19{,}784\tabsd{655} & 1{,}598\tabsd{42} & 2{,}203.6\tabsd{68.8} & -- & -- & \textbf{1{,}563}\tabsd{15} \\
 & Time
& 65.2 & 25.0 & 107.1\tabsd{7.4} & 18.9\tabsd{0.2} & 105.9\tabsd{3.3} & 17.2\tabsd{1.0} & \textbf{17.1}\tabsd{0.7}
& 79.7 & 27.4 & 158.2\tabsd{7.1} & 24.1\tabsd{0.8} & 29.0\tabsd{0.9} & -- & -- & \textbf{23.5}\tabsd{0.6} \\
\addlinespace[2pt]\cmidrule(lr){1-17}\addlinespace[1pt]
\multirow{3}{*}{\makecell[l]{MBPP+}}
 & Acc
& 69.50 & 74.60 & 72.49\tabsd{0.75} & \textbf{75.40}\tabsd{0.22} & 74.69\tabsd{1.19} & 71.78\tabsd{0.45} & 73.02\tabsd{0.57}
& 72.80 & 75.70 & \textbf{79.98}\tabsd{0.12} & 78.13\tabsd{0.12} & 78.22\tabsd{0.33} & 79.54\tabsd{0.25} & -- & 77.87\tabsd{1.47} \\
 & Token
& 3{,}695 & 1{,}164 & 14{,}297\tabsd{370} & 1{,}463\tabsd{19} & 13{,}755\tabsd{454} & \textbf{710}\tabsd{10} & 795\tabsd{10}
& 4{,}971 & 1{,}621 & 21{,}716\tabsd{12} & 1{,}339\tabsd{18} & 1{,}920.3\tabsd{23.5} & 21{,}440\tabsd{336} & -- & \textbf{1{,}314}\tabsd{3} \\
 & Time
& 56.7 & 17.6 & 101.6\tabsd{6.9} & 18.1\tabsd{0.2} & 87.3\tabsd{2.5} & \textbf{13.0}\tabsd{0.3} & 14.4\tabsd{0.7}
& 76.2 & 24.9 & 166.3\tabsd{3.4} & \textbf{21.4}\tabsd{0.2} & 26.8\tabsd{0.5} & 147.2\tabsd{2.2} & -- & 22.0\tabsd{0.3} \\
\addlinespace[2pt]\cmidrule(lr){1-17}\addlinespace[1pt]
\multirow{3}{*}{\makecell[l]{GPQA-\\Diamond}}
 & Acc
& 43.40 & 45.50 & 48.48\tabsd{1.24} & \textbf{49.16}\tabsd{1.04} & 41.58\tabsd{3.15} & 46.63\tabsd{2.08} & 45.12\tabsd{2.68}
& 51.50 & 52.00 & 52.53\tabsd{1.24} & \textbf{56.23}\tabsd{0.48} & 54.71\tabsd{3.10} & 53.70\tabsd{0.86} & -- & 54.55\tabsd{2.51} \\
 & Token
& 17{,}986 & 4{,}571 & 13{,}058\tabsd{315} & 3{,}574\tabsd{87} & 19{,}239\tabsd{502} & \textbf{3{,}091}\tabsd{63} & 3{,}301\tabsd{59}
& 12{,}676 & 5{,}454 & 18{,}443\tabsd{848} & 2{,}999\tabsd{55} & 3{,}499.3\tabsd{32.4} & 18{,}329\tabsd{729} & -- & \textbf{2{,}833}\tabsd{72} \\
 & Time
& 293.9 & 83.6 & 257.4\tabsd{10.7} & 57.6\tabsd{0.6} & 242.5\tabsd{8.4} & \textbf{56.3}\tabsd{2.6} & 57.3\tabsd{1.6}
& 250.2 & 107.7 & 229.7\tabsd{21.1} & \textbf{59.2}\tabsd{1.8} & 67.0\tabsd{1.4} & 196.7\tabsd{14.9} & -- & 63.0\tabsd{0.8} \\
\addlinespace[2pt]\cmidrule(lr){1-17}\addlinespace[1pt]
\multirow{3}{*}{\makecell[l]{MedQA}}
 & Acc
& 75.00 & \textbf{75.30} & 68.67\tabsd{1.44} & 73.22\tabsd{0.87} & 68.22\tabsd{0.57} & 70.33\tabsd{0.82} & 71.33\tabsd{2.16}
& 80.30 & \textbf{80.70} & 78.11\tabsd{0.63} & 80.67\tabsd{0.72} & 79.22\tabsd{0.79} & 78.33\tabsd{1.66} & -- & 78.44\tabsd{1.10} \\
 & Token
& 4{,}260 & 1{,}555 & 15{,}480\tabsd{383} & 1{,}163\tabsd{20} & 14{,}942\tabsd{451} & \textbf{991}\tabsd{32} & 1{,}011\tabsd{28}
& 3{,}444 & 1{,}841 & 22{,}480\tabsd{342} & 896\tabsd{22} & 1{,}509.5\tabsd{21.4} & 21{,}896\tabsd{191} & -- & \textbf{821}\tabsd{15} \\
 & Time
& 39.5 & 17.7 & 107.6\tabsd{2.7} & 13.3\tabsd{0.5} & 99.0\tabsd{4.4} & \textbf{12.4}\tabsd{0.8} & 12.8\tabsd{0.7}
& 41.6 & 22.2 & 147.7\tabsd{3.9} & 12.7\tabsd{0.5} & 16.1\tabsd{0.2} & 129.1\tabsd{0.9} & -- & \textbf{10.3}\tabsd{0.4} \\
\midrule
\multirow{3}{*}{\textbf{AVG}}
 & Acc
& 62.50 & 64.32 & 64.54 & \textbf{67.51} & 64.03 & 61.60 & 60.04
& 68.17 & 70.82 & 70.87 & \textbf{73.13} & 71.49 & -- & -- & 71.24 \\
 & Token
& 19{,}036 & 4{,}468 & 17{,}687 & 4{,}436 & 16{,}876 & 3{,}464 & \textbf{3{,}395}
& 17{,}289 & 5{,}492 & 20{,}386 & \textbf{4{,}159} & 4{,}776.0 & -- & -- & \textbf{4{,}159} \\
 & Time
& 509.4 & 107.5 & 297.3 & 101.4 & 222.4 & \textbf{90.5} & 90.9
& 483.2 & 147.1 & 289.4 & \textbf{114.9} & 131.8 & -- & -- & 115.0 \\
\bottomrule
\end{tabular}}
\par\smallskip
\parbox{\textwidth}{\footnotesize\raggedright
Filtered-Soft uses top-$k=30$, top-$p=0.95$, and top-$n=10$.
The three added columns use the latest available run directory per setting.
Acc superscripts $(n=1)$ or $(n=2)$ indicate the repetition count for all three metrics;
otherwise $n=3$. Single-run results have no standard deviation.
Token is \texttt{output\_tokens}/\texttt{processed}; Time is \texttt{seconds\_per\_sample}.
Fixed-step budgets vary by task and model (recorded in the source comments).
AVG is the unweighted mean over all six tasks and is shown only when all six are available.
Missing results or incomplete task coverage are marked --.}
\end{table*}
\end{comment}


\begin{comment}
\begin{table*}[t]
\centering
\caption{\xiyan{Accuracy (\%, $\uparrow$) of Qwen3-8B under sequential and hierarchical configurations. Code-generation benchmarks are evaluated using pass@1. AVG denotes the unweighted mean across all six benchmarks. $\dagger$ denotes the best result across multiple runs, whereas Exact-Soft Align and Prototype Align report mean $\pm$ standard deviation. Bold indicates the highest reported accuracy among training-free latent alignment methods within each configuration and column.}}
\label{tab:main_results}
\setlength{\tabcolsep}{3pt}
\renewcommand{\arraystretch}{1.2}
\small
\resizebox{\textwidth}{!}{%
\xiyan{%
\begin{tabular}{@{}cl*{7}{c}@{}}
\toprule
Stru & Method
& \makecell{AIME\\2024}
& \makecell{AIME\\2025}
& \makecell{Human\\Eval+}
& MBPP+
& \makecell{GPQA-\\Diamond}
& MedQA
& AVG \\
\midrule
\multirow{4}{*}{\rotatebox[origin=c]{90}{Sequential}}
& Text$^{\dagger}$ & 53.30 & 53.30 & 80.50 & 69.50 & 43.40 & 75.00 & 62.50 \\
\cmidrule(lr){2-9}
& Linear Align$^{\dagger}$ & 56.70 & \textbf{53.30} & 80.50 & 74.60 & 45.50 & \textbf{75.30} & 64.32 \\
& Exact-Soft Align
& $63.75\pm\,\phantom{0.00}$
& $49.17\pm\,\phantom{0.00}$
& $82.62\pm\,\phantom{0.00}$
& $73.15\pm\,\phantom{0.00}$
& $\textbf{49.24}\pm\,\phantom{0.00}$
& $66.67\pm\,\phantom{0.00}$
& 64.10 \\
& Prototype Align
& $\textbf{66.25}\pm\,\phantom{0.00}$
& $49.58\pm\,\phantom{0.00}$
& $\textbf{83.54}\pm\,\phantom{0.00}$
& $\textbf{75.26}\pm\,\phantom{0.00}$
& $48.74\pm\,\phantom{0.00}$
& $73.67\pm\,\phantom{0.00}$
& \textbf{66.17} \\
\midrule\midrule
\multirow{4}{*}{\rotatebox[origin=c]{90}{Hierarchical}}
& Text$^{\dagger}$ & 53.30 & 50.00 & 76.80 & 71.90 & 43.00 & 76.30 & 61.88 \\
\cmidrule(lr){2-9}
& Linear Align$^{\dagger}$ & 53.30 & 50.00 & 78.00 & 72.20 & 46.90 & \textbf{77.00} & 62.90 \\
& Exact-Soft Align
& --
& --
& --
& $\textbf{73.21}\pm\,\phantom{0.00}$
& --
& --
& -- \\
& Prototype Align
& $\textbf{60.00}\pm\,\phantom{0.00}$
& $\textbf{53.33}\pm\,\phantom{0.00}$
& $\textbf{83.84}\pm\,\phantom{0.00}$
& $72.75\pm\,\phantom{0.00}$
& $\textbf{51.39}\pm\,\phantom{0.00}$
& $73.33\pm\,\phantom{0.00}$
& \textbf{65.77} \\
\bottomrule
\end{tabular}%
}%
}
\end{table*}

% Requires: booktabs, multirow, makecell, graphicx
\begin{table*}[t]
\centering
\caption{\xiyan{Computational efficiency of Qwen3-8B under sequential and hierarchical configurations. Token ($\downarrow$) denotes tokens per problem, and Time ($\downarrow$) denotes inference time in seconds per problem. Text and Linear Align times are converted to the throughput of the machine used for Prototype Align. AVG denotes the unweighted mean across all six benchmarks and is reported only when all results are available. Bold indicates the lowest value for each metric within each configuration and column. A dash indicates an unavailable result.}}
\label{tab:efficiency_results}
\setlength{\tabcolsep}{3pt}
\renewcommand{\arraystretch}{1.2}
\small
\resizebox{\textwidth}{!}{%
\xiyan{%
\begin{tabular}{@{}cll*{7}{c}@{}}
\toprule
Stru & Method & Metric
& \makecell{AIME\\2024}
& \makecell{AIME\\2025}
& \makecell{Human\\Eval+}
& MBPP+
& \makecell{GPQA-\\Diamond}
& MedQA
& AVG \\
\midrule
\multirow{8}{*}{\rotatebox[origin=c]{90}{Sequential}}
& \multirow{2}{*}{Text}
& Token & 38596 & 45088 & 4593 & 3695 & 17986 & 4260 & 19036.33 \\
& & Time & 1018.6 & 1159.4 & 57.7 & 52.1 & 292.2 & 37.1 & 436.18 \\
\cmidrule(lr){2-10}
& \multirow{2}{*}{Linear Align}
& Token & 8953 & \textbf{8699} & 1866 & \textbf{1164} & 4571 & 1555 & 4468.00 \\
& & Time & 236.3 & \textbf{223.7} & 23.4 & \textbf{16.4} & 74.3 & 13.5 & \textbf{97.93} \\
\cmidrule(lr){2-10}
& \multirow{2}{*}{Exact-Soft Align}
& Token & 22342 & 24557 & 15055 & 14388 & 13076 & 15972 & 17565.00 \\
& & Time & 549.0 & 597.0 & 112.4 & 101.6 & 257.6 & 111.1 & 288.12 \\
\cmidrule(lr){2-10}
& \multirow{2}{*}{Prototype Align}
& Token & \textbf{8548} & 10158 & \textbf{1637} & 1462 & \textbf{3588} & \textbf{1139} & \textbf{4422.00} \\
& & Time & \textbf{225.6} & 261.2 & \textbf{18.8} & 17.9 & \textbf{58.3} & \textbf{12.5} & 99.05 \\
\midrule
\multirow{8}{*}{\rotatebox[origin=c]{90}{Hierarchical}}
& \multirow{2}{*}{Text}
& Token & 42629 & 53929 & 8768 & 7703 & 22450 & 6893 & 23728.67 \\
& & Time & 1109.7 & 1391.2 & 107.5 & 126.8 & 358.4 & 82.5 & 529.35 \\
\cmidrule(lr){2-10}
& \multirow{2}{*}{Linear Align}
& Token & \textbf{7526} & 13230 & \textbf{1274} & \textbf{1264} & \textbf{3395} & \textbf{1007} & \textbf{4616.00} \\
& & Time & \textbf{195.9} & 341.3 & \textbf{15.6} & 20.8 & \textbf{54.2} & \textbf{12.1} & \textbf{106.65} \\
\cmidrule(lr){2-10}
& \multirow{2}{*}{Exact-Soft Align}
& Token & -- & -- & -- & 12768 & -- & -- & -- \\
& & Time & -- & -- & -- & 92.2 & -- & -- & -- \\
\cmidrule(lr){2-10}
& \multirow{2}{*}{Prototype Align}
& Token & 9926 & \textbf{11207} & 1681 & 1502 & 3752 & 1351 & 4903.17 \\
& & Time & 258.4 & \textbf{289.1} & 20.1 & \textbf{20.1} & 59.9 & 15.7 & 110.55 \\
\bottomrule
\end{tabular}%
}%
}
\end{table*}
\end{comment}

% Requires: booktabs, multirow, makecell, graphicx
% Only the number after \pm uses a smaller font.
\providecommand{\smallerr}[1]{}
%\renewcommand{\smallerr}[1]{\ensuremath{\pm}{\scriptsize #1}}
\renewcommand{\smallerr}[1]{{\scriptsize$\pm#1$}}

\begin{comment}
\begin{table*}[t]
\centering
\caption{Performance and computational efficiency of Qwen3-8B under the sequential configuration. Acc ($\uparrow$) denotes task accuracy (\%), and pass@1 (\%) for code generation. Token ($\downarrow$) denotes tokens per problem, and Time ($\downarrow$) denotes inference time (seconds per problem). Text and Linear Align times are converted to the Kernel machine's throughput. AVG denotes the unweighted mean across all six benchmarks. The highest accuracy among training-free latent alignment methods, the lowest token count, and the shortest inference time in each column are shown in bold.}
\label{tab:main_results}
\setlength{\tabcolsep}{3pt}
\renewcommand{\arraystretch}{1.12}
\small
\resizebox{\textwidth}{!}{%%
\begin{tabular}{@{}ll*{7}{c}@{}}
\toprule
Method & Metric
& \makecell{AIME\\2024}
& \makecell{AIME\\2025}
& \makecell{Human\\Eval+}
& MBPP+
& \makecell{GPQA-\\Diamond}
& MedQA
& AVG \\
\midrule
%Single & Acc & 50.00 & 46.70 & 74.40 & 64.80 & 39.90 & 53.00 & 54.80 \\
% & Token & 12891 & 14692 & 2507 & 2053 & 6435 & 2098 & 6779.33 \\
%& Time & 340.2 & 377.8 & 31.5 & 28.9 & 104.6 & 18.3 & 150.22 \\
%\midrule
Text & Acc & 53.30 & 53.30 & 80.50 & 69.50 & 43.40 & 75.00 & 62.50 \\
 & Token & 38596 & 45088 & 4593 & 3695 & 17986 & 4260 & 19036.33 \\
 & Time & 1018.6 & 1159.4 & 57.7 & 52.1 & 292.2 & 37.1 & 436.18 \\
\midrule\midrule
\multirow{3}{*}{\makecell[l]{Linear Align\\\citep{latentmas}}} & Acc & 56.70 & \textbf{53.30} & 80.50 & 74.60 & 45.50 & \textbf{75.30} & 64.32 \\
 & Token & 8953 & \textbf{8699} & 1866 & \textbf{1164} & 4571 & 1555 & 4468.00 \\
 & Time & 236.3 & \textbf{223.7} & 23.4 & \textbf{16.4} & 74.3 & 13.5 & \textbf{97.93} \\
\midrule
\multirow{3}{*}{\makecell[l]{Exact-Soft Align\\\citep{soft_thinking}}} & Acc & 63.75 & 49.17 & 82.62 & 73.15 & \textbf{49.24} & 66.67 & 64.10 \\
 & Token & 22342 & 24557 & 15055 & 14388 & 13076 & 15972 & 17565.00 \\
 & Time & 549.0 & 597.0 & 112.4 & 101.6 & 257.6 & 111.1 & 288.12 \\
\midrule
Prototype Align & Acc & \textbf{66.25} & 49.58 & \textbf{83.54} & \textbf{75.26} & 48.74 & 73.67 & \textbf{66.17} \\
 & Token & \textbf{8548} & 10158 & \textbf{1637} & 1462 & \textbf{3588} & \textbf{1139} & \textbf{4422.00} \\
 & Time & \textbf{225.6} & 261.2 & \textbf{18.8} & 17.9 & \textbf{58.3} & \textbf{12.5} & 99.05 \\
%\midrule
%Kernel-ES Align & Acc & 64.17 & 50.42 & 83.38 & 75.20 & 48.61 & 72.00 & 65.63 \\
% & Token & 8770 & 9830 & 1696 & 1528 & 3846 & 1154 & 4470.67 \\
% & Time & 225.8 & 255.5 & 21.5 & 19.6 & 59.4 & 13.6 & 99.23 \\
\bottomrule
\end{tabular}%
}%
\end{table*}



\begin{table*}[t]
\centering
\caption{\xiyan{Performance and computational efficiency of Qwen3-14B under the sequential configuration. Acc ($\uparrow$) denotes task accuracy (\%), or pass@1 (\%) for code generation; Token ($\downarrow$) denotes tokens per problem; and Time ($\downarrow$) denotes inference time (seconds per problem). Single, Text, and Linear Align times are converted to the Kernel machine's throughput. AVG denotes the unweighted mean across all six benchmarks. The highest accuracy among training-free latent alignment methods, the lowest token count, and the shortest inference time in each column (including AVG) are shown in bold. A dash indicates an unavailable result; AVG is reported only when all six benchmark results are available.}}
\label{tab:seq_14b_results}
\setlength{\tabcolsep}{3pt}
\renewcommand{\arraystretch}{1.12}
\small
\resizebox{\textwidth}{!}{%
\xiyan{%
\begin{tabular}{@{}ll*{7}{c}@{}}
\toprule
Method & Metric
& \makecell{AIME\\2024}
& \makecell{AIME\\2025}
& \makecell{Human\\Eval+}
& MBPP+
& \makecell{GPQA-\\Diamond}
& MedQA
& AVG \\
\midrule
%Single & Acc & 63.30 & 56.70 & 76.80 & 68.50 & 48.50 & 64.70 & 63.08 \\
% & Token & 11263 & 11298 & 2366 & 1858 & 5547 & 1746 & 5679.67 \\
% & Time & 348.6 & 369.2 & 31.8 & 28.5 & 109.5 & 21.1 & 151.45 \\
%\midrule
Text & Acc & 63.30 & 60.00 & 81.10 & 72.80 & 51.50 & 80.30 & 68.17 \\
 & Token & 32092 & 44618 & 5934 & 4971 & 12676 & 3444 & 17289.17 \\
 & Time & 993.2 & 1458.1 & 79.7 & 76.2 & 250.2 & 41.6 & 483.17 \\
\midrule
\multirow{3}{*}{\makecell[l]{Linear Align\\\citep{latentmas}}} & Acc & 66.70 & \textbf{63.30} & 86.50 & 75.70 & 52.00 & 80.70 & 70.82 \\
 & Token & 10593 & 11402 & 2042 & 1621 & 5454 & 1841 & 5492.17 \\
 & Time & 327.8 & 372.6 & 27.4 & 24.9 & 107.7 & 22.2 & 147.10 \\
\midrule
\multirow{3}{*}{\makecell[l]{Exact-Soft Align\\\citep{soft_thinking}}} & Acc & -- & 49.58 & \textbf{90.70} & -- & 53.54 & 79.00 & -- \\
 & Token & -- & 21403 & 19587 & -- & 18528 & 22875 & -- \\
 & Time & -- & 575.5 & 157.3 & -- & 222.9 & 151.0 & -- \\
\midrule
Prototype Align & Acc & \textbf{72.08} & 56.67 & 88.11 & 78.11 & \textbf{55.81} & \textbf{81.67} & \textbf{72.08} \\
 & Token & \textbf{8259} & \textbf{10266} & \textbf{1594} & \textbf{1336} & 3009 & 922 & \textbf{4231.00} \\
 & Time & \textbf{255.6} & \textbf{335.5} & \textbf{24.1} & \textbf{21.5} & \textbf{59.4} & \textbf{13.2} & \textbf{118.22} \\
\midrule
%Kernel-ES Align & Acc & 69.58 & 60.00 & 86.43 & \textbf{79.43} & 54.42 & 79.33 & 71.53 \\
% & Token & 8489 & 10402 & 1607 & 1346 & \textbf{2995} & \textbf{912} & 4291.83 \\
% & Time & 283.4 & 350.6 & 28.1 & 25.5 & 67.2 & 14.3 & 128.18 \\
\bottomrule
\end{tabular}%
}%
}
\end{table*}
\end{comment}



\begin{table*}[t]
\centering
\setlength{\tabcolsep}{4pt}
\renewcommand{\arraystretch}{1.12}
\small
\resizebox{\textwidth}{!}{%
\begin{tabular}{@{}ll c ccc>{\columncolor{oursbg}}c cc>{\columncolor{oursbg}}c@{}}
\toprule
& & & \multicolumn{4}{c}{Fixed Steps} & \multicolumn{3}{c}{Entropy Stopping} \\
\cmidrule(lr){4-7}\cmidrule(lr){8-10}
Direction & Data & Text & \makecell{Linear\\Align} & \makecell{Exact-Soft\\Align} & \makecell{Filtered-Soft\\Align} & \makecell{\textbf{Prototype}\\\textbf{Align}} & \makecell{Exact-Soft\\Align} & \makecell{Filtered-Soft\\Align} & \makecell{\textbf{Prototype}\\\textbf{Align}} \\
\midrule
\multirow{4}{*}{\makecell{\textbf{Qwen3-8B}\\$\downarrow$\\\textbf{Qwen3-14B}}} & AIME 2025 & 60.00\tabsd{5.44} & 66.67\tabsd{5.44} & 61.11\tabsd{3.14} & 68.89\tabsd{8.31} & \textbf{70.00}\tabsd{2.72} & 58.89\tabsd{5.67} & 57.78\tabsd{5.67} & 56.67\tabsd{2.72} \\
 & MBPP+ & 78.84\tabsd{0.78} & 74.34\tabsd{0.94} & 78.31\tabsd{0.78} & 77.51\tabsd{0.75} & 76.19\tabsd{1.35} & 75.95\tabsd{0.22} & \textbf{82.01}\tabsd{0.65} & 72.31\tabsd{3.88} \\
 & MedQA & 78.11\tabsd{0.68} & 81.33\tabsd{0.82} & 81.11\tabsd{0.96} & 81.56\tabsd{0.57} & \textbf{82.00}\tabsd{0.98} & 78.67\tabsd{0.54} & 78.89\tabsd{1.10} & 80.89\tabsd{0.63} \\
\cmidrule(lr){2-10}
 & \textbf{AVG} & 72.32 & 74.11 & 73.51 & 75.99 & \textbf{76.06} & 71.17 & 72.89 & 69.96 \\
\midrule
\multirow{4}{*}{\makecell{\textbf{Qwen3-14B}\\$\downarrow$\\\textbf{Qwen3-8B}}} & AIME 2025 & 58.89\tabsd{1.57} & \textbf{63.33}\tabsd{5.44} & 58.89\tabsd{3.14} & 60.00\tabsd{4.71} & 61.11\tabsd{4.16} & 53.33\tabsd{2.72} & 57.78\tabsd{3.14} & 40.00\tabsd{4.71} \\
 & MBPP+ & 82.01\tabsd{0.57} & 21.34\tabsd{0.76} & 76.98\tabsd{0.57} & 76.54\tabsd{0.33} & 72.49\tabsd{1.42} & 71.78\tabsd{0.12} & \textbf{80.51}\tabsd{0.87} & 67.90\tabsd{0.76} \\
 & MedQA & 80.67\tabsd{0.27} & 77.00\tabsd{1.19} & 77.00\tabsd{0.72} & 78.33\tabsd{0.47} & 78.11\tabsd{1.29} & 75.89\tabsd{0.31} & \textbf{80.89}\tabsd{0.63} & 69.67\tabsd{0.72} \\
\cmidrule(lr){2-10}
 & \textbf{AVG} & 73.86 & 53.89 & 70.96 & 71.63 & 70.57 & 67.00 & \textbf{73.06} & 59.19 \\
\bottomrule
\end{tabular}}
\caption{Accuracy (\%, $\uparrow$) of cross-model communication between Qwen3-8B and Qwen3-14B (sender\,$\rightarrow$\,receiver; mean\,$\pm$\,population std over three runs). Code generation uses pass@1. \textbf{Bold}: best among latent alignment configurations within each direction and benchmark. AVG denotes the unweighted mean over the three benchmarks.}
\label{tab:cross_model_selected_results_soft_pending}
\end{table*}

\begin{comment}
\begin{table*}[t]
\centering
\caption{Performance and efficiency of cross-model communication between Qwen3-8B and Qwen3-14B (sender\,$\rightarrow$\,receiver; mean\,$\pm$\,std over three runs). Metrics are defined as in Table~\ref{tab:main_results}. \textbf{Bold}: best among latent alignment methods. AVG denotes the mean over the three benchmarks.}
\label{tab:cross_model_selected_results_soft_pending}
\setlength{\tabcolsep}{3pt}\renewcommand{\arraystretch}{1.0}\small
\resizebox{\textwidth}{!}{%
\begin{tabular}{@{}ll ccc>{\columncolor{oursbg}}c ccc>{\columncolor{oursbg}[\tabcolsep][0pt]}c@{}}
\toprule
& & \multicolumn{4}{c}{Qwen3-8B $\rightarrow$ Qwen3-14B} & \multicolumn{4}{c}{Qwen3-14B $\rightarrow$ Qwen3-8B} \\
\cmidrule(lr){3-6}\cmidrule(l){7-10}
Data & Metric & Text & \makecell{Linear\\Align} & \makecell{Exact-Soft\\Align} & \raisebox{-0.3\height}{\shortstack{\textbf{Kernel}\\\textbf{Align}}} & Text & \makecell{Linear\\Align} & \makecell{Exact-Soft\\Align} & \raisebox{-0.3\height}{\shortstack{\textbf{Kernel}\\\textbf{Align}}} \\
\midrule
\multirow{3}{*}{\makecell[l]{AIME\\2025}} & Acc & 60.00\tabsd{5.44} & 66.67\tabsd{5.44} & 58.89\tabsd{5.67} & \textbf{70.00}\tabsd{2.72} & 58.89\tabsd{1.57} & \textbf{63.33}\tabsd{5.44} & 53.33\tabsd{2.72} & 61.11\tabsd{4.16} \\
 & Token & 12{,}116\tabsd{369} & \textbf{11{,}622}\tabsd{438} & 15{,}381\tabsd{211} & 11{,}786\tabsd{127} & 10{,}716\tabsd{101} & 12{,}065\tabsd{466} & 14{,}891\tabsd{103} & \textbf{11{,}802}\tabsd{517} \\
 & Time & 382.6\tabsd{25.4} & 364.2\tabsd{6.4} & 530.3\tabsd{20.6} & \textbf{354.7}\tabsd{8.0} & 346.1\tabsd{17.9} & 289.9\tabsd{6.6} & 475.1\tabsd{9.6} & \textbf{283.4}\tabsd{17.9} \\
\addlinespace[2pt]\cmidrule(lr){1-10}\addlinespace[1pt]
\multirow{3}{*}{\makecell[l]{MBPP+}} & Acc & 78.84\tabsd{0.78} & 74.34\tabsd{0.94} & 75.95\tabsd{0.22} & \textbf{76.19}\tabsd{1.35} & 82.01\tabsd{0.57} & 21.34\tabsd{0.76} & 71.78\tabsd{0.12} & \textbf{72.49}\tabsd{1.42} \\
 & Token & 1{,}976\tabsd{1} & \textbf{1{,}552}\tabsd{22} & 6{,}989\tabsd{5} & 1{,}639\tabsd{60} & 2{,}338\tabsd{24} & 1{,}761\tabsd{29} & 8{,}447\tabsd{5} & \textbf{1{,}727}\tabsd{41} \\
 & Time & 30.3\tabsd{1.1} & \textbf{20.8}\tabsd{1.0} & 72.3\tabsd{0.6} & 21.0\tabsd{0.6} & 31.8\tabsd{0.4} & 18.1\tabsd{0.4} & 97.9\tabsd{0.7} & \textbf{17.6}\tabsd{0.3} \\
\addlinespace[2pt]\cmidrule(lr){1-10}\addlinespace[1pt]
\multirow{3}{*}{\makecell[l]{MedQA}} & Acc & 78.11\tabsd{0.68} & 81.33\tabsd{0.82} & 78.67\tabsd{0.54} & \textbf{82.00}\tabsd{0.98} & 80.67\tabsd{0.27} & 77.00\tabsd{1.19} & 75.89\tabsd{0.31} & \textbf{78.11}\tabsd{1.29} \\
 & Token & 1{,}662\tabsd{10} & 1{,}014\tabsd{13} & 4{,}497\tabsd{3} & \textbf{960}\tabsd{17} & 1{,}549\tabsd{34} & 1{,}403\tabsd{15} & 6{,}826\tabsd{21} & \textbf{1{,}395}\tabsd{25} \\
 & Time & 15.4\tabsd{0.3} & 13.0\tabsd{0.3} & 33.4\tabsd{0.4} & \textbf{11.5}\tabsd{0.7} & 14.9\tabsd{0.5} & \textbf{13.0}\tabsd{0.3} & 68.5\tabsd{0.3} & 13.5\tabsd{0.1} \\
\midrule
\multirow{3}{*}{\textbf{AVG}} & Acc & 72.32 & 74.11 & 71.17 & \textbf{76.06} & 73.86 & 53.89 & 67.00 & \textbf{70.57} \\
 & Token & 5{,}251 & \textbf{4{,}729} & 8{,}956 & 4{,}795 & 4{,}868 & 5{,}076 & 10{,}055 & \textbf{4{,}975} \\
 & Time & 142.8 & 132.7 & 212.0 & \textbf{129.1} & 130.9 & 107.0 & 213.8 & \textbf{104.8} \\
\bottomrule
\end{tabular}}
\end{table*}
\end{comment}


%We evaluate Prototype Align in both homogeneous and heterogeneous multi-agent settings. Table~\ref{tab:main_results} and Table~\ref{tab:cross_model_selected_results_soft_pending} report results on various benchmarks under the sequential configuration.


\noindent\textbf{Homogeneous collaboration.}
Table~\ref{tab:main_results_acc} reports that Prototype Align improves average accuracy under both fixed-budget and entropy-based stopping configurations. On Qwen3-8B, it achieves 67.51\% and 67.69\%, outperforming the strongest baseline Linear Align and Exact-Soft Align, by 3.19 and 3.15 percentage points, respectively. Similar results are also observed on Qwen3-14B, where Prototype Align achieves 73.13\% with fixed budgets and 71.49\% with entropy-based stopping, exceeding Exact-Soft Align by 1.89 and 0.62 percentage points, respectively. These results support the effectiveness of Prototype Align across both model scales and stopping configurations. These gains also extend to the hierarchical collaboration configuration, with detailed results provided in Appendix~\ref{app:hierarchical}.

Figure~\ref{fig:time_8b} compares end-to-end inference time on Qwen3-8B. Under entropy-based stopping, Prototype Align substantially reduces average runtime to 95.8 seconds per problem, compared with 297.3 seconds for Exact-Soft Align and 222.4 seconds for Filtered-Soft Align.
%, corresponding to 3.10$\times$ and 2.32$\times$ speedups, respectively.
Under fixed budgets, Prototype Align takes slightly longer, averaging 101.4 seconds compared with 90.9 and 90.5 seconds, but achieves higher average accuracy by 7.47 and 5.91 percentage points, respectively. These results demonstrate a favorable accuracy and latency trade-off for Prototype Align under both stopping configurations. Similar conclusions regarding accuracy and inference time hold for Qwen3-14B (Appendix \ref{app:sequential_14b_time}) and hierarchical collaboration, with detailed timing results provided in Appendix~\ref{app:results}.

\begin{comment}
Table~\ref{tab:main_results_acc} shows the performance of Prototype Align on homogeneous multi-agent collaboration. Compared with Exact-Soft Align, it achieves 2.93$\times$ and 2.52$\times$ speedups based on average end-to-end inference time in Qwen3-8B and Qwen3-14B, respectively. In addition, its average accuracy reaches 67.51\% and 73.13\%, improving by 2.97 and 2.26 percentage points. These results demonstrate that the Prototype Alignment substantially improves inference efficiency while maintaining strong task performance. It surpasses Exact-Soft Align on most benchmarks and remains competitive on the remaining ones. Compared with the other competitive alignment baseline Linear Align, our method gains 3.19 and 2.31 points, with comparable average inference time. We also observe that all latent alignment methods achieve higher average accuracy compared with conventional text based collaboration. Prototype Alignment yields the largest gains of 5.01 and 4.96 percentage points on Qwen3-8B and Qwen3-14B, respectively, alongside estimated speedups of 5.02$\times$ and 4.21$\times$. These benefits also extend to the hierarchical configuration.
Compared with Exact-Soft Align, Prototype Align improves average accuracy by 3.14 and 4.98 percentage points while achieving 2.87$\times$ and 2.98$\times$ speedups on Qwen3-8B and Qwen3-14B, respectively. Due to space constraints, the hierarchical results and supplementary evaluations on GSM8K, ARC-Easy, and ARC-Challenge are provided in Appendices~\ref{app:hierarchical} and~\ref{app:supplement}, respectively.
\end{comment}

\begin{figure}[t]
\centering
\includegraphics[width=0.48\textwidth]{figs/table1_time_8b.pdf}
\caption{End-to-end inference time (seconds per problem; logarithmic scale) of Qwen3-8B under the sequential configuration.}
\label{fig:time_8b}
\end{figure}

\begin{table}[t]
\centering
\small
\setlength{\tabcolsep}{6pt}
\renewcommand{\arraystretch}{1.12}
\begin{tabular}{lc}
\toprule
Method & Time (ms/step)\,$\downarrow$ \\
\midrule
Text & 1.156\tabsd{0.004} \\
Linear Align & 0.085\tabsd{0.010} \\
Filtered-Soft Align & 1.074\tabsd{0.023} \\
Exact-Soft Align & 1.541\tabsd{0.024} \\
Prototype Align & 0.255\tabsd{0.010} \\
\bottomrule
\end{tabular}
\caption{Alignment time per step for Qwen3-8B
(sequential; mean\,$\pm$\,std over 12 runs).}
\label{tab:qwen3_8b_sequential_alignment_avg}
\end{table}

\noindent\textbf{Heterogeneous collaboration.}
We further evaluate communication between agents with different hidden dimensions but compatible vocabularies under both fixed-budget and entropy-based stopping configurations. Table~\ref{tab:cross_model_selected_results_soft_pending} reports accuracy for both communication directions. With fixed budgets, Prototype Align achieves average accuracies of 76.06\% and 70.57\% for Qwen3-8B$\rightarrow$14B and Qwen3-14B$\rightarrow$8B, respectively, outperforming Linear Align by 1.95 and 16.68 percentage points. 
%{\color{blue}
Under entropy-based stopping, it achieves average accuracies of 69.96\% and 59.19\%, which are 2.93 and 13.87 percentage points below the best value, respectively. 
%These results indicate that the effectiveness of Prototype Align in heterogeneous latent communication depends on the stopping strategy.}
%{\color{red}It is likely because the gap between the intellegence of the larger model and that of smaller model is not well described by our method.}
When transferring from a larger sender to a smaller receiver, Prototype Align may not sufficiently adapt the transferred representations to the smaller model's processing capabilities.


Prototype Align also exhibits a directional difference, where communication from Qwen3-8B to Qwen3-14B achieves higher average accuracy than the reverse direction under fixed budgets (76.06\% vs.\ 70.57\%) and entropy stopping (69.96\% vs.\ 59.19\%). This comparison suggests that transferring information from a larger sender does not necessarily compensate for the limited capabilities of a smaller receiver. In addition, we observe latent alignment still underperforms text-based collaboration in some heterogeneous settings, including MBPP+ in both directions and MedQA in the 14B$\rightarrow$8B direction. These results highlight remaining challenges in cross-model latent communication, suggesting that alignment alone may not fully bridge differences in models' internal representations.
%}
%\noindent\textbf{Heterogeneous collaboration.} We further evaluate whether the alignment interface extends to agents with different hidden dimensions but compatible vocabularies. Table \ref{tab:cross_model_selected_results_soft_pending} shows results. Across both communication directions, Prototype Align achieves higher mean accuracy than Exact-Soft Align on all evaluated benchmarks, while providing end-to-end inference speedups of 1.64$\times$ and 2.04$\times$. Prototype Align also outperforms Linear Align across all evaluated benchmarks at both scales, except for AIME 2025 on Qwen3-14B$\rightarrow$8B. These results support the applicability of Prototype Align to heterogeneous collaboration, enabling efficient communication between agents with different hidden dimensions without additional alignment training. Nevertheless, we notice that latent alignment still underperforms text based collaboration in some heterogeneous settings, including MBPP+ in both directions and MedQA in the 14B$\rightarrow$8B direction. These results indicate the remaining challenges in cross-model latent communication, suggesting that efficient alignment alone may not fully bridge differences in models' internal representations.

\noindent\textbf{Per Step Inference Efficiency}
We examine the computational cost of the alignment operation itself. Table~\ref{tab:qwen3_8b_sequential_alignment_avg} reports the average alignment time per step for Qwen3-8B under the sequential configuration. Further measurement details are provided in Appendix \ref{app:per_step_efficiency}. Prototype Align reduces the average alignment time from 1.54 ms for Exact-Soft Align to 0.26 ms, achieving an approximately $6\times$ speedup. Although Linear Align is faster at 0.09 ms per step, both methods require less than 0.3 ms per alignment step.

%We evaluate kernelized alignment on six benchmarks, using Qwen3-8B and Qwen3-14B under the sequential multi-agent configuration. Table~\ref{tab:main_results} reports task performance and computational efficiency.


%\textbf{Accuracy.} Kernelized alignment achieves the highest average accuracy among the compared methods at both scales, reaching 66.17\% on Qwen3-8B and 72.08\% on Qwen3-14B. These results exceed linear alignment by 1.85 and 1.26 percentage points, respectively, and text-based collaboration by 3.67 and 3.91 points. These results suggest that the Prototype Alignment supports more effective latent reasoning and collaboration across the evaluated model scales.

%\textbf{Inference efficiency.} To account for fair comparison in inference time, we normalize the reported \textit{Text} and \textit{Linear Align} timings. We use a task-specific time-per-token rate, $r=\mathrm{Time}/\mathrm{Token}$, calibrated on the hardware used for Prototype Align. The adjusted times are estimated as $r$ multiplied by the reported token counts. The additional latent-computation and alignment costs of baselines are ignored given its relatively smaller compared to decoding time. Results indicate that Kernelized alignment substantially reduces inference cost relative to text-based collaboration, achieving 4.40$\times$ and 4.09$\times$ speedups on Qwen3-8B and Qwen3-14B, respectively, based on the hardware-adjusted timings in Table~\ref{tab:main_results}. Among latent alignment methods, it avoids the substantial overhead of exact-soft alignment while maintaining inference times comparable to or lower than linear alignment on average. Together, these results suggest that Kernelized alignment provides a more effective interface for latent reasoning and collaboration.

%\xiyan{\paragraph{Cross-model latent communication.}
%Table~\ref{tab:cross_model_results} further evaluates the communication interface when the sender and receiver use different model checkpoints. We use a two-agent Planner--Judger setup with Qwen3-8B and Qwen3-14B, which have compatible vocabularies but different hidden dimensions. The sender's prefill hidden states are aligned and provided to the receiver as additional input. Prototype Align achieves average accuracies of XX\% and XX\% in the 14B$\rightarrow$8B and 8B$\rightarrow$14B directions, respectively, compared with XX\% and XX\% for Linear Align. Relative to Exact-Soft Align, the corresponding accuracy differences are XX and XX percentage points, while end-to-end inference time is reduced by XX\% and XX\%, respectively.}

%\paragraph{Cross-model latent communication.} Agents in a multi-agent system may use models of different scales, making alignment between different representation spaces necessary for latent communication. We therefore investigate whether Prototype Align remains effective when the sender and receiver have different hidden dimensions. To limit additional effects from multi-agent interactions, we restrict the evaluation to a two agent Planner Judger configuration following \citep{fu2026cachetocache}, using Qwen3-8B as the Planner and Qwen3-14B as the Judger. The models have compatible vocabularies, and the Planner's prefill hidden states are aligned and passed to the Judger as additional input. Table~\ref{tab:cross_model_selected_results_soft_pending} shows that Prototype Align outperforms Linear Align on all three benchmarks, improving accuracy by 3.33, 1.85, and 0.67 percentage points on AIME 2025, MBPP+, and MedQA, respectively, with comparable or lower inference times. These results suggest that Prototype Align can support effective latent communication across model scales without additional training.


\subsection{Analysis}
% Preamble
\begin{comment}
\begin{wrapfigure}[15]{r}{0.40\columnwidth}
    \vspace{-8pt}
    \centering
    \includegraphics[width=\linewidth]{figs/tsne-rep.png}
    \setlength{\abovecaptionskip}{3pt}
    \setlength{\belowcaptionskip}{0pt}
    \caption{Joint t-SNE visualization of kernelized and target alignment outputs.}
    \label{fig:alignment_quality}
    \vspace{-8pt}
\end{wrapfigure}

\paragraph{Alignment quality.} We examine how well the estimated alignment preserves the output representations of the target mapping defined in Eq.~\ref{eq:kernel_substituted}. We apply both mappings to the same set of 100 hidden states and jointly project their outputs into two dimensions using t-SNE. As shown in Figure~\ref{fig:alignment_quality}, the kernelized alignment representations and reference representations overlap in part of the visualization. This observation suggests partial correspondence between the two sets of representations in the projected space.
\end{comment}

\begin{figure}[t]
\centering
\includegraphics[width=0.6\columnwidth]{figs/medqa_pca.png}
\caption{Joint PCA projection of 400 Prototype Align outputs and 400 Exact-Soft Align outputs on MedQA with Qwen3-8B.}
\label{fig:alignment_outputs}
\end{figure}

\begin{figure}[t]
\centering
\includegraphics[width=0.9\columnwidth]{figs/medqa_entropy.png}
\caption{Next-token entropy across latent states for Linear Align, Prototype Align, and Exact-Soft Align on MedQA with Qwen3-8B.}
\label{fig:reasoning_dynamics}
\end{figure}

\begin{comment}
\begin{figure*}[t]
\centering
\begin{minipage}[c]{0.36\textwidth}
    \centering
    \includegraphics[width=0.7\linewidth]{figs/medqa_pca.png}

    \small (a) Alignment outputs
\end{minipage}\hfill
\begin{minipage}[c]{0.62\textwidth}
    \centering
    \includegraphics[width=0.7\linewidth]{figs/medqa_entropy.png}

    \small (b) Next-token entropy
\end{minipage}
\caption{Alignment outputs and latent-state dynamics on MedQA
with Qwen3-8B. (a) Joint PCA projection of 400 Prototype Align
outputs and 400 Exact-Soft Align outputs. (b) Next-token
entropy across latent states for Linear Align, Prototype Align,
and Exact-Soft Align.}
\label{fig:alignment_dynamics}
\end{figure*}
\end{comment}

\paragraph{Alignment}
We first examine whether Prototype Align preserves the representations induced by the exact distribution mapping. To this end, we apply Prototype Align and Exact-Soft Align to the same 400 hidden states collected from reasoning trajectories on MedQA. We then jointly project their output representations using PCA \citep{pearson1901liii}. Figure~\ref{fig:alignment_outputs} illustrates results. Results of more datasets are provided in the Appendix \ref{app:alignment}. We observe that the two sets of representations exhibit substantial overlap. This finding suggests that Prototype Align retains broad characteristics of the exact mapping despite using a finite dimensional approximation.

\paragraph{Reasoning Dynamics} To better understand the performance gains of Prototype Align, we examine how its approximate mapping shapes the dynamics of latent reasoning.
%We further investigate whether Prototype Alignment leads to distinct dynamics as the aligned representations are propagated through latent reasoning steps. 
%Specifically, we tracks the next-token entropy across the latent rollout. 
Specifically, at each latent reasoning step, we project the last-layer hidden state through the model's output head and compute the entropy of the resulting next token distribution. This probe is used only for analysis and does not alter the latent trajectory. More details are provided in Appendix \ref{app:dynamics}. Figure~\ref{fig:reasoning_dynamics} illustrates results. We find that Prototype Align exhibits an initial increase followed by a gradual decline, whereas Linear Align maintains relatively high entropy and Exact-Soft Align remains at a lower level. 
%This pattern suggests that Prototype Align supports relatively diverse predictions during intermediate steps before transitioning toward more concentrated predictions. Such dynamics may facilitate exploration early in the reasoning process while allowing the model to consolidate its predictions at later stages.
The higher intermediate entropy indicates that Prototype Align assigns probability mass to a broader set of candidate tokens, while the subsequent decline reflects increasingly confident predictions. It is likely that the random-feature approximation perturbs the latent trajectory, assisting the model avoid committing too early to a narrow set of predictions. The finding is consistent with prior work showing that stochastic perturbations facilitate exploration in latent reasoning \citep{you-etal-2026-parallel}.

%\paragraph{Variance of aligned states.} We assess the sensitivity of Prototype Alignment to the ORF seed by repeatedly aligning hidden states sampled across a multi-agent trajectory. For each state, we measure the dispersion of seed-specific representations around their seed-averaged centroid. This analysis characterizes approximation stability, rather than cross-state geometry preservation (Figure~\ref{fig:orf_seed_robustness}).

\begin{figure}[t]
    \centering
    \includegraphics[width=0.9\linewidth]{figs/orf.png}
    \caption{Distribution of ORF-seed similarity $S_i=1-D_i$
    across 500 hidden states and 50 random seeds.
    Higher values indicate greater stability.}
    \label{fig:orf}
\end{figure}
\paragraph{Sensitivity to random feature initialization.}
To assess the sensitivity of Prototype Alignment to orthogonal random feature (ORF) initialization, we examine how much the aligned representation of a fixed hidden state varies across independently sampled random feature maps. Specifically, we draw a balanced sample of 500 hidden states from Qwen3-8B reasoning trajectories on MedQA and align each state using 50 different ORF seeds, while keeping the feature count, kernel temperature, and embedding and unembedding weights fixed. For each state, we compute the relative root-mean-square (RMS) dispersion $D_i$ of its seed-specific aligned representation around their mean and visualize the corresponding seed similarity $S_i=1-D_i$. Figure \ref{fig:orf} illustrates results. We observe that the median similarity is 85.07\%, and the 5th percentile is 84.56\%. The results suggest that Prototype Align produces reasonably consistent representations across ORF seeds and is relatively robust to random feature initialization. 

%\mingzhe{TODO: I am not sure how to introduce the benefit implied by the result. Maybe we can say that this indicates the consistency across kernel random seeds, but also brings randomness to enhance the capability of exploration.} More details are provided in Appendix~\ref{app:orf}.
%\xiyan{randomness between seeds seem unrelated to exploration}

%\paragraph{Reasoning dynamics.}
%Different alignment methods may shape the reasoning process as latent representations are propagated across intra-model reasoning and inter-model communication. We therefore examine how alignment methods affect latent reasoning. Specifically, we track the entropy of the next-token distribution induced by intermediate hidden states. This provides a diagnostic of how concentrated the model's predictions become as latent computation proceeds. At each step $t$, we compute $p_t = \operatorname{softmax}(W_{\mathrm{out}}h_t+b)$ and $\mathcal{H}_t = -\sum_{v \in \mathcal{V}}p_t(v)\log p_t(v)$. These probes guarantee that the reasoning trajectories are not affected. Figure~\ref{fig:entropy_dynamics} illustrates various alignment under the same latent-step budget. Additional tasks and budgets are reported in the appendix. 

%We observe distinct entropy trajectories across alignment methods. Linear alignment maintains relatively high entropy over extended latent rollouts. Prototype Alignment instead exhibits a transient increase followed by a decline toward a low-entropy regime, while exact-soft alignment remains at relatively low entropy over much of the trajectory. Agent transitions introduce additional transient changes, after which Prototype Alignment often returns to low entropy. The persistence of low entropy under Prototype Alignment motivates us to investigate whether further latent steps can be omitted without sacrificing task accuracy.




\begin{comment}
    

\begin{figure*}[t]
\centering

% Left: entropy figure
\begin{minipage}[t]{0.48\textwidth}
    \vspace{0pt}
    \centering
    \includegraphics[width=\linewidth]{figs/entropy.png}
    \caption{\xiyan{TODO a new graph. Entropy trajectories under different alignment methods on XX using Qwen3-8B. Vertical lines mark agent transitions.}}
    \label{fig:entropy_dynamics}
\end{minipage}
\hfill
% Right: stopping-strategy table
\begin{minipage}[t]{0.48\textwidth}
    \vspace{0pt}
    \centering
    \label{tab:entropy_stopping}

    \setlength{\tabcolsep}{3pt}
    \renewcommand{\arraystretch}{1.1}
    \small
    \resizebox{\linewidth}{!}{%
    \xiyan{%
    \begin{tabular}{@{}llccc@{}}
        \toprule
        Strategy & Metric & GSM8K & AIME 2025 & MBPP+ \\
        \midrule
        Fixed budget
        & Acc $\uparrow$   & -- & -- & -- \\
        & Steps $\downarrow$ & -- & -- & -- \\
        & Time $\downarrow$  & -- & -- & -- \\
        \midrule
        Random stopping
        & Acc $\uparrow$   & -- & -- & -- \\
        & Steps $\downarrow$ & -- & -- & -- \\
        & Time $\downarrow$  & -- & -- & -- \\
        \midrule
        Entropy stopping
        & Acc $\uparrow$   & -- & -- & -- \\
        & Steps $\downarrow$ & -- & -- & -- \\
        & Time $\downarrow$  & -- & -- & -- \\
        \bottomrule
    \end{tabular}%
    }%
    }
    \captionof{table}{\xiyan{Stopping strategies for Prototype Align
    using Qwen3-8B. Acc is accuracy or pass@1 (\%); Steps is
    the average total latent-step count per problem; Time is
    seconds per problem, including entropy checks.}}
\end{minipage}

\end{figure*}
\end{comment}



%\textbf{Scaling with latent steps.} We vary the number of latent reasoning steps to characterize the performance--computation trade-off. This analysis identifies whether additional latent computation continues to improve accuracy or reaches a regime of diminishing returns (Figure~\ref{fig:latent_step_scaling}).


%\subsection{Ablation Study}
\begin{comment}
\paragraph{Latent Step Ablation.}
The latent step budget is an important hyperparameter when applying kernelized alignment to downstream multi-agent tasks. We therefore compare two budget-control strategies: i) \textit{fixed budget reasoning}, which uses the best fixed budget selected on a development set; and ii) \textit{entropy based adaptive stopping}~\citep{soft_thinking}, which checks entropy every 10 latent steps and terminates the current agent's latent reasoning steps after two consecutive checks below 0.25. We also include a \textit{random stopping} baseline, where stopping positions are sampled independently ranging from \{0, 10, 20, 40, 80\} of the latent steps following \citep{latentmas}. All strategies allow at most 200 latent steps per intermediate agent. Table~\ref{tab:entropy_stopping} shows results. Implementation details are provided in Appendix~\ref{app:latent_step}. 

The relative performance of the fixed and entropy strategies varies across benchmarks, but their average accuracies differ by only 0.18 points. Entropy-based stopping reduces average inference time from 101.40 to 95.78 seconds per problem. Both strategies outperform random stopping, improving average
accuracy by 3.00 and 3.18 percentage points, respectively. These results suggest that both strategies are viable for Prototype Align.
%We retain fixed budgets as the default setting given their stability across runs. 
For a fair comparison with other baseline, e.g., linear align, we chose to follow the same fixed budget stopping criteria as theirs.
Beyond fixed stopping, entropy-based stopping on development examples may also help identify candidate fixed budgets and narrow the search range.
\end{comment}

\begin{comment}
    

\begin{table*}[t]
\centering
\caption{Comparison of latent reasoning step control strategies for Prototype Align with Qwen3-8B under the sequential configuration (mean\,$\pm$\,std over three runs). \textbf{Bold}: best per column.}
\label{tab:entropy_stopping}
\setlength{\tabcolsep}{3pt}\renewcommand{\arraystretch}{1.0}\small
\resizebox{\textwidth}{!}{%
\begin{tabular}{@{}ll cccccc c@{}}
\toprule
Strategy & Metric & \makecell{AIME\\2024} & \makecell{AIME\\2025} & \makecell{Human\\Eval+} & MBPP+ & \makecell{GPQA-\\Diamond} & MedQA & AVG \\
\midrule
 & Acc & 62.22\tabsd{3.14} & 46.67\tabsd{7.20} & 83.54\tabsd{0.50} & 74.60\tabsd{0.99} & 47.47\tabsd{0.41} & 72.56\tabsd{1.10} & 64.51 \\
 & Token & 9{,}031\tabsd{466} & 9{,}974\tabsd{175} & 1{,}648\tabsd{30} & 1{,}478\tabsd{19} & \textbf{3{,}534}\tabsd{81} & \textbf{1{,}120}\tabsd{32} & 4{,}464 \\
\multirow{-3}{*}{\makecell[l]{Random stopping}} & Time & 259.6\tabsd{18.6} & 300.9\tabsd{14.7} & \textbf{17.8}\tabsd{0.1} & 18.4\tabsd{0.7} & 59.9\tabsd{3.0} & \textbf{12.7}\tabsd{1.0} & 111.5 \\
\addlinespace[2pt]\cmidrule(lr){1-9}\addlinespace[1pt]
 & Acc & \textbf{68.89}\tabsd{1.57} & 54.44\tabsd{4.16} & \textbf{83.94}\tabsd{0.57} & 75.40\tabsd{0.22} & 49.16\tabsd{1.04} & \textbf{73.22}\tabsd{0.87} & 67.51 \\
 & Token & \textbf{8{,}548}\tabsd{344} & 10{,}230\tabsd{286} & \textbf{1{,}639}\tabsd{42} & \textbf{1{,}463}\tabsd{19} & 3{,}574\tabsd{87} & 1{,}163\tabsd{20} & 4{,}436 \\
\multirow{-3}{*}{\makecell[l]{Fixed budget\\(default)}} & Time & 229.3\tabsd{3.3} & 271.2\tabsd{5.8} & 18.9\tabsd{0.2} & \textbf{18.1}\tabsd{0.2} & \textbf{57.6}\tabsd{0.6} & 13.3\tabsd{0.5} & 101.4 \\
\addlinespace[2pt]\cmidrule(lr){1-9}\addlinespace[1pt]
 & Acc & 65.56\tabsd{1.57} & \textbf{57.78}\tabsd{4.16} & \textbf{83.94}\tabsd{1.60} & \textbf{75.49}\tabsd{0.50} & \textbf{50.17}\tabsd{1.04} & \textbf{73.22}\tabsd{1.29} & \textbf{67.69} \\
 & Token & 8{,}552\tabsd{586} & \textbf{9{,}419}\tabsd{427} & 1{,}692\tabsd{34} & 1{,}532\tabsd{11} & 3{,}865\tabsd{77} & 1{,}148\tabsd{8} & \textbf{4{,}368} \\
\multirow{-3}{*}{\makecell[l]{Entropy stopping}} & Time & \textbf{216.0}\tabsd{13.2} & \textbf{244.2}\tabsd{18.6} & 21.5\tabsd{0.4} & 19.6\tabsd{0.4} & 60.0\tabsd{0.6} & 13.4\tabsd{0.1} & \textbf{95.8} \\
\bottomrule
\end{tabular}}
\end{table*}
\end{comment}

\begin{comment}
\begin{wraptable}{r}{0.48\textwidth}
    \centering
    \caption{Effect of random feature dimension $m$.
    Tokens and time are reported per problem.}
    \label{tab:feature_dimension}
    \small
    \setlength{\tabcolsep}{5pt}
    \renewcommand{\arraystretch}{1.1}
    \begin{tabular}{@{}rrrr@{}}
        \toprule
        $m$ & Acc. (\%) & Tokens & Time (s) \\
        \midrule
        %128  & 71.33 & 1131.9 & 13.45 \\
        %256  & 71.67 & 1150.7 & 12.48 \\
        %512  & 71.67 & 1120.9 & 13.31 \\
        %1024 & 75.00 & 1148.0 & 13.32 \\
        %2048 & 75.33 & 1132.0 & 13.21 \\
        128 &71.33 &1096.9 & 13.26\\
        256 & 73.33 &1089.7 & 13.46\\
        512 & 74.00 &1149.6 & 13.84\\
        1024 & 76.67 &1171.1 & 13.17\\
        2048 & 76.67 &1153.5 & 13.97\\
        \bottomrule
    \end{tabular}
\end{wraptable}
\end{comment}

\begin{comment}
\begin{figure}[t]
    \centering
    \includegraphics[width=\linewidth]{figs/feature_dim.pdf}
    \caption{Effect of random feature dimension $m$.}
    \label{fig:feature_dim}
\end{figure}
\end{comment}

\begin{table}[t]
\centering
\small
\setlength{\tabcolsep}{6pt}
\renewcommand{\arraystretch}{1.1}
\begin{tabular}{@{}rcc@{}}
\toprule
$m$ & Accuracy (\%) $\uparrow$ & Time (s) $\downarrow$ \\
\midrule
128  & 71.33 $\pm$ 2.49 & 13.17 $\pm$ 1.67 \\
256  & 73.33 $\pm$ 0.94 & 13.26 $\pm$ 1.96 \\
512  & 74.00 $\pm$ 1.63 & 13.84 $\pm$ 1.16 \\
1024 & 76.67 $\pm$ 2.49 & 13.97 $\pm$ 2.22 \\
2048 & 76.67 $\pm$ 3.40 & 13.46 $\pm$ 0.68 \\
\bottomrule
\end{tabular}
\caption{Effect of feature dimension $m$ on 50 MedQA development examples with Qwen3-8B. Results are mean $\pm$ population standard deviation over three runs. Time denotes end-to-end seconds per problem.}
\label{tab:feature_dimension}
\end{table}

\paragraph{Orthogonal Random Features Ablation}
The number of orthogonal random features determines the approximation capacity and computational cost of Prototype Align. We therefore investigate how the feature dimension affects downstream task performance and inference efficiency. Specifically, we evaluate five feature dimensions $m$ ranging from 128 to 2048 with other factors fixed on the MedQA development set of Qwen3-8B model. For each setting, we construct the corresponding orthogonal random feature map and precompute the vocabulary-dependent statistics. 
%We keep the alignment temperature, latent-step budgets, model configurations, and generation settings fixed across feature dimensions to isolate the effect of the feature budget. 
%Table~\ref{tab:feature_dimension} shows 
%Fig~\ref{fig:feature_dim} illustrates an overall upward trend in accuracy and tokens as the feature dimension increases. 
Table~\ref{tab:feature_dimension} reports mean accuracy and end-to-end inference time. Mean accuracy increases from 71.33\% at $m=128$ to 76.67\% at $m=1024$, and remains unchanged at $m=2048$. Mean inference time ranges from 13.17 to 13.97 seconds per problem, with no monotonic increase across feature dimensions. We adopt $m=1024$ as the default because it achieves the highest observed mean accuracy with half as many features as $m=2048$.
%Accuracy improves substantially from 71.33\% at $m=128$ to 76.67\% at $m=1024$. 
%Meanwhile, we observe that tokens and end-to-end inference time remain similar.
%We adopt $m=1024$ as the default because it achieves the highest observed accuracy with half as many random features as $m=2048$, while further increasing the feature dimension provides no additional accuracy gain.


% Required packages: booktabs, multirow, makecell, graphicx
\providecommand{\smallerr}[1]{}
\renewcommand{\smallerr}[1]{{\scriptsize$\pm#1$}}


%\paragraph{Noise injection.} We evaluate linear and Prototype Alignment under matched clean and noisy latent trajectories in a Planner--Judger system. Source hidden states are perturbed with norm-scaled Gaussian noise before alignment. We report paired changes in task accuracy and the difference in degradation between methods.



%\clearpage
\begin{comment}

\section{Introduction}
Large language model (LLM) systems are playing an increasingly important role in handling complex tasks, including mathematics, coding, and scientific problems \cite{agentic_ai}. An LLM system consists of one or multiple LLM agent(s), which generate outputs autoregressively over multiple steps. Consecutive steps executed by the same agent constitute intra-model reasoning, whereas steps executed by two agents constitute inter-model communication. Due to the statelessness of LLMs, the system requires an intermediate representation to carry information from the past into the next step, which we call the "medium”. Different mediums have different expressivity and efficiency. As task complexity continues to grow and model capabilities continue to improve, the context required to pass over steps is getting longer, which further emphasizes the importance of the chosen medium.


A good medium should satisfy two criteria: (1) \emph{expressivity}: it should provide a high information-density representation that preserves key information; and (2) \emph{efficiency}: it should reduce computational cost as much as possible for LLM architectures. Text, as a widely available and readily interpretable medium, is the medium for the majority of LLM systems. Yet, a discrete token has to be sampled from a distribution over the vocabulary, i.e., logits, which wastes abundant information in other possible tokens \cite{soft_thinking}; also, a time-consuming vocabulary-sized matrix multiplication is required for decoding, which is to read out logits from hidden states (can account for more than 80\% of inference time \cite{decoding_time}). Prior works \cite{soft_thinking} have explored using logits as the medium to retain the distribution information, but the efficiency issue remained. Coconut first proposes to use hidden states as the medium, which is both information-rich and decoding-free. However, it also introduces a distribution shift problem. That is, the hidden space (the space immediately before unembedding) and the embedding space (the space immediately after embedding) do not completely match. Consequently, either the whole system needs further training, or each passed hidden state has to be aligned as a synonymous embedding state with an adapter.


Training the whole LLM system or neural network (NN) adapters \cite{coconut,c2c,interlat,latent_cache_flow,cache_alignment} is computationally expensive, especially for a large or frequently-changing system, and the alignment is only functionally effective for hidden states following common distributions, ignoring some vital outliers. Some training-free alignment methods \cite{latentmas,later}, which only rely on model parameters, overcome this drawback. However, existing training-free methods are not sufficiently accurate \cite{on_latentmas}, leading to inevitable instability in the reasoning process. Therefore, a training-free and accurate alignment method is required for hidden-state-mediated LLM systems.



In this work, we propose a mathematical framework to unify intra-model reasoning and inter-model communication. Also, noticing the isomorphism between the decoding layer and attention layer, we propose a training-free Kernelized Alignment Method to address the distribution shift problem under the unified framework in hidden-state-mediated LLM systems, inspired by \cite{svm,linear_attentions,performer}. We use vocabulary space as an intermediary between hidden space and embedding space, and merge hidden-vocabulary and vocabulary-embedding maps with orthogonal random features (ORFs) for approximation to cancel explicit decoding. Theoretically, the computational complexity is reduced from $O(|V|(d_1+d_2))$ to $O(m(d_1+d_2))$, where $|V|$ denotes the vocabulary size, $m \ll |V|$ denotes the number of features, and $d_1$, $d_2$ denote the embedding dimension of the source agent and the target agent respectively (the two agents can be the same agent).



Through end-to-end experiments on mathematics \cite{gsm8k,aime2024,aime2025}, code \cite{evalplus}, and science \cite{arc,gpqa,medqa} reasoning tasks using Qwen3-8B and Qwen3-14B \cite{qwen3} and DeepSeek-R1-Distill-Llama-8B \cite{deepseek_r1_distill_llama_8b}, we show that our method can attain higher accuracy, faster inference speed and fewer decoding times than prior approaches. Additionally, we find that our method is more stable than other training-free hidden-state-mediated methods and more exploratory than text- and logit-mediated methods. We further explore whether the entropy of probed logits can be used as the stopping indicator of reasoning. Finally, we discuss the possibility of extending our method to cross-vocabulary systems with optimization transport theory.


\section{Related Work}

%\subsection{Multi-step generation in LLM systems}

Modern LLMs are predominantly based on decoder-only transformer architectures \cite{gpt2} and therefore generate outputs autoregressively and statelessly. Through multi-step generation, information gradually accumulates in the LLM system. In each step, the source agent and the target agent may be the same or different. We define that intra-model reasoning denotes the case where the source and target agents are the same, while inter-model communication denotes the case where they are different. 

\textbf{Intra-model reasoning.} Early works \cite{wei2022cot,kojima2022zeroshot} on intra-model reasoning use prompts to elicit CoT reasoning, which extends the thinking process and also improves the quality of the final answer. Following works \cite{zelikman2022star,hsieh2023distilling,deepseek_r1_distill_llama_8b} improve CoT reasoning through supervised fine-tuning with CoT generated by stronger models or humans, and end-to-end reinforcement learning. To further extend the exploration capability of reasoning, \cite{yao2023tot,besta2024got} propose to use tree- or graph-structured reasoning to search for multiple paths in parallel; to improve the reliability and stability of reasoning, \cite{shinn2023reflexion,yao2023react} use self-reflection mechanism or tools to re-evaluate the reasoning process and correct errors of previous steps. Recently, it has been found that recursive reasoning in one step without explicit CoT can also achieve comparable performance to CoT reasoning \cite{coconut}, which can also be regarded as a special case of latent reasoning.

\textbf{Inter-model communication.} One single agent is insufficient to handle complex tasks, but a multi-agent system (MAS) that consists of different experts assist to bridge the gap. \cite{zhang2024chainofagents,hong2024metagpt} use sequential message flows to enable inter-model communication, and \cite{wang2024moa,liu2024dylan,hu2025owl} create more complex communication topologies, such as hierarchical, star, and more general graph structures. In order to overcome the hallucination and error propagation inherent to individual models, \cite{du2024debate,estornell2024multillm,liang2024divergent} introduce iterative communication rounds, where agents can discuss and refine their responses via adversarial or cooperative dynamics. Recent works such as \cite{anthropic2024mcp,google2025a2a} aim to establish a protocol for communication in MAS, which standardizes the interaction and tool usage of LLM agents.
\end{comment}

\begin{comment}
\subsection{Protocols for LLM multi-step generation}

Since most LLMs are stateless, a medium is required to carry information from the past into the next step. Here we introduce three commonly used mediums: text, logits, and hidden states, and discuss their advantages and disadvantages.

\textbf{Text} is the most common medium for its availability and interpretability. \cite{wei2022cot}, as the foundation of modern LLMs, first use text as the medium in intra-model reasoning, and \cite{hong2024metagpt,zhang2024chainofagents} first use text as the medium in inter-model communication. However, text is sampled from logits, thus does not fully utilize the distribution over the vocabulary \cite{soft_thinking}, and the decoding process is very time-consuming \cite{decoding_time}, which is rather inefficient for long-context tasks. 

\textbf{Logits}, as the medium, is explored by \cite{soft_thinking} to retain the distribution information in intra-model reasoning. \cite{shen2024collaborative} further extend logits-mediated method into inter-model communication. Also, logit-based medium preserves the intepretability and differentiability. However, the decoding process is still required.

\textbf{Hidden states}, which are both information-rich and decoding-free, offer a promising alternative \cite{coconut}. By fine-tuning a model \cite{coconut}, training adapters \cite{c2c,interlat,latent_cache_flow,cache_alignment}, or training-free adapters \cite{latentmas,later,rotation_alignment}, the distribution shift problem between hidden space and embedding space can be resolved. However, training process in a large or frequently-changing MAS is extremely expensive, and current training-free methods are not sufficiently accurate adapters. In this work, we propose a training-free adapter to achieve the efficiency--effectiveness balance.

\end{comment}

\begin{comment}
\section{Method}

We propose a kernelized alignment method that leverages the properties of random features (RFs) to resolve the distribution shifts between hidden space and embedding space, enabling both intra-model CoT and inter-model communication to completely proceed in hidden spaces, which is more efficient than text-based or logit-based approaches and also preserves the information needed for subsequent steps. Figure~\ref{fig:method_overview} contrasts the conventional decoding--re-embedding pipeline with its vocabulary-indexed attention interpretation and summarizes our offline--online kernelized computation. In Problem Formulation, we propose a unified mathematical framework for multi-step generation in LLM systems. In Mathematical Preliminaries, we introduce the linear attention mechanism and orthogonal random features (ORFs) that are used in our method. In Kernelized Alignment, we present our kernelized alignment method under our framework to resolve the distribution shift problem.

%\xiyan{We formalize intra-model reasoning and inter-model communication as a same information transmission problem, which align the hidden state  before the unembedding matrix to the embedding space in Section~\ref{sec:problem_formulation}. Building on this unified view, we propose a kernelized alignment method to solve this transmission problem in Section~\ref{sec:kernelized_alignment}. Figure~\ref{fig:method_overview} illustrates the method.}




\subsection{Problem Formulation}
Consider a multi-step LLM system $S$ comprising multiple agents. We introduce the notation for an LLM agent $A \in S$. Let $\mathcal{V}_A$ denote its vocabulary, and let $d_A$ denote its hidden dimension. Let $V_A$, $E_A$, and $H_A$ denote its one-hot token, embedding, and hidden-state spaces, respectively, while $\tilde{V}_A$ denotes the space of probability distributions over $\mathcal{V}_A$. Specifically, $V_A \cong \{\textbf{e}_i\}_{i=1}^{|\mathcal{V}_A|}$, $E_A \cong \mathbb{R}^{d_A}$, $H_A \cong \mathbb{R}^{d_A}$, and $\tilde{V}_A \cong \Delta(\mathcal{V}_A)$, where $\Delta(\mathcal{V}_A)$ is the probability simplex over $\mathcal{V}_A$. Let $WA^{in} \in \mathbb{R}^{d_A \times |\mathcal{V}_A|}$ and $W_A^{out} \in \mathbb{R}^{|\mathcal{V}_A| \times d_A}$ denote the embedding and decoding matrices, respectively, and let $\Phi_A$ denote the Transformer blocks. A similar definition holds for the notation of agent $B \in S$. The multi-step text-generation process can then be expressed as
\[
V_A \xrightarrow{W_A^{in}} E_A \xrightarrow{\Phi_A} H_A \xrightarrow{Softmax \circ W_A^{out}} \tilde{V}_A \xrightarrow{Sampling} V_A \rightarrow \cdots \rightarrow V_A \xrightarrow{T_{A \rightarrow B}} V_B
\]
where $T_{A \rightarrow B}$ is a retokenization function from $\mathcal{V}_A$ to $\mathcal{V}_B$. In particular, when $A$ and $B$ share the same vocabulary, $T_{A \rightarrow B}$ is the identity map.

The corresponding multi-step generation process over vocabulary distributions, obtained by omitting the sampling operation, is

\[
\tilde{V}_A \xrightarrow{W_A^{in}} E_A \xrightarrow{\Phi_A} H_A \xrightarrow{Softmax \circ W_A^{out}} \tilde{V}_A \rightarrow \cdots \rightarrow \tilde{V}_A \xrightarrow{T_{A \rightarrow B}} \tilde{V}_B
\]

Furthermore, define $F_{A \rightarrow A}=W_A^{in} \circ \operatorname{Softmax} \circ W_A^{out}$ and $F_{A \rightarrow B}=W_B^{in} \circ T_{A \rightarrow B} \circ \operatorname{Softmax} \circ W_A^{out}$. The hidden-state transition process can then be written as
\[
H_A \xrightarrow{\hat{F}_{A \rightarrow A}} E_A \xrightarrow{\Phi_A} H_A \rightarrow \cdots \rightarrow H_A \xrightarrow{\hat{F}_{A \rightarrow B}} E_B \xrightarrow{\Phi_B} H_B
\]
where $\hat{F}_{A \rightarrow A}$ and $\hat{F}_{A \rightarrow B}$ approximate $F_{A \rightarrow A}$ and $F_{A \rightarrow B}$, respectively.

Our objective is to identify an efficient closed-form alignment function $\hat{F}_{A \rightarrow B}$, with $A$ and $B$ potentially denoting the same agent, that approximates $F_{A \rightarrow B}$ to replace the traditional decoding--retokenization--encoding pipeline. 

\subsection{Mathematical Preliminaries}

\textbf{Linear Attention.} The attention mechanism in a typical transformer is defined as
\[
\operatorname{Attention}(Q, K, V) = \operatorname{Softmax}\left(\frac{QK^T}{\sqrt{d_k}}\right)V
\]
with complexity $O(L^2 d)$, where $L$ is the sequence length and $d$ is the hidden dimension. This quadratic complexity becomes a bottleneck for efficient inference in long-context tasks. Therefore, linear attention methods \cite{linear_attentions} have been proposed to reduce the complexity to $O(Ldm)$, where $m$ is the number of features. We treat queries, keys, values, and feature vectors as row vectors: $q_i,k_j \in \mathbb{R}^{d_k}$, $v_j \in \mathbb{R}^{d_v}$, and $\phi(q_i),\phi(k_j) \in \mathbb{R}^{m}$. Linear attention approximates the softmax kernel with a feature map $\phi(\cdot): \mathbb{R}^{d_k} \rightarrow \mathbb{R}^{m}$:
\[
\begin{aligned}
\operatorname{Attn}(q_i, K, V)
&= \frac{\sum_{j=1}^{L} \exp\!\left(q_i k_j^\top / \sqrt{d_k}\right)v_j}
        {\sum_{j=1}^{L} \exp\!\left(q_i k_j^\top / \sqrt{d_k}\right)} \\
&\approx \frac{\sum_{j=1}^{L} \phi(q_i / \sqrt{d_k})\phi(k_j)^\top v_j}
        {\sum_{j=1}^{L} \phi(q_i / \sqrt{d_k})\phi(k_j)^\top} \\
&= \frac{\phi(q_i / \sqrt{d_k})\left(\sum_{j=1}^{L}\phi(k_j)^\top v_j\right)}
        {\phi(q_i / \sqrt{d_k})\left(\sum_{j=1}^{L}\phi(k_j)^\top\right)}.
\end{aligned}
\]
Thus, the key-value summary $S=\sum_{j=1}^{L}\phi(k_j)^\top v_j \in \mathbb{R}^{m \times d_v}$ and the normalization vector $z=\sum_{j=1}^{L}\phi(k_j)^\top \in \mathbb{R}^{m \times 1}$ can be computed once and reused for every query. In the causal setting, these quantities are replaced by prefix sums $S_i=\sum_{j \leq i}\phi(k_j)^\top v_j$ and $z_i=\sum_{j \leq i}\phi(k_j)^\top$, yielding linear-time autoregressive attention.

\textbf{Orthogonal Random Features.}
Kernel methods \cite{svm} represent nonlinear similarities through an inner product in a finite-dimensional feature space. In our setting, the relevant similarity is the softmax kernel $\kappa(q,k)=\exp(qk^\top)$, where the attention scaling is absorbed into $q$ and $k$. Although this kernel admits an infinite-dimensional feature representation, Performer \cite{performer} approximates it with positive random features:
\[
\phi(x)=\frac{1}{\sqrt{m}}\exp\left(-\frac{\|x\|^2}{2}\right)
[\exp(x\omega_1^\top),\ldots,\exp(x\omega_m^\top)],
\]
where $\omega_i\sim\mathcal{N}(0,I)$. This feature map satisfies $\mathbb{E}[\phi(q)\phi(k)^\top]=\exp(qk^\top)$, and thus $\phi(q)\phi(k)^\top$ provides an efficient Monte-Carlo approximation to the softmax kernel. For more stable approximations, Performer replaces independently sampled directions with orthogonal random features: it obtains a row-orthonormal matrix $Q$ from a Gaussian matrix $G \in \mathbb{R}^{m \times d}, G_{ij}\overset{\mathrm{i.i.d.}}{\sim}\mathcal{N}(0,1)$ by RQ decomposition $G=RQ$, where $QQ^\top=I_m$, and restores the Gaussian radial distribution with $r_i\sim\chi_d$:
\[
\omega_i=r_i q_i
\]
where $q_i$ is the $i$-th row of $Q$. This orthogonalization reduces the variance of the kernel approximation, and thus improves the accuracy of linear attention.

\subsection{Kernelized Alignment}

Rather than applying linear attention to a token sequence, we reinterpret the decoding--encoding transition as attention over a fixed key--value cache: the source unembedding matrix provides the keys and the target embedding matrix provides the values. Here we assume that all agents in $S$ share the same vocabulary, i.e. $\forall A, B \in S, \mathcal{V}_A=\mathcal{V}_B$. In this case, the retokenization function $T_{A \rightarrow B}$ is the identity map. Let $w_{A,j}^{out}=(W_A^{out})_{:,j}^\top$ be the row-vector form of the $j$-th output column and $w_{B,j}^{in}=(W_B^{in})_{j,:}$ be the corresponding target embedding row. The source hidden state $h$ acts as a query, $\{w_{A,j}^{out}\}_{j=1}^{|\mathcal{V}_A|}$ as vocabulary-indexed keys, and $\{w_{B,j}^{in}\}_{j=1}^{|\mathcal{V}_A|}$ as the corresponding values. The exact transition is therefore a vocabulary-indexed softmax attention operation:
\[
F_{A \rightarrow B}(h)
= \frac{\sum_{j=1}^{|\mathcal{V}_A|} \exp\!\left(h(w_{A,j}^{out})^\top / \tau\right)w_{B,j}^{in}}
        {\sum_{j=1}^{|\mathcal{V}_A|} \exp\!\left(h(w_{A,j}^{out})^\top / \tau\right)}.
\]

Using the ORF feature map $\phi$, we obtain
\[
\hat{F}_{A \rightarrow B}(h) = \frac{\phi\!\left(h / \sqrt{\tau}\right) S_{A \rightarrow B}}
        {\phi\!\left(h / \sqrt{\tau}\right) z_A}
\]
where
\begin{align*}
S_{A \rightarrow B} = \sum_{j=1}^{|\mathcal{V}_A|}
\phi\!\left(w_{A,j}^{out} / \sqrt{\tau}\right)^\top w_{B,j}^{in} \\
z_A = \sum_{j=1}^{|\mathcal{V}_A|}
\phi\!\left(w_{A,j}^{out} / \sqrt{\tau}\right)^\top.
\end{align*}

Since $W_A^{out}$ and $W_B^{in}$ are fixed matrices that can be directly retrieved from the model, our method is completely training-free. Also, $S_{A \rightarrow B}$ and $z_A$ can be precomputed and cached offline, hence explicit decoding is not required in our method. The computational complexity of the alignment is therefore reduced from $O(|\mathcal{V}_A|(d_A+d_B))$ to $O(m(d_A+d_B))$ with $m \ll |\mathcal{V}_A|$, leading to significant speedup in inference time.

\section{Experiments}

\subsection{Setup}

We briefly describe the experimental setup below; additional implementation details, task-specific hyperparameters, and prompts are provided in Appendix~\ref{app:experimental_setup}.

\textbf{LLM systems.}
We evaluate homogeneous and heterogeneous multi-agent systems built from Qwen3-4B/8B/14B \cite{qwen3}, DeepSeek-R1-Distill-Llama-8B \cite{deepseek_r1_distill_llama_8b}, and Mistral-Nemo-Instruct-2407 \cite{mistral_nemo}. The sequential topology comprises Planner, Critic, Refiner, and Judger roles, whereas the hierarchical topology uses parallel specialist agents whose representations are aggregated by a Judger. The heterogeneous study evaluates bidirectional communication between Qwen3-8B and Qwen3-14B.

\textbf{Baselines.}
We compare the following communication interfaces. 
\emph{Single agent} directly solves the task without intermediate agents. 
\emph{TextMAS} autoregressively generates a natural-language message that is appended to the downstream agent's prompt. 
\emph{Normed alignment} preserves the direction of the source hidden state while rescaling it to the target embedding's mean norm; it is applicable only to dimension-compatible models. 
\emph{Linear alignment} maps source hidden states to target input embeddings with a vocabulary-derived linear map. \cite{latentmas}
\emph{Exact-soft alignment} computes the target embedding expected under the source model's complete next-token distribution, serving as the full-softmax reference. \cite{soft_thinking}
\emph{Prototype Alignment} is our training-free random-feature approximation to exact-soft alignment. 
Finally, \emph{entropy-stopped Prototype Alignment} applies the same map while terminating latent rollout once the probed next-token entropy remains low.

\textbf{Benchmarks.}
We evaluate mathematical reasoning on GSM8K, AIME 2024, and AIME 2025 \cite{gsm8k,aime2024,aime2025}; code generation on HumanEval+ and MBPP+ \cite{evalplus}; and scientific reasoning on ARC-Easy, ARC-Challenge, GPQA-Diamond, and MedQA \cite{arc,gpqa,medqa}.

\textbf{Evaluation settings.}
All experiments run on a single NVIDIA H200 GPU. We use matched prompts, generation budgets, and random seeds within each experimental comparison. We report three metrics: task accuracy, the number of decoding operations, and total wall-clock inference time. Accuracy denotes exact-match accuracy for mathematical tasks, multiple-choice accuracy for scientific tasks, and pass@1 for code-generation tasks. Unless otherwise stated, all results are averaged over repeated runs.

\subsection{Main Results}

\textbf{Homogeneous systems.} The following table compares all communication interfaces across models, tasks, and topologies. It reports accuracy, decoding times, and total inference time.

% TODO: Insert the homogeneous-system results table here.

\begin{table*}[t]
\centering
\caption{\textbf{Main results (accuracy, \%).} All available test-set configurations are shown.
\textbf{H}/\textbf{S} denote hierarchical/sequential topology. For Latent MAS,
\textbf{N}, \textbf{K}, \textbf{K-ES}, \textbf{L}, and \textbf{Sf} denote normed, kernel,
kernel early stopping, linear, and soft alignment, respectively. In the supplied image tables,
\emph{Single} is mapped to Baseline and \emph{LatentMAS} to linear-aligned Latent MAS;
reported image values are prioritized for the linear (L) columns. A dash (-) indicates an
unavailable result. Bold denotes the best reported accuracy within one benchmark--model pair.}
\label{tab:main_results}
\setlength{\tabcolsep}{2.4pt}
\renewcommand{\arraystretch}{1.08}
\scriptsize
\begin{adjustbox}{max width=\textwidth}
\begin{tabular}{ll *{14}{c}}
\toprule
\multirow{3}{*}{\textbf{Benchmark}} & \multirow{3}{*}{\textbf{Model}}
& \multicolumn{2}{c}{\textbf{Baseline}}
& \multicolumn{2}{c}{\textbf{Text MAS}}
& \multicolumn{10}{c}{\textbf{Latent MAS}} \\
\cmidrule(lr){3-4}\cmidrule(lr){5-6}\cmidrule(lr){7-16}
& & \multicolumn{2}{c}{N} & \multicolumn{2}{c}{N}
& \multicolumn{2}{c}{N} & \multicolumn{2}{c}{K}
& \multicolumn{2}{c}{K-ES} & \multicolumn{2}{c}{L}
& \multicolumn{2}{c}{Sf} \\
\cmidrule(lr){3-4}\cmidrule(lr){5-6}\cmidrule(lr){7-8}\cmidrule(lr){9-10}
\cmidrule(lr){11-12}\cmidrule(lr){13-14}\cmidrule(lr){15-16}
& & H & S & H & S & H & S & H & S & H & S & H & S & H & S \\
\midrule

\multirow{3}{*}{ARC-C}
& 4B  & 90.78 & 91.04 & 93.17 & \textbf{93.69} & - & - & - & - & - & - & 91.70 & 92.30 & - & - \\
\cmidrule(lr){2-16}
& 8B  & 91.81 & 91.47 & \textbf{95.65} & 94.97 & 94.20 & 94.54 & 95.05 & 94.80 & 95.05 & 95.56 & 93.90 & 94.40 & - & - \\
\cmidrule(lr){2-16}
& 14B & 93.86 & 93.60 & 95.65 & 95.73 & 95.65 & \textbf{95.99} & - & - & - & - & 95.50 & 95.60 & - & - \\

\midrule
\multirow{3}{*}{ARC-E}
& 4B  & 95.40 & 95.40 & 97.10 & 96.40 & - & - & - & - & - & - & 96.80 & \textbf{98.60} & - & - \\
\cmidrule(lr){2-16}
& 8B  & 97.52 & 97.69 & 98.44 & 98.53 & 98.48 & 98.27 & 98.23 & 98.32 & 98.23 & 98.40 & 98.30 & \textbf{98.80} & - & - \\
\cmidrule(lr){2-16}
& 14B & 97.98 & 97.85 & 98.36 & 98.57 & 98.70 & 98.48 & - & - & - & - & 98.70 & \textbf{99.40} & - & - \\

\midrule
\multirow{3}{*}{GSM8K}
& 4B  & 82.40 & 82.40 & 89.40 & \textbf{89.80} & - & - & - & - & - & - & 88.40 & 88.20 & - & - \\
\cmidrule(lr){2-16}
& 8B  & 88.63 & 89.16 & 93.48 & 93.56 & 91.58 & 91.58 & 91.28 & 91.74 & 92.34 & 92.19 & 89.50 & \textbf{93.80} & - & - \\
\cmidrule(lr){2-16}
& 14B & 89.99 & 90.37 & 92.49 & 92.27 & 91.96 & 92.12 & - & - & - & - & 91.60 & \textbf{95.20} & - & - \\

\midrule
\multirow{3}{*}{HumanEval+}
& 4B  & 75.00 & 75.00 & 76.20 & 79.70 & - & - & - & - & - & - & 79.30 & \textbf{79.90} & - & - \\
\cmidrule(lr){2-16}
& 8B  & 77.29 & 76.98 & 85.52 & \textbf{89.63} & 82.62 & 84.15 & 83.84 & 83.54 & 83.23 & 83.38 & 78.00 & 80.50 & - & 82.62 \\
\cmidrule(lr){2-16}
& 14B & 81.55 & 84.15 & 89.94 & \textbf{90.70} & 88.41 & 86.89 & - & - & - & - & 86.60 & 86.50 & - & - \\

\midrule
\multirow{3}{*}{MBPP+}
& 4B  & 63.50 & 63.50 & 69.30 & 69.80 & - & - & - & - & - & - & 70.60 & \textbf{73.50} & - & - \\
\cmidrule(lr){2-16}
& 8B  & 73.74 & 74.21 & 79.43 & \textbf{79.89} & 70.97 & - & 72.75 & 75.26 & 72.16 & 75.20 & 72.20 & 74.60 & 73.21 & 73.15 \\
\cmidrule(lr){2-16}
& 14B & 76.59 & 75.46 & 80.49 & \textbf{81.88} & 76.92 & 78.17 & - & - & - & - & 73.80 & 75.70 & - & - \\

\midrule
\multirow{3}{*}{MedQA}
& 4B  & 47.70 & 47.70 & 65.00 & 65.30 & - & - & - & - & - & - & \textbf{67.30} & 66.30 & - & - \\
\cmidrule(lr){2-16}
& 8B  & 76.67 & 75.33 & \textbf{78.67} & 75.67 & 75.33 & 76.33 & 73.33 & 73.67 & 72.67 & 72.00 & 77.00 & 75.30 & - & 66.67 \\
\cmidrule(lr){2-16}
& 14B & \textbf{81.67} & 80.33 & 78.00 & 81.33 & 81.33 & 81.00 & - & - & - & - & 78.30 & 80.70 & - & - \\

\midrule
\multirow{2}{*}{GPQA-Diamond}
& 8B  & 39.90 & 39.90 & 43.00 & 43.40 & - & - & \textbf{51.39} & 48.74 & 49.75 & 48.61 & 46.90 & 45.50 & - & - \\
\cmidrule(lr){2-16}
& 14B & 48.50 & 48.50 & 52.00 & 51.50 & - & - & - & - & - & - & \textbf{53.00} & 52.00 & - & - \\

\midrule
\multirow{2}{*}{AIME'24}
& 8B  & 50.00 & 50.00 & 53.30 & 53.30 & - & - & 60.00 & \textbf{66.25} & 65.83 & 64.17 & 53.30 & 56.70 & - & 63.75 \\
\cmidrule(lr){2-16}
& 14B & 63.30 & 63.30 & 70.00 & 63.30 & - & - & - & 72.08 & - & - & \textbf{73.30} & 66.70 & - & - \\

\midrule
\multirow{2}{*}{AIME'25}
& 8B  & 46.70 & 46.70 & 50.00 & 53.30 & - & - & \textbf{53.33} & 49.58 & 52.50 & 50.42 & 50.00 & 53.30 & - & 49.17 \\
\cmidrule(lr){2-16}
& 14B & 56.70 & 56.70 & \textbf{66.70} & 60.00 & - & - & - & - & - & - & \textbf{66.70} & 63.30 & - & - \\

\bottomrule
\end{tabular}
\end{adjustbox}
\vspace{-1mm}
\end{table*}

\textbf{Heterogeneous systems.} The following table reports cross-model communication results for Qwen3-8B and Qwen3-14B. Same-model directions act as controls; cross-scale directions evaluate whether the aligned latent representation improves a receiver that also observes the original question.

% TODO: Insert the heterogeneous-system results table here.

\subsection{Reasoning Dynamics and Scaling with Latent Steps}

\textbf{Entropy of probed logits.} We probe the next-token entropy induced by intermediate hidden states during sequential latent reasoning. This diagnostic reveals whether different feedback interfaces yield stable trajectories that return to low-entropy predictions. We compare normed, linear, exact-soft, kernel, and hard-token feedback on MBPP+ and AIME 2025. The probe is observational and does not modify generation; the resulting trajectories are shown in Figure~\ref{fig:entropy_trajectory}.

\begin{figure*}[t]
\centering
\begin{minipage}[t]{0.49\textwidth}
\centering
\includegraphics[width=\linewidth]{figs/c1_entropy_vs_cumulative_step_mbpp+.pdf}
\small (a) MBPP+
\end{minipage}\hfill
\begin{minipage}[t]{0.49\textwidth}
\centering
\includegraphics[width=\linewidth]{figs/c1_entropy_vs_cumulative_step_aime2025.pdf}
\small (b) AIME 2025
\end{minipage}
\caption{Entropy of the probed next-token distribution over cumulative latent reasoning steps. Curves compare feedback interfaces across the Planner, Critic, and Refiner stages.}
\label{fig:entropy_trajectory}
\end{figure*}

\textbf{Scaling with latent steps.} We vary the number of latent reasoning steps to characterize the performance--computation trade-off. This analysis identifies whether additional latent computation continues to improve accuracy or reaches a regime of diminishing returns (Figure~\ref{fig:latent_step_scaling}).

\begin{figure}[t]
\centering
\includegraphics[width=\linewidth]{figs/m1_accuracy_vs_steps_entropy.pdf}
\caption{Accuracy as a function of the latent-step budget.}
\label{fig:latent_step_scaling}
\end{figure}

\subsection{Robustness Analysis}

\textbf{Variance of aligned states.} We assess the sensitivity of Prototype Alignment to the ORF seed by repeatedly aligning hidden states sampled across a multi-agent trajectory. For each state, we measure the dispersion of seed-specific representations around their seed-averaged centroid. This analysis characterizes approximation stability, rather than cross-state geometry preservation (Figure~\ref{fig:orf_seed_robustness}).

\begin{figure}[t]
\centering
\includegraphics[width=\linewidth]{figs/per_state_seed_robustness.png}
\caption{Seed sensitivity of Prototype Alignment measured by the state-wise relative RMS dispersion across ORF initializations.}
\label{fig:orf_seed_robustness}
\end{figure}

\textbf{Noise injection.} We evaluate linear and Prototype Alignment under matched clean and noisy latent trajectories in a Planner--Judger system. Source hidden states are perturbed with norm-scaled Gaussian noise before alignment. We report paired changes in task accuracy and the difference in degradation between methods.

% TODO: Insert the noise-robustness results table here.

\FloatBarrier
\end{comment}
\section{Conclusion}
%In this work, we introduce Prototype Align, which is an efficient and training-free interface for latent reasoning and communication in multi-agent systems. Specifically, we reformulate the vocabulary-mediated soft alignment as attention over fixed model parameters. It enables orthogonal random features to move vocabulary dependent computation offline. Experiments across homogeneous and heterogeneous multi-agent settings demonstrate the effectiveness and efficiency of Prototype Align for latent reasoning and communication. Future work will explore more complex heterogeneous settings, including communication between models with different families.
%Our analyses suggest that kernel approximation can alter latent reasoning dynamics while retaining useful characteristics of the exact mapping. 


We introduce Prototype Align, a training-free method that performs this alignment through nonlinear mixtures of vocabulary-derived prototypes constructed using orthogonal random features. Experiments on six benchmarks show improved average accuracy in homogeneous collaboration under both stopping strategies and across sequential and hierarchical configurations. These results support prototype-based alignment as a promising interface for continuous latent information transmission in multi-agent systems.

\section*{Limitations}
Prototype Align requires the sender and receiver to share a compatible vocabulary. Our experiments are limited to Qwen3-8B and Qwen3-14B. Extending the method to different model families and architectures is left for future work.

%The random-feature approximation introduces seed-dependent variation, and its error grows with the norms of hidden states and unembedding vectors. Finally, the entropy-based stopping threshold was chosen from Prototype Align's entropy dynamics and may not be optimal for the other methods.

%{\color{blue}
Finite number of random features may introduce bias from pretrained semantic distribution, which accumulates with thinking steps and thus limits the scalability. 
%}

%\xiyan{todo, future work}
%\mingzhe{Our analyses indicates that Prototype Align can bring more effective latent reasoning dynamics while retaining useful features from exact-soft alignment. We would also suggest that some cross-vocabulary mappings (e.g. optimal transportation between vocabulary tables) possibly help to generalize our approach to even broader multi-agent systems that consists of models of different families.}

\begin{comment}
    

\section*{AI use statement}

\section*{Reproducibility statement}
We provide the complete implementation of Prototype Align in an anonymous repository at \url{https://anonymous.4open.science/r/kernel_align-258F}. The paper details the models, datasets, evaluation protocols, and experimental configurations to facilitate independent reproduction and verification of our results.
\end{comment}

%Together, these findings establish Prototype Align as a practical approach to reducing alignment overhead while maintaining strong downstream performance.


%We introduced kernelized alignment, a training-free method for hidden-state-mediated LLM systems. By viewing the decoding--encoding transition as attention over the source unembedding and target embedding matrices, our method uses orthogonal random features to approximate the full-softmax transition without explicit vocabulary decoding. This formulation applies uniformly to intra-model latent reasoning and inter-model communication, and enables vocabulary-dependent quantities to be precomputed offline. Our experiments evaluate the resulting accuracy--efficiency trade-off across reasoning tasks, system topologies, and model pairings, while further examining its dynamics and robustness. An important limitation is the shared-vocabulary assumption used in the present formulation. Extending kernelized alignment to heterogeneous vocabularies, for example through optimal transport between vocabulary spaces, is a promising direction for future work.



\FloatBarrier
\bibliography{iclr2027_conference}
%\bibliographystyle{iclr2027_conference}

\clearpage
\appendix


\section{Experiment Setup}
\label{app:experimental_setup}



\begin{comment}
\begin{table}[t]
\centering
\small
\begin{tabular}{ll}
\toprule
\textbf{Model} & \textbf{Hugging Face checkpoint} \\
\midrule
Qwen3-4B  & \texttt{Qwen/Qwen3-4B} \\
Qwen3-8B  & \texttt{Qwen/Qwen3-8B} \\
Qwen3-14B & \texttt{Qwen/Qwen3-14B} \\
DeepSeek-R1-Distill-Llama-8B
    & \texttt{deepseek-ai/DeepSeek-R1-Distill-Llama-8B} \\
Mistral-Nemo-Instruct-2407
    & \texttt{mistralai/Mistral-Nemo-Instruct-2407} \\
\bottomrule
\end{tabular}
\caption{Model checkpoints used in our experiments.}
\label{tab:model_checkpoints}
\end{table}
\end{comment}



\subsection{Task Configuration}
\label{app:config_task}
To evaluate the performance of Prototype Align, we conduct experiments with a multi-agent reasoning and communication system. Specifically, we adopt two commonly used multi-agent settings: sequential and hierarchical, following \citet{latentmas}. For the \textbf{sequential} setting, we use a chain-of-thought design that includes four LLM agents: i) a \emph{Planner}, which proposes a solution plan; ii) a \emph{Critic}, which evaluates the correctness of the plan; iii) a \emph{Refiner}, which provides a refined plan; and iv) a \emph{Judger}, which produces the final answer used for evaluation. These four agents are organized sequentially. For the \textbf{hierarchical} setting, we use domain specialists in code, mathematics, and science. Each specialist reasons about the problem independently. Their intermediate outputs are then aggregated by a summarizer to produce the final answer.

We evaluate Prototype Align in two multi-agent settings: \textit{homogeneous} systems, in which agents share the same model, and \textit{heterogeneous} systems, in which agents use models with different hidden-state dimensions. For the heterogeneous setting, we use two agents: a \emph{Planner} processes the question and transmits its hidden states, while a \emph{Judger} uses the aligned states to produce the final answer. This setup allows us to assess cross-model latent communication.

\begin{comment}
\subsection{Models and Inference Infrastructure}
We conduct experiments with Qwen3-8B and Qwen3-14B \cite{qwen3}. 
\paragraph{Homogeneous systems.}
In homogeneous experiments, all four agents share the same model checkpoint. Qwen3-8B and Qwen3-14B are evaluated on all nine benchmarks. Qwen3-4B is evaluated on ARC-Easy, ARC-Challenge, GSM8K, HumanEval+, MBPP+, and MedQA. DeepSeek-R1-Distill-Llama-8B and Mistral-Nemo-Instruct-2407 are evaluated using the same sequential and hierarchical configurations.

\paragraph{Heterogeneous systems.}
\mingzhe{TODO: Do we need 8B - 8B and 14B - 14B config?}
For heterogeneous communication, we use Qwen3-14B and Qwen3-8B because they provide tokenizer-compatible vocabularies while having different hidden-state dimensions. We evaluate all four ordered sender--receiver pairs:
\[
14\text{B}\rightarrow14\text{B},\quad
14\text{B}\rightarrow8\text{B},\quad
8\text{B}\rightarrow14\text{B},\quad
8\text{B}\rightarrow8\text{B}.
\]
The two same-model directions serve as controls, whereas
$14\text{B}\rightarrow8\text{B}$ and
$8\text{B}\rightarrow14\text{B}$ measure cross-scale communication in both directions.

In this experiment, the sender acts as a Planner and receives the complete task question. We extract the complete sequence of last-layer hidden states produced during the sender's prompt prefill. After mapping these states into the receiver's input-embedding space, we prepend them to the receiver prompt. The receiver acts as the Judger, observes the original question, and uses the aligned sender states as additional evidence. Thus, the experiment measures whether latent communication provides an improvement beyond the receiver's own access to the problem, rather than testing question recovery from a hidden message.

\end{comment}






\subsection{Baselines}
We compare our method with the following training free baselines: \par
%\paragraph{Single agent.} A single model directly solves the input problem without intermediate agents.
\noindent\textbf{Text} Agents communicate by autoregressively generating natural-language messages. The generated output of an upstream agent is inserted as text into the prompt of the downstream agent. \par
\begin{comment}
\textbf{Normed alignment.}
This baseline preserves the source hidden state's direction but replaces its norm with the mean $\ell_2$ norm of target token embeddings. Specifically, for source state $h$ and target input-embedding rows $\{w_{B,v}^{in}\}_{v\in\mathcal{V}_B}$, it computes
\[
\widetilde h
=
h\,\frac{\bar r_B}{\max(\lVert h\rVert_2,10^{-6})},
\qquad
\bar r_B
=
\frac{1}{|\mathcal{V}_B|}
\sum_{v\in\mathcal{V}_B}\lVert w_{B,v}^{in}\rVert_2.
\]
It therefore preserves directions and pairwise cosine similarities while discarding the original state norm. This baseline is applicable when the source and target have compatible hidden dimensions.
\end{comment}
\noindent\textbf{Linear Align} \citep{latentmas}, which learn a ridge-regression mapping from the sender model's output representation space to the receiver model's input-embedding space using their vocabulary representations.
%\paragraph{Exact soft alignment} \citep{soft_thinking}, which computes the exact distribution mapping $F_{A\rightarrow B}$ in Eq.~\ref{eq:vocab_attention}. It approximately evaluates the sender's full next-token distribution and takes the probability weighted average of the receiver's input embeddings.
\paragraph{Exact-Soft Align.}
This baseline computes the full-vocabulary soft mapping $F_{A\rightarrow B}$, taking the probability-weighted average of receiver input embeddings \citep{soft_thinking}.
\paragraph{Filtered-Soft Align.}
This variant filters the sender's next-token distribution using top-$k=30$, top-$p=0.95$, and top-$n=10$, then renormalizes the retained probabilities before embedding aggregation \citep{soft_thinking}.

\begin{comment}
Following Soft Thinking \cite{soft_thinking}, we first compute the sender's full next-token distribution and then take the expectation of the receiver's input embeddings:
\[
F_{\mathrm{soft}}(h_A)
=
\sum_{i=1}^{V}
p_A(i\mid h_A)\,W_{\mathrm{in},i}^{B},
\]
where
\[
p_A(i\mid h_A)
=
\operatorname{softmax}
\left(
\frac{W_{\mathrm{out}}^{A}h_A+b_A}{\tau}
\right)_i.
\]
We use $\tau=0.6$. This method is a vocabulary-distribution-mediated reference and does not sample or apply an argmax over the vocabulary.
\end{comment}

%\paragraph{Prototype Alignment.} Our method approximates the exact soft mapping using positive orthogonal random features. Vocabulary-dependent terms are preaggregated before latent rollout, reducing the online dependence on vocabulary size. Unless otherwise specified, we use $m=1024$ random features in the end-to-end experiments and $m=2048$ in the dedicated heterogeneous communication study.

%Normed alignment is excluded from cross-model 14B--8B comparisons because the two models have different hidden-state dimensions. The heterogeneous evaluation therefore compares linear, exact soft, kernel, and direct-text conditions.



%For registered reasoning models, we rely on the reasoning cue inserted by the native chat template; for other models, we explicitly add the corresponding reasoning instruction.


%Unless otherwise specified, we use Hugging Face Transformers for both text generation and latent-state rollout. We use a fixed random seed for sampling and alignment construction, as detailed below.


\subsection{Benchmarks and Evaluation}
We evaluate our method on nine benchmarks spanning three domains: i) mathematical reasoning with GSM8K, AIME 2024, and AIME 2025 \citep{gsm8k,aime2024,aime2025}; ii) code generation with HumanEval+ and MBPP+ \citep{evalplus}; and iii) scientific and medical question answering with ARC-Easy, ARC-Challenge, GPQA-Diamond, and MedQA \citep{arc,gpqa,medqa}. AIME 2024 and AIME 2025 provide their competition problems through a \texttt{train}-named split, which we use as the evaluation set. For other benchmarks, we use the test split for evaluations.



\begin{comment}
\begin{table}[t]
\centering
\small
\begin{tabular}{llll}
\toprule
\textbf{Category} & \textbf{Benchmark} & \textbf{Split} & \textbf{Metric} \\
\midrule
Mathematics
    & GSM8K & Test & Exact-match accuracy \\
    & AIME 2024 & Train/evaluation & Exact-match accuracy \\
    & AIME 2025 & Train/evaluation & Exact-match accuracy \\
\midrule
Code
    & HumanEval+ & Test & Execution-based pass@1 \\
    & MBPP+ & Test & Execution-based pass@1 \\
\midrule
Science
    & ARC-Easy & Test & Accuracy \\
    & ARC-Challenge & Test & Accuracy \\
    & GPQA-Diamond & Test & Accuracy \\
    & MedQA & Local evaluation set & Accuracy \\
\bottomrule
\end{tabular}
\caption{Benchmarks, evaluation splits, and task-specific definitions of accuracy.}
\label{tab:benchmark_details}
\end{table}
\end{comment}

We measure task performance across three categories. i) \textit{Mathematical reasoning} includes GSM8K and AIME. We extract and normalize the answer enclosed in \verb|\boxed{}| and compare it with the reference answer using exact-match accuracy. ii) \textit{Question answering} includes ARC, GPQA-Diamond, and MedQA. We extract the selected option and report answer-choice accuracy. iii) \textit{Code generation} includes HumanEval+ and MBPP+. We extract the Python code block, append the EvalPlus tests, and execute the resulting program with a 10-second timeout. A solution is counted as correct only if it passes all associated tests. Because each run produces one candidate per problem, this corresponds to pass@1. Parsing and execution failures are counted as incorrect.

\subsection{Inference Efficiency}
\label{app:inference_efficiency}

% \mingzhe{DOING}
Table~\ref{tab:baseline_reference_throughput} reports the reference
throughputs used to convert the published Text and Linear Align token counts
into the time values in the main evaluation.  For baseline $b$, multi-agent
setting $s$, model $m$, and benchmark $d$, the conversion is
$\widehat{T}_{b,s,m,d}=N_{b,s,m,d}/r_{b,s,m,d}$.  The tabulated values are
reported in tokens per second and are recovered from the displayed token and
time values at their reported precision.

\begin{table*}[t]
\centering
\setlength{\tabcolsep}{3pt}
\renewcommand{\arraystretch}{1.0}
\small
\resizebox{0.8\textwidth}{!}{%
\begin{tabular}{@{}lllrrrrrr@{}}
\toprule
Setting & Model & Baseline
& \makecell{AIME\\2024}
& \makecell{AIME\\2025}
& \makecell{Human\\Eval+}
& MBPP+
& \makecell{GPQA-\\Diamond}
& MedQA \\
\midrule
\multirow{4}{*}{Sequential}
& \multirow{2}{*}{Qwen3-8B}
& Text         & 32.79 & 31.67 & 70.43 & 65.20 & 61.20 & 107.91 \\
& & Linear & 33.24 & 37.49 & 74.56 & 66.25 & 54.69 & 87.69 \\
\cmidrule(lr){2-9}
& \multirow{2}{*}{Qwen3-14B}
& Text         & 32.31 & 30.60 & 74.45 & 65.24 & 50.66 & 82.79 \\
& & Linear & 32.32 & 30.60 & 74.53 & 65.10 & 50.64 & 82.93 \\
\midrule
\multirow{4}{*}{Hierarchical}
& \multirow{2}{*}{Qwen3-8B}
& Text         & 38.41 & 38.76 & 81.56 & 60.75 & 62.64 & 83.55 \\
& & Linear & 38.42 & 38.76 & 81.67 & 60.77 & 62.64 & 83.22 \\
\cmidrule(lr){2-9}
& \multirow{2}{*}{Qwen3-14B}
& Text         & 31.49 & 31.00 & 89.56 & 114.04 & 45.02 & 55.28 \\
& & Linear & 31.50 & 31.00 & 89.47 & 114.13 & 45.02 & 55.15 \\
\bottomrule
\end{tabular}%
}
\caption{Reference throughput $r_{b,s,m,d}$ (output tokens per second) used to convert the published Text and Linear Align token counts into the inference times reported in Tables~\ref{tab:main_results_acc} and~\ref{tab:hier_results}.}
\label{tab:baseline_reference_throughput}
\end{table*}

\subsection{Experimental Setting}
% \mingzhe{TODO: I will validate the exact settings here. Also, the choice of hyperparameter, as mentioned in the main body as well, is tuned on a development set. I would suggest that the approach to obtain dev set needs to be provided. Utilizing codex-gpt-5.6-sol-high to imitate the style of corresponding datasets, we create questions in dev set for hyperparameter tuning. The dev set can be published on github.}
%Unless otherwise specified, 
%For generation, we use nucleus-sampling threshold $p=0.95$. We set the ridge regularization coefficient to $\lambda_{\mathrm{ridge}}=10^{-5}$ for Linear Align. The temperatures for Prototype Align and Exact-Soft Align are both set to $0.6$, i.e., $\tau_{\mathrm{kernel}}=\tau_{\mathrm{soft}}=0.6$.

We use $m=1024$ orthogonal random features for Prototype Align and set the ridge regularization coefficient to $\lambda_{\mathrm{ridge}}=10^{-5}$ for Linear Align. The temperatures for Prototype Align and Exact-Soft Align are both set to $0.6$, i.e., $\tau_{\mathrm{kernel}}=\tau_{\mathrm{soft}}=0.6$. For generation, we use nucleus sampling with a top-$p$ threshold of $0.95$. 
We construct the orthogonal random features using the same seed as the corresponding generation run. Vocabulary statistics are precomputed in chunks of 4096 tokens. Exact soft alignment processes at most 32 hidden-state queries per chunk.

\begin{comment}
The default alignment parameters are:
\[
m=1024,\qquad
\lambda_{\mathrm{ridge}}=10^{-5},\qquad
\tau_{\mathrm{kernel}}=\tau_{\mathrm{soft}}=0.6
\]
where \(m\) is the dimension of orthogonal random features, \(\lambda_{\mathrm{ridge}}\) the ridge coefficient of linear alignment, \(\tau_{\mathrm{kernel}}\) and \(\tau_{\mathrm{soft}}\) the temperature of kernel and exact-soft alignment respectively.
\end{comment}



Task-dependent generation and latent-rollout settings are listed in Table~\ref{tab:task_hyperparameters}. 
%A latent-step value of zero indicates that no intermediate latent rollout is performed for that topology and task. 
%We tune the best hyperparameters for each task and topology on development sets that is extended from the original datasets. 
For each task and topology, we select the optimal number of latent reasoning steps from $\{0, 10, 20, 40, 80\}$ using development sets generated from the original datasets.
%More specifically, we randomly sample 50 problems from each dataset, and use gpt-5.6-sol \footnote{\url{https://developers.openai.com/api/docs/models/gpt-5.6-sol}} to create 50 new problems to comprise of a development set for that dataset, following the style of the sampled problems.
For each dataset, we use GPT-5.6-sol\footnote{\url{https://developers.openai.com/api/docs/models/gpt-5.6-sol}} to generate 50 new problems in the style of the sampled problems, forming a development set for that dataset.

%\subsection{Experimental Setting} 
We conduct experiments with Qwen3-8B \footnote{\url{https://huggingface.co/Qwen/Qwen3-8B}} and Qwen3-14B \footnote{\url{https://huggingface.co/Qwen/Qwen3-14B}}\citep{qwen3}. We use the model-specific chat template provided by its tokenizer. All experiments are run on a single NVIDIA H200 GPU. Models are evaluated in bfloat16 and use KV caching during generation. 

\begin{comment}

\begin{table}[t]
\centering
\scriptsize
\begin{tabular}{lrrrrr}
\toprule
\textbf{Benchmark}
& \textbf{Max tokens}
& \textbf{Batch size}
& \textbf{Sequential steps}
& \textbf{Hierarchical steps}
& \textbf{Runs} \\
\midrule
GSM8K        & 2048  & 16 & 20 & 20 & 1 \\
AIME 2024    & 20000 & 2  & 20 & 10 & 8 \\
AIME 2025    & 20000 & 2  & 0  & 10 & 8 \\
ARC-Easy     & 2048  & 16 & 10 & 10 & 1 \\
ARC-Challenge& 2048  & 16 & 20 & 10 & 1 \\
GPQA-Diamond & 8000  & 5  & 10 & 80 & 4 \\
HumanEval+   & 4096  & 8  & 20 & 40 & 4 \\
MBPP+        & 4096  & 8  & 10 & 80 & 4 \\
MedQA        & 4096  & 8  & 0  & 10 & 1 \\
\bottomrule
\end{tabular}
\caption{Task-specific generation and latent-rollout hyperparameters. ``Runs'' denotes the number of generation seeds over which results are averaged.}
\label{tab:task_hyperparameters}
\end{table}
\end{comment}

\begin{table*}[t]
\centering
\begin{adjustbox}{max width=\textwidth}
\begin{tabular}{lrrrrrrrrr}
\toprule
& GSM8K
& \makecell{AIME\\2024}
& \makecell{AIME\\2025}
& \makecell{ARC-\\Easy}
& \makecell{ARC-\\Challenge}
& \makecell{GPQA-\\Diamond}
& \makecell{Human\\Eval+}
& MBPP+
& MedQA \\
\midrule
Max tokens         & 2048 & 20000 & 20000 & 2048 & 2048 & 8000 & 4096 & 4096 & 4096 \\
Batch size         & 16   & 2     & 2     & 16   & 16   & 5    & 8    & 8    & 8 \\
%Sequential steps   & 20   & 20    & 10     & 10   & 20   & 10   & 20   & 10   & 10 \\
%Hierarchical steps & 20   & 10    & 10    & 10   & 10   & 80   & 40   & 80   & 10 \\
\bottomrule
\end{tabular}
\end{adjustbox}
\caption{Task-specific generation and latent-rollout hyperparameters.}
\label{tab:task_hyperparameters}
\end{table*}

\begin{comment}
    

\subsection{Heterogeneous Communication Protocol}

We evaluate all four ordered Qwen3 model pairs, $14\mathrm{B}\rightarrow14\mathrm{B}$, $14\mathrm{B}\rightarrow8\mathrm{B}$, $8\mathrm{B}\rightarrow14\mathrm{B}$, and $8\mathrm{B}\rightarrow8\mathrm{B}$. For each question, the sender processes the Planner prompt in one prefill pass, yielding the complete sequence of last-layer states $H_A\in\mathbb{R}^{L\times d_A}$. We align every state into the receiver embedding space and prepend the resulting sequence to the Judger's ordinary input embeddings. The Judger receives the original question in all conditions.

We evaluate 30 questions per benchmark from GPQA-Diamond, AIME 2024, AIME 2025, GSM8K, MedQA, and ARC-Challenge, sampled without replacement with seed 42. Prototype Alignment uses $m=2048$, $\tau=0.6$, and ORF seed 101. Receiver generation uses seed 77 and a maximum of 512 new tokens. We report accuracy, parsing success, and the number of communicated hidden states, with paired question-level 95\% bootstrap confidence intervals based on 1,000 replicates.
\end{comment}

\begin{table*}[t]
\centering
\setlength{\tabcolsep}{5pt}
\renewcommand{\arraystretch}{1.12}
\small
\resizebox{\textwidth}{!}{%
%\xiyan{%
\begin{tabular}{@{}lccccc@{}}
\toprule
Method
& \makecell{AIME\\2024}
& \makecell{Human\\Eval+}
& MBPP+
& \makecell{GPQA-\\Diamond}
& AVG \\
\midrule
Linear Align
& 0.06953$\pm$0.00079
& 0.08970$\pm$0.00030
& 0.09706$\pm$0.00115
& 0.08524$\pm$0.00074
& 0.08538 \\
Exact-Soft Align
& 1.54775$\pm$0.00045
& 1.54572$\pm$0.03707
& 1.52102$\pm$0.01314
& 1.54910$\pm$0.01519
& 1.54090 \\
Prototype Align
& 0.24285$\pm$0.00514
& 0.25071$\pm$0.00229
& 0.26837$\pm$0.00273
& 0.25926$\pm$0.00446
& 0.25530 \\
\bottomrule
\end{tabular}%
}%
%}
\caption{Alignment time per step (milliseconds) for Qwen3-8B under the sequential configuration.}
\label{tab:qwen3_8b_sequential_alignment_ms}
\end{table*}

\section{Results}
\label{app:results}

\subsection{Per Step Inference Efficiency}
\label{app:per_step_efficiency}
Table~\ref{tab:qwen3_8b_sequential_alignment_ms} provides the per-benchmark breakdown of the alignment times summarized in Table~\ref{tab:qwen3_8b_sequential_alignment_avg}. We evaluate Qwen3-8B under the sequential configuration on AIME 2024, HumanEval+, MBPP+, and GPQA-Diamond. For each repeat, the effective number of steps is the total number of latent output steps divided by \texttt{generate\_bs}; per-step time is then computed as $1000\times$ total alignment seconds divided by the effective number of steps. Each benchmark entry reports the mean $\pm$ population standard deviation over three repeat files. AVG pools all 12 repeat-level per-step measurements from AIME 2024, HumanEval+, MBPP+, and GPQA-Diamond for each method and reports their combined mean value.


Prototype Align consistently reduces alignment overhead across all four benchmarks, requiring 0.24 - 0.27 ms per step, compared with 1.52 - 1.55 ms for Exact-Soft Align. This corresponds to approximately $5.67$ - $6.37\times$ speedups. Linear Align remains the fastest at 0.07 - 0.10 ms per step, while both Linear Align and Prototype Align stay below 0.3 ms across all evaluated benchmarks.


\begin{table}[t]
\centering
\setlength{\tabcolsep}{3pt}
\renewcommand{\arraystretch}{1.1}
\small
\begin{tabular}{@{}llccc@{}}
\toprule
Method & Metric & GSM8K
& \makecell{ARC-\\Easy}
& \makecell{ARC-\\Challenge} \\
\midrule
\multirow{2}{*}{Text}
& Acc.  & 92.30 & 99.10 & 94.60 \\
%& Token & 2324  & 2085  & 2252  \\
& Time  & 14.3  & 11.0  & 14.8  \\
\midrule
\multirow{2}{*}{Linear Align}
& Acc.  & 93.80 & 98.80 & 94.40 \\
%& Token & 860   & 490   & 529   \\
& Time  & 5.3   & 2.6   & 3.5   \\
\midrule
\multirow{2}{*}{Exact-Soft Align}
& Acc.  & 90.52 & 97.94 & 94.45 \\
%& Token & 11369 & 10190 & 12970 \\
& Time  & 58.2  & 48.8  & 65.3  \\
\midrule
\multirow{2}{*}{Prototype Align}
& Acc.  & 91.74 & 98.32 & 94.86 \\
%& Token & 642   & 472   & 570   \\
& Time  & 4.1   & 2.5   & 3.7   \\
\bottomrule
\end{tabular}
\caption{Performance and computational efficiency of Qwen3-8B
on three supplementary benchmarks under the sequential configuration.}
\label{tab:supplement_8b_results}
\end{table}


\subsection{Supplement Results}
\label{app:supplement}
We report results on GSM8K \citep{gsm8k}, ARC-Easy, and ARC-Challenge \citep{arc} in Table~\ref{tab:supplement_8b_results}. Accuracy is generally high across methods on these benchmarks, particularly on ARC-Easy, leaving limited room for improvement. Given space constraints, we therefore focus the main comparison on more challenging benchmarks and present these results here for completeness. 

%Prototype Align achieves higher accuracy than Exact-Soft Align on all three benchmarks while substantially reducing inference time. Specifically, the end-to-end speedups range from $14.2\times$ to $19.5\times$. Compared with Text and Linear Align, its accuracy is slightly lower on GSM8K and ARC-Easy but higher on ARC-Challenge. Its inference time is also comparable to Linear Align, remaining within 0.2 seconds per problem on both ARC benchmarks. These results show that its efficiency benefits extend to these benchmarks, while accuracy differences remain modest.

Prototype Align achieves higher accuracy than Exact-Soft Align on all three supplementary benchmarks, with gains of 1.22, 0.38, and 0.41 percentage points on GSM8K, ARC-Easy, and ARC-Challenge, respectively. Compared with Text and Linear Align, its accuracy is lower on GSM8K and ARC-Easy but higher on ARC-Challenge. These results indicate that the advantage over Exact-Soft Align also holds on these benchmarks, although Prototype Align does not consistently outperform the other baselines.

\begin{comment}
\begin{table*}[t]
\centering
\caption{Performance and efficiency of Qwen3-8B and Qwen3-14B under the hierarchical configuration (mean\,$\pm$\,std over three runs). \textit{Acc} (\%, $\uparrow$) denotes task accuracy, \textit{Token} ($\downarrow$) the total number of vocabulary level decoding operations per problem, and \textit{Time} ($\downarrow$) inference time in seconds per problem. \textit{Time} of Text and Linear Align is converted to the throughput of our machine (Appendix~\ref{app:inference_efficiency}). \textbf{Bold}: best among latent alignment methods. Colored numbers in AVG: change of Prototype Align relative to Exact-Soft Align. $^\dagger$Published results from LatentMAS~\citep{latentmas}.}
\label{tab:hier_results}
\setlength{\tabcolsep}{3pt}
\renewcommand{\arraystretch}{1.12}
\small
\resizebox{\textwidth}{!}{%
\begin{tabular}{@{}l|l|*{4}{c}|*{4}{c}@{}}
\toprule
\multirow{2}{*}{Data}
& \multirow{2}{*}{Metric}
& \multicolumn{4}{c}{Qwen3-8B}
& \multicolumn{4}{c}{Qwen3-14B} \\
\cmidrule(lr){3-6}\cmidrule(lr){7-10}
& & Text
& \makecell{Linear\\Align}
& \makecell{Exact-Soft\\Align}
& \makecell{Kernel\\Align}
& Text
& \makecell{Linear\\Align}
& \makecell{Exact-Soft\\Align}
& \makecell{Kernel\\Align} \\
\midrule

\multirow{3}{*}{\makecell[l]{AIME\\2024}}
& Acc.
& 53.30
& 53.30
& 64.44\smallerr{1.93}
& \textbf{65.56\smallerr{1.57}}
& 70.00
& \textbf{73.30}
& 71.11\smallerr{5.67}
& 72.22\smallerr{1.57} \\
& Token
& 42629
& \textbf{7526}
& 22952\smallerr{748}
& 9102\smallerr{389}
& 29025
& 10230
& 23892\smallerr{660}
& \textbf{8844\smallerr{104}} \\
& Time
& 1109.7
& \textbf{195.9}
& 523.6\smallerr{66.9}
& 240.8\smallerr{5.2}
& 921.7
& 324.8
& 640.9\smallerr{33.3}
& \textbf{282.7\smallerr{2.9}} \\
\midrule

\multirow{3}{*}{\makecell[l]{AIME\\2025}}
& Acc.
& 50.00
& 50.00
& 55.56\smallerr{6.29}
& \textbf{58.89\smallerr{3.14}}
& 66.70
& \textbf{66.70}
& 62.22\smallerr{6.29}
& 64.44\smallerr{1.57} \\
& Token
& 53929
& 13230
& 29403\smallerr{987}
& \textbf{11207\smallerr{253}}
& 50003
& \textbf{9527}
& 24307\smallerr{148}
& 10817\smallerr{88} \\
& Time
& 1391.2
& 341.3
& 699.0\smallerr{32.0}
& \textbf{292.3\smallerr{10.8}}
& 1612.8
& \textbf{307.3}
& 689.2\smallerr{11.8}
& 346.2\smallerr{8.1} \\
\midrule

\multirow{3}{*}{\makecell[l]{Human\\Eval+}}
& Acc.
& 76.80
& 78.00
& 83.94\smallerr{0.35}
& \textbf{84.35\smallerr{0.76}}
& 84.10
& 86.60
& \textbf{90.85\smallerr{1.72}}
& 87.60\smallerr{1.25} \\
& Token
& 8768
& \textbf{1274}
& 11118\smallerr{488}
& 1697\smallerr{27}
& 8114
& 1512
& 21263\smallerr{1076}
& \textbf{1439\smallerr{19}} \\
& Time
& 107.5
& \textbf{15.6}
& 73.90\smallerr{1.1}
& 20.5\smallerr{0.4}
& 90.6
& \textbf{16.9}
& 185.4\smallerr{12.7}
& 23.7\smallerr{0.4} \\
\midrule

\multirow{3}{*}{MBPP+}
& Acc.
& 71.90
& 72.20
& \textbf{73.72\smallerr{0.25}}
& 73.02\smallerr{0.86}
& 73.00
& 73.80
& 78.57\smallerr{0.50}
& \textbf{78.66\smallerr{0.90}} \\
& Token
& 7703
& \textbf{1264}
& 12809\smallerr{472}
& 1495\smallerr{13}
& 7458
& 1187
& 26844\smallerr{125}
& \textbf{1133\smallerr{10}} \\
& Time
& 126.8
& 20.8
& 92.1\smallerr{5.7}
& \textbf{20.0\smallerr{0.0}}
& 65.4
& \textbf{10.4}
& 241.5\smallerr{0.1}
& 21.6\smallerr{0.2} \\
\midrule

\multirow{3}{*}{\makecell[l]{GPQA-\\Diamond}}
& Acc.
& 43.00
& 46.90
& 43.77\smallerr{2.38}
& \textbf{52.36\smallerr{2.34}}
& 52.00
& 53.00
& 54.38\smallerr{3.60}
& \textbf{56.06\smallerr{1.80}} \\
& Token
& 22450
& \textbf{3395}
& 20504\smallerr{412}
& 3728\smallerr{106}
& 20931
& 3606
& 21220\smallerr{530}
& \textbf{2749\smallerr{69}} \\
& Time
& 358.4
& \textbf{54.2}
& 316.1\smallerr{11.2}
& 59.4\smallerr{2.0}
& 464.9
& 80.1
& 296.5\smallerr{4.6}
& \textbf{61.2\smallerr{0.9}} \\
\midrule

\multirow{3}{*}{MedQA}
& Acc.
& 76.30
& \textbf{77.00}
& 66.44\smallerr{0.51}
& 72.55\smallerr{1.35}
& 78.00
& 78.30
& 51.89\smallerr{3.08}
& \textbf{79.89}\smallerr{1.64} \\
& Token
& 6893
& \textbf{1007}
& 19628\smallerr{31}
& 1351\smallerr{16}
& 5473
& 899
& 23619\smallerr{70}
& \textbf{865}\smallerr{6} \\
& Time
& 82.5
& \textbf{12.1}
& 157.5\smallerr{2.4}
& 15.3\smallerr{0.5}
& 99.0
& 16.3
& 178.3\smallerr{4.5}
& \textbf{12.3}\smallerr{0.6} \\
\midrule\midrule

\multirow{3}{*}{AVG}
& Acc.
& 61.88 & 62.90 & 64.65 & \textbf{67.79}
& 70.63 & 71.95 & 68.17 & \textbf{73.15} \\
& Token
& 23728.67 & \textbf{4616.00} & 19402.33 & 4763.33
& 20167.33 & 4493.50 & 23524.17 & \textbf{4307.83} \\
& Time
& 529.35 & \textbf{106.65} & 310.37 & 108.05
& 542.40 & 125.97 & 371.97 & \textbf{124.63} \\
\bottomrule
\end{tabular}%
}
\end{table*}
\end{comment}

\begin{table*}[t]
\centering
\setlength{\tabcolsep}{3pt}
\renewcommand{\arraystretch}{1.08}
\small
\resizebox{\textwidth}{!}{%
\begin{tabular}{@{}ll c ccc>{\columncolor{oursbg}}c cc>{\columncolor{oursbg}}c@{}}
\toprule
& & & \multicolumn{4}{c}{Fixed Steps} & \multicolumn{3}{c}{Entropy Stopping} \\
\cmidrule(lr){4-7}\cmidrule(lr){8-10}
Data & Metric & Text$^\dagger$ & \makecell{Linear\\Align$^\dagger$} & \makecell{Exact-Soft\\Align} & \makecell{Filtered-Soft\\Align} & \makecell{\textbf{Prototype}\\\textbf{Align}} & \makecell{Exact-Soft\\Align} & \makecell{Filtered-Soft\\Align} & \makecell{\textbf{Prototype}\\\textbf{Align}} \\
\midrule
\multirow{2}{*}{AIME 2024} & Acc & 53.30 & 53.30 & 54.44\tabsd{5.67} & 46.67\tabsd{2.72} & \textbf{65.56}\tabsd{1.57} & 64.44\tabsd{1.57} & 63.33\tabsd{2.72} & 62.22\tabsd{4.16} \\
 & Time & 1109.7 & \textbf{195.9} & 253.3\tabsd{34.7} & 263.3\tabsd{25.4} & 240.8\tabsd{5.2} & 523.6\tabsd{66.9} & 476.3\tabsd{12.0} & 283.1\tabsd{33.6} \\
\midrule
\multirow{2}{*}{AIME 2025} & Acc & 50.00 & 50.00 & 40.00\tabsd{11.86} & 37.78\tabsd{5.67} & \textbf{58.89}\tabsd{3.14} & 55.56\tabsd{6.29} & 51.11\tabsd{4.16} & 53.33\tabsd{5.44} \\
 & Time & 1391.2 & 341.3 & \textbf{261.4}\tabsd{10.4} & 303.2\tabsd{12.4} & 292.3\tabsd{10.8} & 699.0\tabsd{32.0} & 596.5\tabsd{79.2} & 318.7\tabsd{20.8} \\
\midrule
\multirow{2}{*}{HumanEval+} & Acc & 76.80 & 78.00 & 56.30\tabsd{1.88} & 54.67\tabsd{0.57} & \textbf{84.35}\tabsd{0.76} & 83.94\tabsd{0.29} & 84.15\tabsd{1.49} & 83.54\tabsd{2.59} \\
 & Time & 107.5 & \textbf{15.6} & 18.2\tabsd{0.8} & 18.0\tabsd{0.7} & 20.5\tabsd{0.4} & 73.90\tabsd{1.1} & 86.8\tabsd{13.2} & 19.3\tabsd{0.0} \\
\midrule
\multirow{2}{*}{MBPP+} & Acc & 71.90 & 72.20 & 65.43\tabsd{0.25} & 65.34\tabsd{1.51} & 73.02\tabsd{0.86} & 73.72\tabsd{0.25} & \textbf{73.99}\tabsd{0.82} & 70.37\tabsd{0.78} \\
 & Time & 126.8 & 20.8 & \textbf{16.5}\tabsd{0.4} & 17.1\tabsd{0.2} & 20.0\tabsd{0.0} & 92.1\tabsd{5.7} & 76.4\tabsd{4.7} & 20.1\tabsd{0.9} \\
\midrule
\multirow{2}{*}{\makecell[l]{GPQA-\\Diamond}} & Acc & 43.00 & 46.90 & 38.89\tabsd{2.30} & 36.70\tabsd{1.04} & \textbf{52.36}\tabsd{2.34} & 43.77\tabsd{2.38} & 43.27\tabsd{2.12} & 49.49\tabsd{0.41} \\
 & Time & 358.4 & \textbf{54.2} & 62.2\tabsd{1.9} & 66.0\tabsd{0.3} & 59.4\tabsd{2.0} & 316.1\tabsd{11.2} & 257.4\tabsd{7.9} & 66.1\tabsd{3.1} \\
\midrule
\multirow{2}{*}{MedQA} & Acc & 76.30 & \textbf{77.00} & 68.00\tabsd{1.09} & 67.56\tabsd{2.08} & 72.56\tabsd{1.10} & 66.44\tabsd{0.42} & 65.33\tabsd{0.00} & 74.56\tabsd{0.87} \\
 & Time & 82.5 & \textbf{12.1} & 16.7\tabsd{0.8} & 16.6\tabsd{0.5} & 15.3\tabsd{0.5} & 157.5\tabsd{2.4} & 142.5\tabsd{4.1} & 15.1\tabsd{0.3} \\
\midrule
\multirow{2}{*}{\textbf{AVG}} & Acc & 61.88 & 62.90 & 53.84 & 51.45 & \textbf{67.79} & 64.65 & 63.53 & 65.59 \\
 & Time & 529.35 & 106.65 & \textbf{104.7} & 114.0 & 108.05 & 310.37 & 272.6 & 120.40 \\
\bottomrule
\end{tabular}}
\caption{Accuracy and end-to-end inference time of Qwen3-8B under the hierarchical configuration (mean $\pm$ population standard deviation over available runs). Acc (\%, $\uparrow$) denotes task accuracy or pass@1 for code generation; Time (s, $\downarrow$) denotes seconds per problem. Text and Linear Align times are estimated using reference throughput on our hardware (Appendix~\ref{app:inference_efficiency}). Bold indicates the highest accuracy among available latent alignment configurations and the lowest time among all available configurations for each benchmark. $^\dagger$Published results from LatentMAS~\citep{latentmas}.}
\label{tab:hier_results}
\end{table*}

\begin{table*}[t]
\centering
\setlength{\tabcolsep}{3pt}
\renewcommand{\arraystretch}{1.08}
\small
\resizebox{\textwidth}{!}{%
\begin{tabular}{@{}ll c ccc>{\columncolor{oursbg}}c cc>{\columncolor{oursbg}}c@{}}
\toprule
& & & \multicolumn{4}{c}{Fixed Steps} & \multicolumn{3}{c}{Entropy Stopping} \\
\cmidrule(lr){4-7}\cmidrule(lr){8-10}
Data & Metric & Text$^\dagger$ & \makecell{Linear\\Align$^\dagger$} & \makecell{Exact-Soft\\Align} & \makecell{Filtered-Soft\\Align} & \makecell{\textbf{Prototype}\\\textbf{Align}} & \makecell{Exact-Soft\\Align} & \makecell{Filtered-Soft\\Align} & \makecell{\textbf{Prototype}\\\textbf{Align}} \\
\midrule
\multirow{2}{*}{AIME 2024} & Acc & 70.00 & 73.30 & \textbf{74.44}\tabsd{1.57} & 68.89\tabsd{4.16} & 72.22\tabsd{1.57} & 71.11\tabsd{5.67} & 68.89\tabsd{1.57} & 66.67\tabsd{4.71} \\
 & Time & 921.7 & 324.8 & 287.5\tabsd{9.8} & 292.6\tabsd{23.7} & \textbf{282.7}\tabsd{2.9} & 640.9\tabsd{33.3} & 554.2\tabsd{34.4} & 334.6\tabsd{31.7} \\
\midrule
\multirow{2}{*}{AIME 2025} & Acc & 66.70 & \textbf{66.70} & 61.11\tabsd{3.14} & 63.33\tabsd{5.44} & 64.44\tabsd{1.57} & 62.22\tabsd{6.29} & 60.00\tabsd{2.72} & 57.78\tabsd{4.16} \\
 & Time & 1612.8 & \textbf{307.3} & 366.2\tabsd{14.3} & 376.8\tabsd{12.0} & 346.2\tabsd{8.1} & 689.2\tabsd{11.8} & 657.6\tabsd{32.4} & 392.6\tabsd{20.0} \\
\midrule
\multirow{2}{*}{HumanEval+} & Acc & 84.10 & 86.60 & 86.99\tabsd{0.29} & 87.40\tabsd{0.76} & 87.60\tabsd{1.25} & \textbf{90.85}\tabsd{1.72} & 89.84\tabsd{1.60} & 86.99\tabsd{1.44} \\
 & Time & 90.6 & \textbf{16.9} & 24.0\tabsd{0.6} & 23.4\tabsd{0.6} & 23.7\tabsd{0.4} & 185.4\tabsd{12.7} & 150.1\tabsd{3.8} & 25.9\tabsd{1.5} \\
\midrule
\multirow{2}{*}{MBPP+} & Acc & 73.00 & 73.80 & 78.04\tabsd{0.22} & 76.72\tabsd{0.22} & 78.66\tabsd{0.90} & \textbf{78.75}\tabsd{0.12} & 78.13\tabsd{1.41} & 78.31\tabsd{0.75} \\
 & Time & 65.4 & \textbf{10.4} & 21.4\tabsd{0.5} & 22.3\tabsd{1.3} & 21.6\tabsd{0.2} & 241.5\tabsd{0.1} & 206.0\tabsd{8.6} & 25.6\tabsd{0.3} \\
\midrule
\multirow{2}{*}{\makecell[l]{GPQA-\\Diamond}} & Acc & 52.00 & 53.00 & 51.01\tabsd{1.89} & 50.00\tabsd{3.11} & 56.06\tabsd{1.80} & 54.38\tabsd{3.60} & 53.54\tabsd{1.89} & \textbf{56.57}\tabsd{1.43} \\
 & Time & 464.9 & 80.1 & 59.8\tabsd{4.2} & \textbf{57.8}\tabsd{0.3} & 61.2\tabsd{0.9} & 296.5\tabsd{4.6} & 279.4\tabsd{19.2} & 63.2\tabsd{3.4} \\
\midrule
\multirow{2}{*}{MedQA} & Acc & 78.00 & 78.30 & 76.78\tabsd{0.96} & 74.44\tabsd{0.87} & \textbf{79.89}\tabsd{1.64} & 51.89\tabsd{2.51} & 51.67\tabsd{2.37} & 77.00\tabsd{2.33} \\
 & Time & 99.0 & 16.3 & 11.7\tabsd{0.5} & \textbf{11.0}\tabsd{0.9} & 12.3\tabsd{0.6} & 178.3\tabsd{4.5} & 145.2\tabsd{4.0} & 14.6\tabsd{0.1} \\
\midrule
\multirow{2}{*}{\textbf{AVG}} & Acc & 70.63 & 71.95 & 71.40 & 70.13 & \textbf{73.15} & 68.20 & 67.01 & 70.55 \\
 & Time & 542.40 & 125.97 & 128.4 & 130.6 & \textbf{124.63} & 371.97 & 332.1 & 142.75 \\
\bottomrule
\end{tabular}}
\caption{Accuracy and end-to-end inference time of Qwen3-14B under the hierarchical configuration (mean $\pm$ population standard deviation over available runs). Acc (\%, $\uparrow$) denotes task accuracy or pass@1 for code generation; Time (s, $\downarrow$) denotes seconds per problem. Text and Linear Align times are estimated using reference throughput on our hardware (Appendix~\ref{app:inference_efficiency}). Bold indicates the highest accuracy among available latent alignment configurations and the lowest time among all available configurations for each benchmark. $^\dagger$Published results from LatentMAS~\citep{latentmas}.}
\label{tab:hier_results_14b}
\end{table*}

\begin{comment}
\begin{table*}[t]
\centering
\caption{Performance and efficiency of Qwen3-8B and Qwen3-14B
under the hierarchical configuration (mean $\pm$ population std over
available runs). \textit{Acc} (\%, $\uparrow$) denotes task accuracy,
\textit{Token} ($\downarrow$) the total number of vocabulary-level
decoding operations per problem, and \textit{Time} ($\downarrow$)
inference time in seconds per problem.
\textit{Time} of Text and Linear Align is converted to the
throughput of our machine
(Appendix~\ref{app:inference_efficiency}).
\textbf{Bold}: best among available latent alignment methods.
$^\dagger$Published results from LatentMAS~\citep{latentmas}.}
\label{tab:hier_results}
\setlength{\tabcolsep}{3pt}
\renewcommand{\arraystretch}{1.08}
\small
\resizebox{\textwidth}{!}{%
\begin{tabular}{@{}ll ccc>{\columncolor{oursbg}}c c ccc ccc>{\columncolor{oursbg}}c c ccc@{}}
\toprule
& & \multicolumn{8}{c}{Qwen3-8B}
  & \multicolumn{8}{c}{Qwen3-14B} \\
\cmidrule(lr){3-10}\cmidrule(l){11-18}
Data & Metric
& Text$^\dagger$
& \makecell{Linear\\Align$^\dagger$}
& \makecell{Exact-Soft\\Align}
& \raisebox{-0.3\height}{\shortstack{\textbf{Kernel}\\\textbf{Align}}}
& \makecell{Kernel-Align\\entropy}
& \makecell{Filtered-Soft\\Align}
& \makecell{Filtered-Soft\\fixed steps}
& \makecell{Exact-Soft\\fixed steps}
& Text$^\dagger$
& \makecell{Linear\\Align$^\dagger$}
& \makecell{Exact-Soft\\Align}
& \raisebox{-0.3\height}{\shortstack{\textbf{Kernel}\\\textbf{Align}}}
& \makecell{Kernel-Align\\entropy}
& \makecell{Filtered-Soft\\Align}
& \makecell{Filtered-Soft\\fixed steps}
& \makecell{Exact-Soft\\fixed steps} \\
\midrule

\multirow{3}{*}{\makecell[l]{AIME\\2024}}
& Acc
& 53.30 & 53.30 & 64.44\tabsd{1.57} & \textbf{65.56}\tabsd{1.57} & 62.22\tabsd{4.16} & 63.33\tabsd{2.72} & 46.67\tabsd{2.72} & 54.44\tabsd{5.67}
& 70.00 & 73.30 & 71.11\tabsd{5.67} & 72.22\tabsd{1.57} & 66.67\tabsd{4.71} & -- & -- & \textbf{74.44}\tabsd{1.57} \\
& Token
& 42629 & \textbf{7526} & 22952\tabsd{748} & 9102\tabsd{389} & 9968\tabsd{485} & 22{,}163\tabsd{475} & 9{,}680\tabsd{949} & 8{,}889\tabsd{1239}
& 29025 & 10230 & 23892\tabsd{660} & \textbf{8844}\tabsd{104} & 9831\tabsd{623} & -- & -- & 9{,}071\tabsd{459} \\
& Time
& 1109.7 & \textbf{195.9} & 523.6\tabsd{66.9} & 240.8\tabsd{5.2} & 283.1\tabsd{33.6} & 476.3\tabsd{12.0} & 263.3\tabsd{25.4} & 253.3\tabsd{34.7}
& 921.7 & 324.8 & 640.9\tabsd{33.3} & \textbf{282.7}\tabsd{2.9} & 334.6\tabsd{31.7} & -- & -- & 287.5\tabsd{9.8} \\
\addlinespace[2pt]\cmidrule(lr){1-18}\addlinespace[1pt]

\multirow{3}{*}{\makecell[l]{AIME\\2025}}
& Acc
& 50.00 & 50.00 & 55.56\tabsd{6.29} & \textbf{58.89}\tabsd{3.14} & 53.33\tabsd{5.44} & 51.11\tabsd{4.16} & 37.78\tabsd{5.67} & 40.00\tabsd{11.86}
& 66.70 & \textbf{66.70} & 62.22\tabsd{6.29} & 64.44\tabsd{1.57} & 57.78\tabsd{4.16} & 60.00\tabsd{2.72} & -- & 61.11\tabsd{3.14} \\
& Token
& 53929 & 13230 & 29403\tabsd{987} & 11207\tabsd{253} & 10828\tabsd{506} & 24{,}130\tabsd{1359} & 10{,}711\tabsd{438} & \textbf{9{,}422}\tabsd{355}
& 50003 & \textbf{9527} & 24307\tabsd{148} & 10817\tabsd{88} & 11167\tabsd{462} & 25{,}647\tabsd{105} & -- & 11{,}517\tabsd{368} \\
& Time
& 1391.2 & 341.3 & 699.0\tabsd{32.0} & 292.3\tabsd{10.8} & 318.7\tabsd{20.8} & 596.5\tabsd{79.2} & 303.2\tabsd{12.4} & \textbf{261.4}\tabsd{10.4}
& 1612.8 & \textbf{307.3} & 689.2\tabsd{11.8} & 346.2\tabsd{8.1} & 392.6\tabsd{20.0} & 657.6\tabsd{32.4} & -- & 366.2\tabsd{14.3} \\
\addlinespace[2pt]\cmidrule(lr){1-18}\addlinespace[1pt]

\multirow{3}{*}{\makecell[l]{Human\\Eval+}}
& Acc
& 76.80 & 78.00 & 83.94\tabsd{0.29} & \textbf{84.35}\tabsd{0.76} & 83.54\tabsd{2.59} & 84.15\tabsd{1.49} & 54.67\tabsd{0.57} & 56.30\tabsd{1.88}
& 84.10 & 86.60 & \textbf{90.85}\tabsd{1.72} & 87.60\tabsd{1.25} & 86.99\tabsd{1.44} & 89.84\tabsd{1.60} & -- & 86.99\tabsd{0.29} \\
& Token
& 8768 & 1274 & 11118\tabsd{488} & 1697\tabsd{27} & 1907\tabsd{71} & 12{,}845\tabsd{1512} & 1{,}264\tabsd{18} & \textbf{1{,}231}\tabsd{44}
& 8114 & 1512 & 21263\tabsd{1076} & \textbf{1439}\tabsd{19} & 1988\tabsd{38} & 19{,}657\tabsd{535} & -- & 1{,}547\tabsd{33} \\
& Time
& 107.5 & \textbf{15.6} & 73.90\tabsd{1.1} & 20.5\tabsd{0.4} & 19.3\tabsd{0.0} & 86.8\tabsd{13.2} & 18.0\tabsd{0.7} & 18.2\tabsd{0.8}
& 90.6 & \textbf{16.9} & 185.4\tabsd{12.7} & 23.7\tabsd{0.4} & 25.9\tabsd{1.5} & 150.1\tabsd{3.8} & -- & 24.0\tabsd{0.6} \\
\addlinespace[2pt]\cmidrule(lr){1-18}\addlinespace[1pt]

\multirow{3}{*}{MBPP+}
& Acc
& 71.90 & 72.20 & 73.72\tabsd{0.25} & 73.02\tabsd{0.86} & 70.37\tabsd{0.78} & \textbf{73.99}\tabsd{0.82} & 66.40\tabsd{0.26}\textsuperscript{($n=2$)} & 65.43\tabsd{0.25}
& 73.00 & 73.80 & 78.75\tabsd{0.12} & 78.66\tabsd{0.90} & 78.31\tabsd{0.75} & \textbf{79.89}\textsuperscript{($n=1$)} & -- & 78.04\tabsd{0.22} \\
& Token
& 7703 & 1264 & 12809\tabsd{472} & 1495\tabsd{13} & 1740\tabsd{142} & 11{,}789\tabsd{328} & 1{,}161\tabsd{8} & \textbf{1{,}153}\tabsd{28}
& 7458 & 1187 & 26844\tabsd{125} & \textbf{1133}\tabsd{10} & 1703\tabsd{71} & 26{,}061 & -- & 1{,}268\tabsd{11} \\
& Time
& 126.8 & 20.8 & 92.1\tabsd{5.7} & 20.0\tabsd{0.0} & 20.1\tabsd{0.9} & 76.4\tabsd{4.7} & 17.0\tabsd{0.0} & \textbf{16.5}\tabsd{0.4}
& 65.4 & \textbf{10.4} & 241.5\tabsd{0.1} & 21.6\tabsd{0.2} & 25.6\tabsd{0.3} & 203.1 & -- & 21.4\tabsd{0.5} \\
\addlinespace[2pt]\cmidrule(lr){1-18}\addlinespace[1pt]

\multirow{3}{*}{\makecell[l]{GPQA-\\Diamond}}
& Acc
& 43.00 & 46.90 & 43.77\tabsd{2.38} & \textbf{52.36}\tabsd{2.34} & 49.49\tabsd{0.41} & 43.27\tabsd{2.12} & 36.70\tabsd{1.04} & 38.89\tabsd{2.30}
& 52.00 & 53.00 & 54.38\tabsd{3.60} & 56.06\tabsd{1.80} & \textbf{56.57}\tabsd{1.43} & 53.54\tabsd{1.89} & -- & 51.01\tabsd{1.89} \\
& Token
& 22450 & \textbf{3395} & 20504\tabsd{412} & 3728\tabsd{106} & 3994\tabsd{98} & 20{,}367\tabsd{1014} & 4{,}025\tabsd{39} & 3{,}540\tabsd{48}
& 20931 & 3606 & 21220\tabsd{530} & 2749\tabsd{69} & 3316\tabsd{61} & 20{,}905\tabsd{678} & -- & \textbf{2{,}430}\tabsd{106} \\
& Time
& 358.4 & \textbf{54.2} & 316.1\tabsd{11.2} & 59.4\tabsd{2.0} & 66.1\tabsd{3.1} & 257.4\tabsd{7.9} & 66.0\tabsd{0.3} & 62.2\tabsd{1.9}
& 464.9 & 80.1 & 296.5\tabsd{4.6} & 61.2\tabsd{0.9} & 63.2\tabsd{3.4} & 279.4\tabsd{19.2} & -- & \textbf{59.8}\tabsd{4.2} \\
\addlinespace[2pt]\cmidrule(lr){1-18}\addlinespace[1pt]

\multirow{3}{*}{MedQA}
& Acc
& 76.30 & \textbf{77.00} & 66.44\tabsd{0.42} & 72.56\tabsd{1.10} & 74.56\tabsd{0.87} & 65.33\tabsd{0.00} & 67.56\tabsd{2.08} & 68.00\tabsd{1.09}
& 78.00 & 78.30 & 51.89\tabsd{2.51} & \textbf{79.89}\tabsd{1.64} & 77.00\tabsd{2.33} & 51.67\tabsd{2.37} & -- & 76.78\tabsd{0.96} \\
& Token
& 6893 & \textbf{1007} & 19628\tabsd{31} & 1351\tabsd{16} & 1516\tabsd{18} & 18{,}707\tabsd{814} & 1{,}210\tabsd{3} & 1{,}235\tabsd{13}
& 5473 & 899 & 23619\tabsd{70} & 865\tabsd{6} & 1470\tabsd{7} & 22{,}708\tabsd{763} & -- & \textbf{640}\tabsd{7} \\
& Time
& 82.5 & \textbf{12.1} & 157.5\tabsd{2.4} & 15.3\tabsd{0.5} & 15.1\tabsd{0.3} & 142.5\tabsd{4.1} & 16.6\tabsd{0.5} & 16.7\tabsd{0.8}
& 99.0 & 16.3 & 178.3\tabsd{4.5} & 12.3\tabsd{0.6} & 14.6\tabsd{0.1} & 145.2\tabsd{4.0} & -- & \textbf{11.7}\tabsd{0.5} \\
\midrule

\multirow{3}{*}{\textbf{AVG}}
& Acc
& 61.88 & 62.90 & 64.65 & \textbf{67.79} & 65.59 & 63.53 & 51.63 & 53.84
& 70.63 & 71.95 & 68.20 & \textbf{73.15} & 70.55 & -- & -- & 71.40 \\
& Token
& 23729 & 4616 & 19402 & 4763 & 4992 & 18{,}334 & 4{,}675 & \textbf{4{,}245}
& 20167 & 4494 & 23524 & \textbf{4308} & 4912 & -- & -- & 4{,}412 \\
& Time
& 529.35 & 106.65 & 310.37 & 108.05 & 120.40 & 272.6 & 114.0 & \textbf{104.7}
& 542.40 & 125.97 & 371.97 & \textbf{124.63} & 142.75 & -- & -- & 128.4 \\
\bottomrule
\end{tabular}}
\par\smallskip
\parbox{\textwidth}{\footnotesize\raggedright
Filtered-Soft uses top-$k=30$, top-$p=0.95$, and top-$n=10$.
The three added columns use the latest available run directory per setting.
Acc superscripts $(n=1)$ or $(n=2)$ indicate the repetition count for all three metrics;
otherwise $n=3$. Single-run results have no standard deviation.
Token is \texttt{output\_tokens}/\texttt{processed}; Time is \texttt{seconds\_per\_sample}.
Fixed-step budgets vary by task and model (recorded in the source comments).
AVG is the unweighted mean over all six tasks and is shown only when all six are available.
Missing results or incomplete task coverage are marked --.}
\end{table*}

\end{comment}

\subsection{Sequential Results}
\label{app:sequential_14b_time}
\begin{figure}[t]
\centering
\includegraphics[width=0.48\textwidth]{figs/table1_time_14b.pdf}
\caption{End-to-end inference time (seconds per problem; logarithmic scale) of Qwen3-14B under the sequential configuration.}
\label{fig:time_14b}
\end{figure}
Figure~\ref{fig:time_14b} reports the end-to-end inference time of Qwen3-14B under sequential collaboration. Under fixed budgets, Prototype Align has comparable average inference time to Exact-Soft Align and Filtered-Soft Align (114.9 vs.\ 115.0 and 116.0 seconds per problem). Under entropy-based stopping, it reduces average inference time to 131.8 seconds, compared with 289.4 and 264.5 seconds for the two soft alignment baselines, respectively.

\subsection{Hierarchical Results}
\label{app:hierarchical}
%Table~\ref{tab:hier_results} reports results under the hierarchical configuration, where specialist agents reason independently and their outputs are aggregated to produce the final answer. Prototype Align achieves average accuracies of 67.79\% and 73.15\% on Qwen3-8B and Qwen3-14B, respectively. Compared with Exact-Soft Align, Prototype Align improves average accuracy by 3.14 percentage points on Qwen3-8B while achieving a 2.87$\times$ speedup based on average end-to-end inference time. On Qwen3-14B, it achieves higher accuracy on most benchmarks and reduces inference time on all. Compared with the published Linear Align results, Prototype Align improves average accuracy by 4.89 and 1.20 points at the two scales, respectively, with similar average inference time. These results suggest that the benefits of prototype alignment extend beyond sequential collaboration to hierarchical multi-agent systems.

Tables~\ref{tab:hier_results} and~\ref{tab:hier_results_14b} report results under the hierarchical configuration, where specialist agents reason independently and their outputs are aggregated to produce the final answer. Under fixed-step budgets, Prototype Align achieves the highest average accuracy of 67.79\% on Qwen3-8B and 73.15\% on Qwen3-14B, exceeding the strongest baseline at each scale by 4.89 and 1.20 percentage points, respectively. With entropy-based stopping, it also outperforms both soft alignment baselines on average, achieving 65.59\% and 70.55\%, with gains of 0.94 and 2.35 percentage points over Exact-Soft Align. These results suggest that the average accuracy benefits of prototype alignment extend beyond sequential collaboration to hierarchical multi-agent systems under both stopping strategies.

\begin{comment}
\subsection{Baseline Reproduction}
\label{app:reproduction}
We evaluate the Linear Align baselines using the publicly released code following \cite{latentmas}. Table~\ref{tab:baseline_reproduction} reports the results obtained under our experimental setup, which differ from the published numbers. We therefore retain the published numbers as reference results in the main comparison. Standard deviations are omitted where unavailable in the source.
\end{comment}

\begin{comment}
\begin{table*}[t]
\centering
\caption{Reproduced results of Linear Align on six
benchmarks. Acc. denotes accuracy (\%), or pass@1 (\%) for code
generation; Token denotes tokens per problem; and Time denotes
inference time in seconds per problem. AVG denotes the unweighted
mean across benchmarks.}
\label{tab:baseline_reproduction}
\setlength{\tabcolsep}{4pt}
\renewcommand{\arraystretch}{1.1}
\small
\resizebox{\textwidth}{!}{%
\begin{tabular}{@{}llccccccc@{}}
\toprule
Method & Metric
& \makecell{AIME\\2024}
& \makecell{AIME\\2025}
& \makecell{HumanEval+}
& MBPP+
& \makecell{GPQA-\\Diamond}
& MedQA
& AVG \\
\midrule
\multirow{3}{*}{Linear Align}
& Acc.
& $33.33\pm0.00$ & $46.67\pm5.44$ & $17.28\pm1.88$
& $29.37\pm0.78$ & $36.70\pm3.31$ & $74.33\pm1.09$
& 39.61 \\
& Token
& $8854\pm1005$ & $9756\pm500$ & $1605\pm43$
& $1289\pm19$ & $3052\pm267$ & $1142\pm24$
& 4282.96 \\
& Time
& $266.3\pm17.0$ & $260.2\pm18.1$ & $21.5\pm0.6$
& $19.5\pm0.4$ & $55.8\pm3.1$ & $13.0\pm0.7$
& 106.06 \\
\bottomrule
\end{tabular}%
}
\end{table*}
\end{comment}

\section{Analysis}
\subsection{Alignment Analysis}
\label{app:alignment}
We extend the alignment analysis to an additional dataset HumanEval using the same protocol as in the main text. We apply Prototype Align and Exact-Soft Align to the same collected hidden states and jointly project their outputs using PCA. Figure \ref{fig:app_alignment_dynamics} (a) illustrates the results. The two sets of representations again exhibit substantial overlap in the projected space, providing further qualitative support for the representational agreement observed on MedQA.

\begin{figure*}[t]
\centering
\begin{minipage}[c]{0.36\textwidth}
    \centering
    \includegraphics[width=\linewidth]{figs/humanevalplus_pca.png}
    \small (a) Alignment outputs
\end{minipage}\hfill
\begin{minipage}[c]{0.62\textwidth}
    \centering
    \includegraphics[width=\linewidth]{figs/humanevalplus_entropy.png}
    \small (b) Next-token entropy
\end{minipage}
\caption{Alignment outputs and latent-state dynamics on HumanEval+ with Qwen3-8B. (a) Joint PCA projection of 400 Prototype Align outputs and 400 Exact-Soft Align outputs. (b) Next-token entropy across latent states for Linear Align, Prototype Align, and Exact-Soft Align.}
\label{fig:app_alignment_dynamics}
\end{figure*}

\subsection{Reasoning Dynamics Analyses}
\label{app:dynamics}
At each latent reasoning step, we project the last-layer hidden state through the model's output head and compute the entropy of the resulting next token distribution. This process is denoted as:
\[
\begin{aligned}
p_t &= \operatorname{softmax}(W_{\mathrm{out}}h_t+b), \\
H_t &= -\sum_{v=1}^{V}p_t(v)\log p_t(v).
\end{aligned}
\]
where $h_t$ denotes the last-layer hidden state at latent reasoning step $t$, and $W_{\mathrm{out}}$ and $b$ are the weight matrix and bias of the model's output head, respectively. The $p_t$ is the resulting next-token probability distribution over the vocabulary of size $V$, and $p_t(v)$ denotes the probability assigned to token $v$. This probe is used only for analysis and does not alter the latent trajectory. For experiments, we use a single Qwen3-8B agent testing on MedQA and HumanEval+ model with fixed latent step budget $K=150$. At each step, we report the mean entropy across repetitions and questions, and then estimate 95 percent confidence intervals with $1.96 \times$ standard error across questions. 

% \xiyan{TODO: K number for these two datasets. How to get the shaded band in the figure.}


%The probe is not used to modify the latent trajectory. We compare normed, linear, exact-soft, kernel, and hard-token feedback. For MBPP+, we evaluate $K\in\{20,40,60,80,100,120,140,160,180\}$ latent steps per reasoning role; for AIME 2025, we use $K\in\{20,60,100,140,180\}$. Entropy trajectories are indexed cumulatively across the Planner, Critic, and Refiner and summarized with question-level bootstrap confidence intervals.

%For latent-step scaling, we vary the latent-step budget in the same multi-agent setting and report accuracy as a function of the available latent computation. This experiment uses the same prompts, generation settings, and evaluation protocol across step budgets.

\subsection{Entropy-based Stopping}
\label{app:latent_step}
%We discover that the probed entropy of Prototype Align on development sets converges to around 0.25, which we choose as the stopping signal. We follow the original implementation in \citep{soft_thinking}: stop where entropy of consecutive 256 steps are below 0.01. The difference of entropy-based stopping signal is originated from the observable reasoning dynamics difference indicated by Figure~\ref{fig:reasoning_dynamics}.

%{\color{blue}
\paragraph{Prototype Align.}
We choose an entropy threshold of $0.25$ based on Prototype Align's entropy trajectories on development sets. During latent reasoning, we check the vocabulary entropy every 10 steps and stop after two consecutive checks below this threshold. Each intermediate agent uses at most 200 latent steps according to the typical thinking steps.

\paragraph{Soft alignment.}
For Exact-Soft Align and Filtered-Soft Align, we follow the stopping rule of \citet{soft_thinking}. Vocabulary entropy is evaluated at every latent step, and reasoning terminates after 256 consecutive steps with entropy below $0.01$, with a maximum of 10{,}000 latent steps per intermediate agent. These method-specific criteria reflect the different entropy dynamics observed in Figure~\ref{fig:reasoning_dynamics}.
%}

\subsection{Robustness Protocols}
\paragraph{ORF-seed variance.} 
\label{app:orf}
We study ORF-seed sensitivity using Qwen3-8B kernel-aligned trajectories on MedQA generated with trajectory seed 42. We draw a balanced sample of $N=500$ hidden states using sampling seed 20260821. For every hidden state $h_i$, we construct 50 independent ORF maps, each with $m=1024$ random features and kernel temperature $\tau=0.6$, while holding the embedding and unembedding weights fixed. Let $y_{i,s}=f_s(h_i)$ denote the alignment output for state $i$ under ORF seed $s$, and let
\[
\bar y_i=\frac{1}{50}\sum_{s=1}^{50}y_{i,s}.
\]
We quantify variation across seeds by the relative RMS dispersion
\[
D_i=
\frac{\sqrt{\frac{1}{50}\sum_{s=1}^{50}\lVert y_{i,s}-\bar y_i\rVert_2^2}}
{\lVert\bar y_i\rVert_2}.
\]
Smaller $D_i$ indicates greater invariance to the ORF seed. For visualization, we transform dispersion into the seed-similarity score
\[
S_i=1-D_i,
\]
so that larger values indicate greater stability. This score is defined from relative RMS dispersion.

We plot the empirical distribution of the 500 values $\{S_i\}_{i=1}^{500}$ as a weighted histogram. Both axes are displayed as percentages. Dashed reference lines mark the median and 5th percentile of $S_i$. 

Across the 500 states, $S_i$ ranges from 84.160\% to 85.775\%, with 5th percentile 84.56\% and median 85.07\%. The two reference lines in the histogram therefore appear at 84.56\% and 85.07\%, respectively.

%\subsection{E}

\begin{comment}
    

\paragraph{Noise injection.} We evaluate linear and Prototype Alignment in a two-role Qwen3-8B Planner--Judger system on 30 fixed questions from each of GSM8K, MBPP+, and AIME 2024. The Planner uses 20, 80, and 10 latent steps for these tasks, respectively. For every causally used Planner state $h_t\in\mathbb{R}^{d}$, we inject norm-scaled Gaussian noise before alignment,
\[
\widetilde h_t=h_t+\delta_t,
\qquad
\delta_t=\alpha\frac{\lVert h_t\rVert_2}{\sqrt{d}}\epsilon_t,
\qquad
\epsilon_t\sim\mathcal{N}(0,I_d),
\]
with $\alpha=0.05$. A keyed random-number generator ensures that both alignment methods receive the same perturbation for each dataset, example, role, step, and repetition without altering generation randomness. We use eight paired repetitions with generation seeds 42--49 and kernel seeds 101--108. The primary statistic is the paired accuracy change $\Delta_{\mathrm{noise}}^{(a)}=\operatorname{Acc}_{\mathrm{noisy}}^{(a)}-\operatorname{Acc}_{\mathrm{clean}}^{(a)}$, together with $\Delta_{\mathrm{robust}}=\Delta_{\mathrm{noise}}^{(\mathrm{kernel})}-\Delta_{\mathrm{noise}}^{(\mathrm{linear})}$. We compute paired bootstrap confidence intervals over examples and repetitions.
\end{comment}


\begin{comment}
\subsection{Latent Step Ablation.}
\label{app:latent_step}
We additionally evaluate an entropy-based early-stopping variant. The vocabulary entropy is checked every 10 latent steps. Rollout terminates after two consecutive low-entropy checks, corresponding to a default low-entropy span of 20 steps, using an entropy threshold of $0.25$. The maximum rollout length for this variant is 200 latent steps.
\end{comment}

% {\color{blue}
% \section{Proof of Prototype-Based Approximation}
% \label{app:prototype_approximation}

% We prove the approximation theorem in Section~\ref{sec:prototype_alignment} by first rewriting the prototype mixture as normalized kernel attention, then controlling the kernel estimation error. All statements concern a fixed query and a finite vocabulary.

% \paragraph{Equivalence to positive random-feature attention.}
% Define the kernel estimate
% \begin{equation}
% \begin{aligned}
%     &\widehat{\kappa}_m(q,k_j) \\
%     &=\frac{1}{m}\sum_{i=1}^{m}
%       \exp\!\left(\omega_i^\top q-\frac{\|q\|^2}{2}\right) \\
%     &\quad\times
%       \exp\!\left(\omega_i^\top k_j-\frac{\|k_j\|^2}{2}\right) \\
%     &=\frac{\exp(-\|q\|^2/2)}{m}
%       \sum_{i=1}^{m}\exp(\omega_i^\top q)b_{ij}.
% \end{aligned}
% \label{eq:prototype_kernel_estimator}
% \end{equation}
% Using $Z_i u_i=\sum_j b_{ij}v_j$ and $Z_i=\sum_j b_{ij}$, we obtain
% \begin{equation}
% \begin{aligned}
%     &\frac{\sum_i a_i(q)u_i}{\sum_i a_i(q)} \\
%     &=\frac{\sum_i\exp(\omega_i^\top q)\sum_j b_{ij}v_j}
%             {\sum_i\exp(\omega_i^\top q)\sum_j b_{ij}} \\
%     &=\frac{\sum_j\widehat{\kappa}_m(q,k_j)v_j}
%             {\sum_j\widehat{\kappa}_m(q,k_j)}.
% \end{aligned}
% \label{eq:prototype_attention_equivalence}
% \end{equation}
% The common factor $\exp(-\|q\|^2/2)/m$ cancels between the numerator and denominator. This establishes the exact equivalence between prototype mixing and normalized positive random-feature attention, and explains why the activation includes the vocabulary mass $Z_i$.

% \paragraph{Kernel estimation error.}
% For each marginal $\omega_i\sim\mathcal{N}(0,I_{d_A})$, the Gaussian moment-generating function gives
% \begin{equation}
% \begin{aligned}
%     &\mathbb{E}\!\left[
%       \exp\!\left(\omega_i^\top(q+k_j)
%              -\frac{\|q\|^2+\|k_j\|^2}{2}\right)
%       \right] \\
%     &=\exp(q^\top k_j) \\
%     &=:\kappa_j.
% \end{aligned}
% \label{eq:prototype_kernel_expectation}
% \end{equation}
% Thus $\widehat{\kappa}_m(q,k_j)$ is unbiased even though the directions within an ORF block are dependent. The blockwise variance bound in Appendix~\ref{app:orf_accuracy} implies
% \begin{equation}
% \begin{aligned}
%     &\operatorname{Var}[\widehat{\kappa}_m(q,k_j)] \\
%     &\leq\frac{\kappa_j^2}{m}
%           \left(\exp(\|q+k_j\|^2)-1\right).
% \end{aligned}
% \label{eq:prototype_kernel_variance}
% \end{equation}
% For any $0<\epsilon<1$, Chebyshev's inequality and a union bound over the finite vocabulary therefore yield
% \begin{equation}
% \begin{aligned}
%     &\mathbb{P}\!\left[
%        \max_{j\in\mathcal{V}}
%        \frac{|\widehat{\kappa}_m(q,k_j)-\kappa_j|}{\kappa_j}
%        >\epsilon\right] \\
%     &\leq\frac{C(q,K)}{m\epsilon^2}, \\
%     &C(q,K) \\
%     &:=\sum_{j\in\mathcal{V}}
%                  \left(\exp(\|q+k_j\|^2)-1\right).
% \end{aligned}
% \label{eq:prototype_relative_error}
% \end{equation}
% This constant is finite for fixed $q$ and keys $K=\{k_j\}_{j\in\mathcal{V}}$.

% \paragraph{Error after normalization.}
% Let $D=\sum_j\kappa_j>0$, $N=\sum_j\kappa_jv_j$, and $R=\max_j\|v_j\|_2$. Denote their estimated counterparts by $\widehat D=\sum_j\widehat{\kappa}_m(q,k_j)$ and $\widehat N=\sum_j\widehat{\kappa}_m(q,k_j)v_j$. On the event that every relative kernel error is at most $\epsilon$, we have
% \[
% \begin{aligned}
%     |\widehat D-D|&\leq\epsilon D, \\
%     \|\widehat N-N\|_2&\leq\epsilon D R, \\
%     \widehat D&\geq(1-\epsilon)D.
% \end{aligned}
% \]
% Since $\|N\|_2\leq DR$, it follows that
% \begin{equation}
% \begin{aligned}
%     &\left\|\frac{\widehat N}{\widehat D}
%              -\frac{N}{D}\right\|_2 \\
%     &\leq\frac{\|\widehat N-N\|_2}{\widehat D} \\
%     &\quad+\frac{\|N\|_2|D-\widehat D|}{D\widehat D} \\
%     &\leq\frac{2R\epsilon}{1-\epsilon}.
% \end{aligned}
% \label{eq:prototype_normalized_error}
% \end{equation}
% Combining Eqs.~\ref{eq:prototype_relative_error} and \ref{eq:prototype_normalized_error} proves convergence in probability. Taking $\epsilon=L/\sqrt{m}$ for sufficiently large $m$ bounds the probability of an output error greater than $4RL/\sqrt{m}$ by $C(q,K)/L^2$, establishing the stated $O_{\mathbb{P}}(m^{-1/2})$ rate. Finally, $N/D$ equals the readout--embedding mapping because $q^\top k_j=h_t^A(w_{A,j}^{\mathrm{out}})^\top/\tau$. This mapping is precisely softmax attention over the sender's unembedding keys and receiver's embedding values. \hfill$\square$

% The theorem does not assert that the normalized finite-$m$ estimator is unbiased, that the bound is uniform over changing queries, or that closer approximation guarantees higher downstream accuracy. It establishes a training-free approximation of the target interface; the behavior of finite prototype mixtures in latent reasoning is evaluated empirically.
% }


% {\color{blue}
\section{Proof of Prototype-Based Approximation}
\label{app:prototype_approximation}
\label{app:orf_accuracy}

We prove the approximation theorem in Section~\ref{sec:prototype_alignment} by first rewriting the prototype mixture as normalized kernel attention, then controlling the kernel estimation error. All statements concern a fixed query and a finite vocabulary.

\paragraph{ORF block construction.}
Following \citet{performer}, partition the $m$ probes into $B$ independent blocks with sizes $s_b\leq d_A$ and $\sum_{b=1}^{B}s_b=m$. In each block, draw a Haar-distributed orthonormal frame $\{\theta_{b,i}\}_{i=1}^{s_b}$ and independent radii $r_{b,i}\sim\chi_{d_A}$, independent of the frame, and set $\omega_{b,i}=r_{b,i}\theta_{b,i}$. The probes are orthogonal within each block, and each has marginal distribution $\mathcal{N}(0,I_{d_A})$. Different blocks are sampled independently. We relabel the concatenated probes as $\{\omega_i\}_{i=1}^{m}$ below. This construction covers both a single block when $m\leq d_A$ and multiple blocks when $m>d_A$.

\paragraph{Equivalence to positive random-feature attention.}
Define the kernel estimate
\begin{equation}
\begin{aligned}
    &\widehat{\kappa}_m(q,k_j) \\
    &=\frac{1}{m}\sum_{i=1}^{m}
      \exp\!\left(\omega_i^\top q-\frac{\|q\|^2}{2}\right) \\
    &\quad\times
      \exp\!\left(\omega_i^\top k_j-\frac{\|k_j\|^2}{2}\right) \\
    &=\frac{\exp(-\|q\|^2/2)}{m}
      \sum_{i=1}^{m}\exp(\omega_i^\top q)b_{ij}.
\end{aligned}
\label{eq:prototype_kernel_estimator}
\end{equation}
Using $Z_i u_i=\sum_j b_{ij}v_j$ and $Z_i=\sum_j b_{ij}$, we obtain
\begin{equation}
\begin{aligned}
    &\frac{\sum_i a_i(q)u_i}{\sum_i a_i(q)} \\
    &=\frac{\sum_i\exp(\omega_i^\top q)\sum_j b_{ij}v_j}
            {\sum_i\exp(\omega_i^\top q)\sum_j b_{ij}} \\
    &=\frac{\sum_j\widehat{\kappa}_m(q,k_j)v_j}
            {\sum_j\widehat{\kappa}_m(q,k_j)}.
\end{aligned}
\label{eq:prototype_attention_equivalence}
\end{equation}
The common factor $\exp(-\|q\|^2/2)/m$ cancels between the numerator and denominator. This establishes the exact equivalence between prototype mixing and normalized positive random-feature attention, and explains why the activation includes the vocabulary mass $Z_i$.

\paragraph{Kernel estimation error.}
For each marginal $\omega_i\sim\mathcal{N}(0,I_{d_A})$, the Gaussian moment-generating function gives
\begin{equation}
\begin{aligned}
    &\mathbb{E}\!\left[
      \exp\!\left(\omega_i^\top(q+k_j)
             -\frac{\|q\|^2+\|k_j\|^2}{2}\right)
      \right] \\
    &=\exp(q^\top k_j) \\
    &=:\kappa_j.
\end{aligned}
\label{eq:prototype_kernel_expectation}
\end{equation}
Thus $\widehat{\kappa}_m(q,k_j)$ is unbiased even though the directions within an ORF block are dependent. A second Gaussian moment calculation gives the single-feature variance $V_j=\kappa_j^2(\exp(\|q+k_j\|^2)-1)$. The orthogonal Monte-Carlo bound of \citet{performer} implies that the mean kernel estimate $\widehat{\kappa}^{(b)}(q,k_j)$ from a block of $s_b$ probes has variance at most $V_j/s_b$. Since $\widehat{\kappa}_m=\sum_b(s_b/m)\widehat{\kappa}^{(b)}$ and the blocks are independent,
\begin{equation}
\begin{aligned}
    &\operatorname{Var}[\widehat{\kappa}_m(q,k_j)] \\
    &\leq\sum_{b=1}^{B}\left(\frac{s_b}{m}\right)^2\frac{V_j}{s_b} \\
    &=\frac{\kappa_j^2}{m}
          \left(\exp(\|q+k_j\|^2)-1\right).
\end{aligned}
\label{eq:prototype_kernel_variance}
\end{equation}
For any $0<\epsilon<1$, Chebyshev's inequality and a union bound over the finite vocabulary therefore yield
\begin{equation}
\begin{aligned}
    &\mathbb{P}\!\left[
       \max_{j\in\mathcal{V}}
       \frac{|\widehat{\kappa}_m(q,k_j)-\kappa_j|}{\kappa_j}
       >\epsilon\right] \\
    &\leq\frac{C(q,K)}{m\epsilon^2}, \\
    &C(q,K) \\
    &:=\sum_{j\in\mathcal{V}}
                 \left(\exp(\|q+k_j\|^2)-1\right).
\end{aligned}
\label{eq:prototype_relative_error}
\end{equation}
This constant is finite for fixed $q$ and keys $K=\{k_j\}_{j\in\mathcal{V}}$.

\paragraph{Error after normalization.}
Let $D=\sum_j\kappa_j>0$, $N=\sum_j\kappa_jv_j$, and $R=\max_j\|v_j\|_2$. Denote their estimated counterparts by $\widehat D=\sum_j\widehat{\kappa}_m(q,k_j)$ and $\widehat N=\sum_j\widehat{\kappa}_m(q,k_j)v_j$. On the event that every relative kernel error is at most $\epsilon$, we have
\[
\begin{aligned}
    |\widehat D-D|&\leq\epsilon D, \\
    \|\widehat N-N\|_2&\leq\epsilon D R, \\
    \widehat D&\geq(1-\epsilon)D.
\end{aligned}
\]
Since $\|N\|_2\leq DR$, it follows that
\begin{equation}
\begin{aligned}
    &\left\|\frac{\widehat N}{\widehat D}
             -\frac{N}{D}\right\|_2 \\
    &\leq\frac{\|\widehat N-N\|_2}{\widehat D} \\
    &\quad+\frac{\|N\|_2|D-\widehat D|}{D\widehat D} \\
    &\leq\frac{2R\epsilon}{1-\epsilon}.
\end{aligned}
\label{eq:prototype_normalized_error}
\end{equation}
Combining Eqs.~\ref{eq:prototype_relative_error} and \ref{eq:prototype_normalized_error} proves convergence in probability. Taking $\epsilon=L/\sqrt{m}$ for sufficiently large $m$ bounds the probability of an output error greater than $4RL/\sqrt{m}$ by $C(q,K)/L^2$, establishing an $O_{\mathbb{P}}(m^{-1/2})$ rate. Finally, $N/D$ equals the readout--embedding mapping because $q^\top k_j=h_t^A(w_{A,j}^{\mathrm{out}})^\top/\tau$. This mapping is precisely softmax attention over the sender's unembedding keys and receiver's embedding values. \hfill$\square$

The theorem does not assert that the normalized finite-$m$ estimator is unbiased, that the bound is uniform over changing queries, or that closer approximation guarantees higher downstream accuracy. It establishes a training-free approximation of the target interface; the behavior of finite prototype mixtures in latent reasoning is evaluated empirically.
% }

% Required packages:
% \usepackage{booktabs,multirow,makecell,graphicx}

\providecommand{\smallerr}[1]{}
\renewcommand{\smallerr}[1]{{\scriptsize$\pm#1$}}




\begin{comment}
    

\section{Norm Discrepancy as Evidence of Distribution Shift}
\label{app:norm_shift}

An input embedding and a final-layer hidden state occupy vector spaces of the same dimension in a homogeneous model, but they are generated by different mechanisms and need not follow the same distribution. In particular, token embeddings are learned lookup vectors presented at the input of the Transformer, whereas a hidden state is the contextual output of all Transformer blocks. Passing a hidden state directly as a continuous input therefore induces an out-of-distribution input for the receiving model whenever its statistics differ from those of ordinary embeddings.

The norm provides a simple diagnostic of this mismatch. For target embedding rows $\{w_{B,v}^{in}\}_{v\in\mathcal{V}_B}$, define their mean norm as
\[
\bar r_B
=
\frac{1}{|\mathcal{V}_B|}
\sum_{v\in\mathcal{V}_B}\lVert w_{B,v}^{in}\rVert_2.
\]
For an incoming hidden state $h$, its relative norm discrepancy is
\[
D_{\mathrm{norm}}(h)
=
\frac{\left|\lVert h\rVert_2-\bar r_B\right|}{\bar r_B}.
\]
A systematic difference between the distribution of $\lVert h\rVert_2$ and that of $\lVert w_{B,v}^{in}\rVert_2$ indicates that the receiving Transformer encounters radial statistics not observed during ordinary token-input processing. Figure~\ref{fig:norm_distribution_shift} visualizes this discrepancy. This is direct evidence of distribution shift along the radial direction. It does not, by itself, characterize angular or higher-order differences between the two distributions.

Normed alignment isolates this diagnostic by retaining the direction of $h$ while setting its norm to $\bar r_B$. Consequently, comparing normed alignment with other interfaces tests whether correcting this radial mismatch alone improves downstream behavior; any remaining gap reflects distributional differences beyond vector norm.

\begin{figure*}[t]
\centering
\begin{minipage}[t]{0.49\textwidth}
\centering
\includegraphics[width=\linewidth]{figs/s0_embedding_norm_hist.pdf}
\small (a) Input-embedding norms
\end{minipage}\hfill
\begin{minipage}[t]{0.49\textwidth}
\centering
\includegraphics[width=\linewidth]{figs/s0_hidden_norm_hist.pdf}
\small (b) Hidden-state norms
\end{minipage}
\caption{Norm distributions of target input embeddings and communicated hidden states. Their discrepancy provides a radial diagnostic of the hidden-to-embedding distribution shift.}
\label{fig:norm_distribution_shift}
\end{figure*}

\end{comment}







\end{document}
